\documentclass[10pt]{article}
\usepackage[utf8]{inputenc}
\usepackage[T1]{fontenc}
\usepackage{graphicx}
\usepackage[export]{adjustbox}
\graphicspath{ {./images/} }
\usepackage{hyperref}
\hypersetup{colorlinks=true, linkcolor=blue, filecolor=magenta, urlcolor=cyan,}
\urlstyle{same}
\usepackage{amsmath}
\usepackage{amsfonts}
\usepackage{amssymb}
\usepackage[version=4]{mhchem}
\usepackage{stmaryrd}
\usepackage{bbold}
\usepackage{mathrsfs}
\usepackage{caption}
\usepackage{multirow}

%New command to display footnote whose markers will always be hidden
\let\svthefootnote\thefootnote
\newcommand\blfootnotetext[1]{%
  \let\thefootnote\relax\footnote{#1}%
  \addtocounter{footnote}{-1}%
  \let\thefootnote\svthefootnote%
}
%Overriding the \footnotetext command to hide the marker if its value is `0`
\let\svfootnotetext\footnotetext
\renewcommand\footnotetext[2][?]{%
  \if\relax#1\relax%
    \ifnum\value{footnote}=0\blfootnotetext{#2}\else\svfootnotetext{#2}\fi%
  \else%
    \if?#1\ifnum\value{footnote}=0\blfootnotetext{#2}\else\svfootnotetext{#2}\fi%
    \else\svfootnotetext[#1]{#2}\fi%
  \fi
}
\DeclareUnicodeCharacter{22C5}{\ifmmode\cdot\else{$\cdot$}\fi}
\DeclareUnicodeCharacter{0394}{\ifmmode\Delta\else{$\Delta$}\fi}
\DeclareUnicodeCharacter{22EE}{\ifmmode\vdots\else{$\vdots$}\fi}
\title{Chapter 13: Maximum Likelihood Methods}

\author{}
\date{}

\begin{document}
\maketitle

\subsection*{13.1 Introduction}
This chapter contains a general treatment of maximum likelihood estimation (MLE) under random sampling. All the models we considered in Part I could be estimated without making full distributional assumptions about the endogenous variables conditional on the exogenous variables: maximum likelihood methods were not needed. Instead, we focused primarily on zero-covariance and zero-conditional-mean assumptions, and secondarily on assumptions about conditional variances and covariances. These assumptions were sufficient for obtaining consistent, asymptotically normal estimators, some of which were shown to be efficient within certain classes of estimators.

Some texts on advanced econometrics take MLE as the unifying theme, and then most models are estimated by maximum likelihood. In addition to providing a unified approach to estimation, MLE has some desirable efficiency properties: it is generally the most efficient estimation procedure in the class of estimators that use information on the distribution of the endogenous variables given the exogenous variables. (We formalize the efficiency of MLE in Section 14.4.) So why not always use MLE?

As we saw in Part I, efficiency usually comes at the price of nonrobustness, and this is certainly the case for maximum likelihood. Maximum likelihood estimators are generally inconsistent if some part of the specified distribution is misspecified. As an example, consider from Section 9.5 a simultaneous equations model that is linear in its parameters but nonlinear in some endogenous variables. There, we discussed estimation by instrumental variables methods. We could estimate SEMs nonlinear in endogenous variables by maximum likelihood if we assumed independence between the structural errors and the exogenous variables and if we assumed a particular distribution for the structural errors, say, multivariate normal. The MLE would be asymptotically more efficient than the best GMM estimator, but failure of normality generally results in inconsistent estimators of all parameters.

As a second example, suppose we wish to estimate $\mathrm{E}(y \mid \mathbf{x})$, where $y$ is bounded between zero and one. The logistic function, $\exp (\mathbf{x} \boldsymbol{\beta}) /[1+\exp (\mathbf{x} \boldsymbol{\beta})]$, is a reasonable model for $\mathrm{E}(y \mid \mathbf{x})$, and, as we discussed in Chapter 12, nonlinear least squares provides consistent, $\sqrt{N}$-asymptotically normal estimators under weak regularity conditions. We can easily make inference robust to arbitrary heteroskedasticity in $\operatorname{Var}(y \mid \mathbf{x})$. An alternative approach is to model the density of $y$ given $\mathbf{x}$-which, of course, implies a particular model for $\mathrm{E}(y \mid \mathbf{x})$-and use MLE. As we will see, the strength of MLE is that, under correct specification of the density, we would have\\
the asymptotically efficient estimators, and we would be able to estimate any feature of the conditional distribution, such as $\mathrm{P}(y=1 \mid \mathbf{x})$. The drawback is that, except in special cases, if we have misspecified the density in any way, we will not be able to consistently estimate the conditional mean.

In most applications, specifying the distribution of the endogenous variables conditional on exogenous variables must have a component of arbitrariness, as economic theory rarely provides guidance. Our perspective is that, for robustness reasons, it is desirable to make as few assumptions as possible-at least until relaxing them becomes practically difficult. There are cases in which MLE turns out to be robust to failure of certain assumptions, but these must be examined on a case-by-case basis, a process that detracts from the unifying theme provided by the MLE approach. (One such example is nonlinear regression under a homoskedastic normal assumption; the MLE of the parameters $\boldsymbol{\beta}_{\mathrm{o}}$ is identical to the NLS estimator, and we know the latter is consistent and asymptotically normal quite generally. We will cover some other leading cases in Section 13.11 and Chapter 18.)

Maximum likelihood plays an important role in modern econometric analysis, for good reason. There are many problems for which it is indispensable. For example, in Chapters 15, 16, and 17 we study various limited dependent variable models, and MLE plays a central role.

\subsection*{13.2 Preliminaries and Examples}
Traditional maximum likelihood theory for independent, identically distributed observations $\left\{\mathbf{y}_{i} \in \mathbb{R}^{G}: i=1,2, \ldots\right\}$ starts by specifying a family of densities for $\mathbf{y}_{i}$. This is the framework used in introductory statistics courses, where $\mathbf{y}_{i}$ is a scalar with a normal or Poisson distribution. But in almost all economic applications, we are interested in estimating parameters of conditional distributions. Therefore, we assume that each random draw is partitioned as $\left(\mathbf{x}_{i}, \mathbf{y}_{i}\right)$, where $\mathbf{x}_{i} \in \mathbb{R}^{K}$ and $\mathbf{y}_{i} \in \mathbb{R}^{G}$, and we are interested in estimating a model for the conditional distribution of $\mathbf{y}_{i}$ given $\mathbf{x}_{i}$. We are not interested in the distribution of $\mathbf{x}_{i}$, so we will not specify a model for it. Consequently, the method of this chapter is properly called conditional maximum likelihood estimation (CMLE). By taking $\mathbf{x}_{i}$ to be null we cover unconditional MLE as a special case.

An alternative to viewing $\left(\mathbf{x}_{i}, \mathbf{y}_{i}\right)$ as a random draw from the population is to treat the conditioning variables $\mathbf{x}_{i}$ as nonrandom vectors that are set ahead of time and that appear in the unconditional distribution of $\mathbf{y}_{i}$. (This setup is analogous to the fixed regressor assumption in classical regression analysis.) Then the $\mathbf{y}_{i}$ cannot be identi-\\
cally distributed, and this fact complicates the asymptotic analysis. More important, treating the $\mathbf{x}_{i}$ as nonrandom is much too restrictive for all uses of maximum likelihood. In fact, later on we will cover methods where $\mathbf{x}_{i}$ contains what are endogenous variables in a structural model, but where it is convenient to obtain the distribution of one set of endogenous variables conditional on another set. Once we know how to analyze the general CMLE case, applications follow fairly directly.

It is important to understand that the subsequent results apply any time we have random sampling in the cross section dimension. Thus, the general theory applies to system estimation, as in Chapters 7 and 9, provided we are willing to assume a distribution for $\mathbf{y}_{i}$ given $\mathbf{x}_{i}$. In addition, panel data settings with large cross sections and relatively small time periods are encompassed, since the appropriate asymptotic analysis is with the time dimension fixed and the cross section dimension tending to infinity.

In order to perform maximum likelihood analysis we need to specify, or derive from an underlying (structural) model, the density of $\mathbf{y}_{i}$ given $\mathbf{x}_{i}$. We assume this density is known up to a finite number of unknown parameters, with the result that we have a parametric model of a conditional density. The vector $\mathbf{y}_{i}$ can be continuous or discrete, or it can have both discrete and continuous characteristics. In many of our applications, $\mathbf{y}_{i}$ is a scalar, but this feature does not simplify the general treatment.

We will carry along two examples to illustrate the general theory of conditional maximum likelihood. The first example is a binary response model, specifically the probit model. We postpone the uses and interepretation of binary response models until Chapter 15.

Example 13.1 (Probit): Suppose that the latent variable $y_{i}^{*}$ follows


\begin{equation*}
y_{i}^{*}=\mathbf{x}_{i} \boldsymbol{\theta}+e_{i} \tag{13.1}
\end{equation*}


where $e_{i}$ is independent of $\mathbf{x}_{i}$ (which is a $1 \times K$ vector with first element equal to unity for all $i), \boldsymbol{\theta}$ is a $K \times 1$ vector of parameters, and $e_{i} \sim \operatorname{Normal}(0,1)$. Instead of observing $y_{i}^{*}$, we observe only a binary variable indicating the sign of $y_{i}^{*}$ :

\[
y_{i}=\left\{\begin{array}{ll}
1 & \text { if } y_{i}^{*}>0  \tag{13.2}\\
0 & \text { if } y_{i}^{*} \leq 0
\end{array} .\right.
\]

To be succinct, it is useful to write equations (13.2) and (13.3) in terms of the indicator function, denoted $1[\cdot]$. Recall that this function is unity whenever the statement in brackets is true, and zero otherwise. Thus, equations (13.2) and (13.3) are\\
equivalently written as $y_{i}=1\left[y_{i}^{*}>0\right]$. Because $e_{i}$ is normally distributed, it is irrelevant whether the strict inequality is in equation (13.2) or (13.3).

We can easily obtain the distribution of $y_{i}$ given $\mathbf{x}_{i}$ :


\begin{align*}
\mathrm{P}\left(y_{i}=1 \mid \mathbf{x}_{i}\right) & =\mathrm{P}\left(y_{i}^{*}>0 \mid \mathbf{x}_{i}\right)=\mathrm{P}\left(\mathbf{x}_{i} \boldsymbol{\theta}+e_{i}>0 \mid \mathbf{x}_{i}\right) \\
& =\mathrm{P}\left(e_{i}>-\mathbf{x}_{i} \boldsymbol{\theta} \mid \mathbf{x}_{i}\right)=1-\Phi\left(-\mathbf{x}_{i} \boldsymbol{\theta}\right)=\Phi\left(\mathbf{x}_{i} \boldsymbol{\theta}\right), \tag{13.4}
\end{align*}


where $\Phi(\cdot)$ denotes the standard normal cumulative distribution function (cdf). We have used Property CD. 4 in the chapter appendix along with the symmetry of the normal distribution. Therefore,\\
$\mathrm{P}\left(y_{i}=0 \mid \mathbf{x}_{i}\right)=1-\Phi\left(\mathbf{x}_{i} \boldsymbol{\theta}\right)$.

We can combine equations (13.4) and (13.5) into the density of $y_{i}$ given $\mathbf{x}_{i}$ :\\
$f\left(y \mid \mathbf{x}_{i}\right)=\left[\Phi\left(\mathbf{x}_{i} \boldsymbol{\theta}\right)\right]^{y}\left[1-\Phi\left(\mathbf{x}_{i} \boldsymbol{\theta}\right)\right]^{1-y}, \quad y=0,1$.

That $f\left(y \mid \mathbf{x}_{i}\right)$ is zero when $y \notin\{0,1\}$ is obvious, so we will not be explicit about this in the future.

Our second example is useful when the variable to be explained takes on nonnegative integer values. Such a variable is called a count variable. We will discuss the use and interpretation of count data models in Chapter 18. For now, it suffices to note that a linear model for $\mathrm{E}(y \mid \mathbf{x})$ when $y$ takes on nonnegative integer values is not ideal because it can lead to negative predicted values. Further, since $y$ can take on the value zero with positive probability, the transformation $\log (y)$ cannot be used to obtain a model with constant elasticities or constant semielasticities. A functional form well suited for $\mathrm{E}(y \mid \mathbf{x})$ is $\exp (\mathbf{x} \boldsymbol{\theta})$. We could estimate $\boldsymbol{\theta}$ by using NLS, but all of the standard distributions for count variables imply heteroskedasticity (see Chapter 18). Thus, we can hope to do better. A traditional approach to regression models with count data is to assume that $y_{i}$ given $\mathbf{x}_{i}$ has a Poisson distribution.

Example 13.2 (Poisson Regression): Let $y_{i}$ be a nonnegative count variable; that is, $y_{i}$ can take on integer values $0,1,2, \ldots$. Denote the conditional mean of $y_{i}$ given the vector $\mathbf{x}_{i}$ as $\mathrm{E}\left(y_{i} \mid \mathbf{x}_{i}\right)=\mu\left(\mathbf{x}_{i}\right)$. A natural distribution for $y_{i}$ given $\mathbf{x}_{i}$ is the Poisson distribution:\\
$f\left(y \mid \mathbf{x}_{i}\right)=\exp \left[-\mu\left(\mathbf{x}_{i}\right)\right]\left\{\mu\left(\mathbf{x}_{i}\right)\right\}^{y} / y!, \quad y=0,1,2, \ldots$.\\
(We use $y$ as the dummy argument in the density, not to be confused with the random variable $y_{i}$.) Once we choose a form for the conditional mean function, we have\\
completely determined the distribution of $y_{i}$ given $\mathbf{x}_{i}$. For example, from equation (13.7), $\mathrm{P}\left(y_{i}=0 \mid \mathbf{x}_{i}\right)=\exp \left[-\mu\left(\mathbf{x}_{i}\right)\right]$. An important feature of the Poisson distribution is that the variance equals the mean: $\operatorname{Var}\left(y_{i} \mid \mathbf{x}_{i}\right)=\mathrm{E}\left(y_{i} \mid \mathbf{x}_{i}\right)=\mu\left(\mathbf{x}_{i}\right)$. The usual choice for $\mu(\cdot)$ is $\mu(\mathbf{x})=\exp (\mathbf{x} \boldsymbol{\theta})$, where $\boldsymbol{\theta}$ is $K \times 1$ and $\mathbf{x}$ is $1 \times K$ with first element unity.

\subsection*{13.3 General Framework for Conditional Maximum Likelihood Estimation}
Let $p_{\mathrm{o}}(\mathbf{y} \mid \mathbf{x})$ denote the conditional density of $\mathbf{y}_{i}$ given $\mathbf{x}_{i}=\mathbf{x}$, where $\mathbf{y}$ and $\mathbf{x}$ are dummy arguments. We index this density by "o" to emphasize that it is the true density of $\mathbf{y}_{i}$ given $\mathbf{x}_{i}$, and not just one of many candidates. It will be useful to let $\mathscr{X} \subset \mathbb{R}^{K}$ denote the possible values for $\mathbf{x}_{i}$ and $\mathscr{Y}$ denote the possible values of $\mathbf{y}_{i} ; \mathscr{X}$ and $\mathscr{Y}$ are called the supports of the random vectors $\mathbf{x}_{i}$ and $\mathbf{y}_{i}$, respectively.

For a general treatment, we assume that, for all $\mathbf{x} \in \mathscr{X}, p_{\mathrm{o}}(\cdot \mid \mathbf{x})$ is a density with respect to a $\sigma$-finite measure, denoted $v(d \mathbf{y})$. Defining a $\sigma$-finite measure would take us too far afield. We will say little more about the measure $v(d \mathbf{y})$ because it does not play a crucial role in applications. It suffices to know that $v(d \mathbf{y})$ can be chosen to allow $\mathbf{y}_{i}$ to be discrete, continuous, or some mixture of the two. When $\mathbf{y}_{i}$ is discrete, the measure $v(d \mathbf{y})$ simply turns all integrals into sums; when $\mathbf{y}_{i}$ is purely continuous, we obtain the usual Riemann integrals. Even in more complicated cases-where, say, $\mathbf{y}_{i}$ has both discrete and continuous characteristics-we can get by with tools from basic probability without ever explicitly defining $v(d \mathbf{y})$. For more on measures and general integrals, you are referred to Billingsley (1979) and Davidson (1994, Chaps. 3 and 4).

In Chapter 12 we saw how NLS can be motivated by the fact that $\mu_{\mathrm{o}}(\mathbf{x}) \equiv \mathrm{E}(y \mid \mathbf{x})$ minimizes $\mathrm{E}\left\{[y-m(\mathbf{x})]^{2}\right\}$ for all other functions $m(\mathbf{x})$ with $\mathrm{E}\left\{[m(\mathbf{x})]^{2}\right\}<\infty$. Conditional maximum likelihood has a similar motivation. The result from probability that is crucial for applying the analogy principle is the conditional Kullback-Leibler information inequality. Although there are more general statements of this inequality, the following suffices for our purpose: for any nonnegative function $f(\cdot \mid \mathbf{x})$ such that


\begin{equation*}
\int_{\mathscr{Y}} f(\mathbf{y} \mid \mathbf{x}) v(d \mathbf{y})=1, \quad \text { all } \mathbf{x} \in \mathscr{X} \tag{13.8}
\end{equation*}


Property CD. 1 in the chapter appendix implies that


\begin{equation*}
\mathscr{K}(f ; \mathbf{x}) \equiv \int_{\mathscr{Y}} \log \left[p_{\mathrm{o}}(\mathbf{y} \mid \mathbf{x}) / f(\mathbf{y} \mid \mathbf{x})\right] p_{\mathrm{o}}(\mathbf{y} \mid \mathbf{x}) v(d \mathbf{y}) \geq 0, \quad \text { all } \mathbf{x} \in \mathscr{X} \tag{13.9}
\end{equation*}


Because the integral is identically zero for $f=p_{\mathrm{o}}$, expression (13.9) says that, for each $\mathbf{x}, \mathscr{K}(f ; \mathbf{x})$ is minimized at $f=p_{\mathrm{o}}$.

We can apply inequality (13.9) to a parametric model for $p_{\mathrm{o}}(\cdot \mid \mathbf{x})$,


\begin{equation*}
\left\{f(\cdot \mid \mathbf{x} ; \boldsymbol{\theta}), \boldsymbol{\theta} \in \boldsymbol{\Theta}, \boldsymbol{\Theta} \subset \mathbb{R}^{\mathrm{P}}\right\} \tag{13.10}
\end{equation*}


which we assume satisfies condition (13.8) for each $\mathbf{x} \in \mathscr{X}$ and each $\boldsymbol{\theta} \in \boldsymbol{\Theta}$; if it does not, then $f(\cdot \mid \mathbf{x} ; \boldsymbol{\theta})$ does not integrate to unity (with respect to the measure $v$ ), and as a result it is a very poor candidate for $p_{\mathrm{o}}(\mathbf{y} \mid \mathbf{x})$. Model (13.10) is a correctly specified model of the conditional density, $p_{\mathrm{o}}(\cdot \mid \cdot)$, if, for some $\boldsymbol{\theta}_{\mathrm{o}} \in \boldsymbol{\Theta}$,


\begin{equation*}
f\left(\cdot \mid \mathbf{x} ; \boldsymbol{\theta}_{\mathrm{o}}\right)=p_{\mathrm{o}}(\cdot \mid \mathbf{x}), \quad \text { all } \mathbf{x} \in \mathscr{X} . \tag{13.11}
\end{equation*}


As we discussed in Chapter 12, it is useful to use $\boldsymbol{\theta}_{\mathrm{o}}$ to distinguish the true value of the parameter from a generic element of $\boldsymbol{\Theta}$. In particular examples, we will not bother making this distinction unless it is needed to make a point.

For each $\mathbf{x} \in \mathscr{X}, \mathscr{K}(f, \mathbf{x})$ can be written as $\mathrm{E}\left\{\log \left[p_{\mathrm{o}}\left(\mathbf{y}_{i} \mid \mathbf{x}_{i}\right)\right] \mid \mathbf{x}_{i}=\mathbf{x}\right\}- \mathrm{E}\left\{\log \left[f\left(\mathbf{y}_{i} \mid \mathbf{x}_{i}\right)\right] \mid \mathbf{x}_{i}=\mathbf{x}\right\}$. Therefore, if the parametric model is correctly specified, then $\mathrm{E}\left\{\log \left[f\left(\mathbf{y}_{i} \mid \mathbf{x}_{i} ; \boldsymbol{\theta}_{\mathrm{o}}\right)\right] \mid \mathbf{x}_{i}\right\} \geq \mathrm{E}\left\{\log \left[f\left(\mathbf{y}_{i} \mid \mathbf{x}_{i} ; \boldsymbol{\theta}\right)\right] \mid \mathbf{x}_{i}\right\}$, or


\begin{equation*}
\mathrm{E}\left[\ell_{i}\left(\boldsymbol{\theta}_{\mathrm{o}}\right) \mid \mathbf{x}_{i}\right] \geq \mathrm{E}\left[\ell_{i}(\boldsymbol{\theta}) \mid \mathbf{x}_{i}\right], \quad \boldsymbol{\theta} \in \boldsymbol{\Theta}, \tag{13.12}
\end{equation*}


where


\begin{equation*}
\ell_{i}(\boldsymbol{\theta}) \equiv \ell\left(\mathbf{y}_{i}, \mathbf{x}_{i}, \boldsymbol{\theta}\right) \equiv \log f\left(\mathbf{y}_{i} \mid \mathbf{x}_{i} ; \boldsymbol{\theta}\right) \tag{13.13}
\end{equation*}


is the conditional $\log$ likelihood for observation $\mathbf{i}$. Note that $\ell_{i}(\boldsymbol{\theta})$ is a random function of $\boldsymbol{\theta}$, since it depends on the random vector $\left(\mathbf{x}_{i}, \mathbf{y}_{i}\right)$. By taking the expected value of expression (13.12) and using iterated expectations, we see that $\boldsymbol{\theta}_{\mathrm{o}}$ solves


\begin{equation*}
\max _{\boldsymbol{\theta} \in \boldsymbol{\Theta}} \mathrm{E}\left[\ell_{i}(\boldsymbol{\theta})\right], \tag{13.14}
\end{equation*}


where the expectation is with respect to the joint distribution of $\left(\mathbf{x}_{i}, \mathbf{y}_{i}\right)$. The sample analogue of expression (13.14) is


\begin{equation*}
\max _{\boldsymbol{\theta} \in \boldsymbol{\Theta}} N^{-1} \sum_{i=1}^{N} \log f\left(\mathbf{y}_{i} \mid \mathbf{x}_{i} ; \boldsymbol{\theta}\right) \tag{13.15}
\end{equation*}


A solution to problem (13.15), assuming that one exists, is the conditional maximum likelihood estimator of $\boldsymbol{\theta}_{\mathrm{o}}$, which we denote as $\hat{\boldsymbol{\theta}}$. We will sometimes drop "conditional" when it is not needed for clarity.

The CMLE is clearly an M-estimator, since a maximization problem is easily turned into a minimization problem: in the notation of Chapter 12, take $\mathbf{w}_{i} \equiv\left(\mathbf{x}_{i}, \mathbf{y}_{i}\right)$\\
and $q\left(\mathbf{w}_{i}, \boldsymbol{\theta}\right) \equiv-\log f\left(\mathbf{y}_{i} \mid \mathbf{x}_{i} ; \boldsymbol{\theta}\right)$. As long as we keep track of the minus sign in front of the log likelihood, we can apply the results in Chapter 12 directly.

The motivation for the conditional MLE as a solution to problem (13.15) may appear backward if you learned about MLE in an introductory statistics course. In a traditional framework, we would treat the $\mathbf{x}_{i}$ as constants appearing in the distribution of $\mathbf{y}_{i}$, and we would define $\hat{\boldsymbol{\theta}}$ as the solution to


\begin{equation*}
\max _{\boldsymbol{\theta} \in \boldsymbol{\Theta}} \prod_{i=1}^{N} f\left(\mathbf{y}_{i} \mid \mathbf{x}_{i} ; \boldsymbol{\theta}\right) . \tag{13.16}
\end{equation*}


Under independence, the product in expression (13.16) is the model for the joint density of $\left(\mathbf{y}_{1}, \ldots, \mathbf{y}_{N}\right)$, evaluated at the data. Because maximizing the function in (13.16) is the same as maximizing its natural log, we are led to problem (13.15). However, the arguments explaining why solving (13.16) should lead to a good estimator of $\boldsymbol{\theta}_{\mathrm{o}}$ are necessarily heuristic. By contrast, the analogy principle applies directly to problem (13.15), and we need not assume that the $\mathbf{x}_{i}$ are fixed.

In our two examples, the conditional log likelihoods are fairly simple.\\
Example 13.1 (continued): In the probit example, the log likelihood for observation $i$ is $\ell_{i}(\boldsymbol{\theta})=y_{i} \log \Phi\left(\mathbf{x}_{i} \boldsymbol{\theta}\right)+\left(1-y_{i}\right) \log \left[1-\Phi\left(\mathbf{x}_{i} \boldsymbol{\theta}\right)\right]$.

Example 13.2 (continued): In the Poisson example, $\ell_{i}(\boldsymbol{\theta})=-\exp \left(\mathbf{x}_{i} \boldsymbol{\theta}\right)+y_{i} \mathbf{x}_{i} \boldsymbol{\theta}- \log \left(y_{i}!\right)$. Normally, we would drop the last term in defining $\ell_{i}(\boldsymbol{\theta})$ because it does not affect the maximization problem.

\subsection*{13.4 Consistency of Conditional Maximum Likelihood Estimation}
In this section we state a formal consistency result for the CMLE, which is a special case of the M-estimator consistency result Theorem 12.2.\\
theorem 13.1 (Consistency of CMLE): Let $\left\{\left(\mathbf{x}_{i}, \mathbf{y}_{i}\right): i=1,2, \ldots\right\}$ be a random sample with $\mathbf{x}_{i} \in \mathscr{X} \subset \mathbb{R}^{K}, \mathbf{y}_{i} \in \mathscr{Y} \subset \mathbb{R}^{G}$. Let $\boldsymbol{\Theta} \subset \mathbb{R}^{P}$ be the parameter set, and denote the parametric model of the conditional density as $\{f(\cdot \mid \mathbf{x} ; \boldsymbol{\theta}): \mathbf{x} \in \mathscr{X}, \boldsymbol{\theta} \in \boldsymbol{\Theta}\}$. Assume that (a) $f(\cdot \mid \mathbf{x} ; \boldsymbol{\theta})$ is a true density with respect to the measure $v(d \mathbf{y})$ for all $\mathbf{x}$ and $\boldsymbol{\theta}$, so that condition (13.8) holds; (b) for some $\boldsymbol{\theta}_{\mathrm{o}} \in \boldsymbol{\Theta}, p_{\mathrm{o}}(\cdot \mid \mathbf{x})=f\left(\cdot \mid \mathbf{x} ; \boldsymbol{\theta}_{\mathrm{o}}\right)$, all $\mathbf{x} \in \mathscr{X}$, and $\boldsymbol{\theta}_{\mathrm{o}}$ is the unique solution to problem (13.14); (c) $\boldsymbol{\Theta}$ is a compact set; (d) for each $\boldsymbol{\theta} \in \boldsymbol{\Theta}$, $\ell(\cdot, \boldsymbol{\theta})$ is a Borel measurable function on $\mathscr{Y} \times \mathscr{X} ;(\mathrm{e})$ for each $(\mathbf{y}, \mathbf{x}) \in \mathscr{Y} \times \mathscr{X}, \ell(\mathbf{y}, \mathbf{x}, \cdot)$ is a continuous function on $\boldsymbol{\Theta}$; and (f) $|\ell(\mathbf{w}, \boldsymbol{\theta})| \leq b(\mathbf{w})$, all $\boldsymbol{\theta} \in \boldsymbol{\Theta}$, and $\mathrm{E}[b(\mathbf{w})]<\infty$. Then there exists a solution to problem (13.15), the CMLE $\hat{\boldsymbol{\theta}}$, and $\operatorname{plim} \hat{\boldsymbol{\theta}}=\boldsymbol{\theta}_{\mathrm{o}}$.

As we discussed in Chapter 12, the measurability assumption in part d is purely technical and does not need to be checked in practice. Compactness of $\boldsymbol{\Theta}$ can be relaxed, but doing so usually requires considerable work. The continuity assumption holds in most econometric applications, but there are cases where it fails, such as when estimating certain models of auctions-see, for example, Donald and Paarsch (1996) and Paarsch and Hong (2006). The moment assumption in part f typically restricts the distribution of $\mathbf{x}_{i}$ in some way, but such restrictions are rarely a serious concern. For the most part, the key assumptions are that the parametric model is correctly specified, that $\boldsymbol{\theta}_{\mathrm{o}}$ is identified, and that the log-likelihood function is continuous in $\boldsymbol{\theta}$.

For the probit and Poisson examples, the log likelihoods are clearly continuous in $\boldsymbol{\theta}$. We can verify the moment condition (f) if we bound certain moments of $\mathbf{x}_{i}$ and make the parameter space compact. But our primary concern is that densities are correctly specified. For example, in the probit case, the density for $y_{i}$ given $\mathbf{x}_{i}$ will be incorrect if the latent error $e_{i}$ is not independent of $\mathbf{x}_{i}$ and normally distributed, or if the latent variable model is not linear to begin with. For identification we must rule out perfect collinearity in $\mathbf{x}_{i}$. The Poisson CMLE turns out to have desirable properties even if the Poisson distributional assumption does not hold, but we postpone a discussion of the robustness of the Poisson CMLE until Section 13.11 and Chapter 18.

\subsection*{13.5 Asymptotic Normality and Asymptotic Variance Estimation}
Under the differentiability and moment assumptions that allow us to apply the theorems in Chapter 12, we can show that the MLE is generally asymptotically normal. Naturally, the computational methods discussed in Section 12.7, including concentrating parameters out of the log likelihood, apply directly.

\subsection*{13.5.1 Asymptotic Normality}
We can derive the limiting distribution of the MLE by applying Theorem 12.3. We will have to assume the regularity conditions there; in particular, we assume that $\boldsymbol{\theta}_{\mathrm{o}}$ is in the interior of $\boldsymbol{\Theta}$, and $\ell_{i}(\boldsymbol{\theta})$ is twice continuously differentiable on the interior of $\boldsymbol{\Theta}$.

The score of the log likelihood for observation $i$ is simply


\begin{equation*}
\mathbf{s}_{i}(\boldsymbol{\theta}) \equiv \nabla_{\theta} \ell_{i}(\boldsymbol{\theta})^{\prime}=\left(\frac{\partial \ell_{i}}{\partial \theta_{1}}(\boldsymbol{\theta}), \frac{\partial \ell_{i}}{\partial \theta_{2}}(\boldsymbol{\theta}), \ldots, \frac{\partial \ell_{i}}{\partial \theta_{P}}(\boldsymbol{\theta})\right)^{\prime}, \tag{13.17}
\end{equation*}


a $P \times 1$ vector as in Chapter 12.

Example 13.1 (continued): For the probit case, $\boldsymbol{\theta}$ is $K \times 1$ and\\
$\nabla_{\theta} \ell_{i}(\boldsymbol{\theta})=y_{i}\left[\frac{\phi\left(\mathbf{x}_{i} \boldsymbol{\theta}\right) \mathbf{x}_{i}}{\Phi\left(\mathbf{x}_{i} \boldsymbol{\theta}\right)}\right]-\left(1-y_{i}\right)\left\{\frac{\phi\left(\mathbf{x}_{i} \boldsymbol{\theta}\right) \mathbf{x}_{i}}{\left[1-\Phi\left(\mathbf{x}_{i} \boldsymbol{\theta}\right)\right]}\right\}$.\\
Transposing this equation, and using a little algebra, gives\\
$\mathbf{s}_{i}(\boldsymbol{\theta})=\frac{\phi\left(\mathbf{x}_{i} \boldsymbol{\theta}\right) \mathbf{x}_{i}^{\prime}\left[y_{i}-\Phi\left(\mathbf{x}_{i} \boldsymbol{\theta}\right)\right]}{\Phi\left(\mathbf{x}_{i} \boldsymbol{\theta}\right)\left[1-\Phi\left(\mathbf{x}_{i} \boldsymbol{\theta}\right)\right]}$.

Recall that $\mathbf{x}_{i}^{\prime}$ is a $K \times 1$ vector.\\
Example 13.2 (continued): The score for the Poisson case, where $\boldsymbol{\theta}$ is again $K \times 1$, is\\
$\mathbf{s}_{i}(\boldsymbol{\theta})=-\exp \left(\mathbf{x}_{i} \boldsymbol{\theta}\right) \mathbf{x}_{i}^{\prime}+y_{i} \mathbf{x}_{i}^{\prime}=\mathbf{x}_{i}^{\prime}\left[y_{i}-\exp \left(\mathbf{x}_{i} \boldsymbol{\theta}\right)\right]$.

In the vast majority of cases, the score of the log-likelihood function has an important zero conditional mean property:\\
$\mathrm{E}\left[\mathbf{s}_{i}\left(\boldsymbol{\theta}_{\mathrm{o}}\right) \mid \mathbf{x}_{i}\right]=\mathbf{0}$.

In other words, when we evaluate the $P \times 1$ score at $\boldsymbol{\theta}_{\mathrm{o}}$, and take its expectation with respect to $f\left(\cdot \mid \mathbf{x}_{i} ; \boldsymbol{\theta}_{\mathrm{o}}\right)$, the expectation is zero. Under condition (13.20), $\mathrm{E}\left[\mathbf{s}_{i}\left(\boldsymbol{\theta}_{\mathrm{o}}\right)\right]=\mathbf{0}$, which was a key condition in deriving the asymptotic normality of the M-estimator in Chapter 12.

To show condition (13.20) generally, let $\mathrm{E}_{\theta}\left[\cdot \mid \mathbf{x}_{i}\right]$ denote conditional expectation with respect to the density $f\left(\cdot \mid \mathbf{x}_{i} ; \boldsymbol{\theta}\right)$ for any $\boldsymbol{\theta} \in \boldsymbol{\Theta}$. Then, by definition,\\
$\mathrm{E}_{\theta}\left[\mathbf{s}_{i}(\boldsymbol{\theta}) \mid \mathbf{x}_{i}\right]=\int_{\mathscr{Y}} \mathbf{s}\left(\mathbf{y}, \mathbf{x}_{i}, \boldsymbol{\theta}\right) f\left(\mathbf{y} \mid \mathbf{x}_{i} ; \boldsymbol{\theta}\right) v(d \mathbf{y})$.\\
If integration and differentation can be interchanged on $\operatorname{int}(\boldsymbol{\Theta})$-that is, if


\begin{equation*}
\nabla_{\theta}\left(\int_{\mathscr{Y}} f\left(\mathbf{y} \mid \mathbf{x}_{i} ; \boldsymbol{\theta}\right) v(d \mathbf{y})\right)=\int_{\mathscr{Y}} \nabla_{\theta} f\left(\mathbf{y} \mid \mathbf{x}_{i} ; \boldsymbol{\theta}\right) v(d \mathbf{y}) \tag{13.21}
\end{equation*}


for all $\mathbf{x}_{i} \in \mathscr{X}, \boldsymbol{\theta} \in \operatorname{int}(\boldsymbol{\Theta})$-then\\
$\mathbf{0}=\int_{\mathscr{Y}} \nabla_{\theta} f\left(\mathbf{y} \mid \mathbf{x}_{i} ; \boldsymbol{\theta}\right) v(d \mathbf{y})$,\\
since $\int_{\mathscr{y}} f\left(\mathbf{y} \mid \mathbf{x}_{i} ; \boldsymbol{\theta}\right) v(d \mathbf{y})$ is unity for all $\boldsymbol{\theta}$, and therefore the partial derivatives with respect to $\boldsymbol{\theta}$ must be identically zero. But the right-hand side of equation (13.22) can be written as $\int_{\mathscr{y}}\left[\nabla_{\theta} \ell\left(\mathbf{y}, \mathbf{x}_{i}, \boldsymbol{\theta}\right)\right] f\left(\mathbf{y} \mid \mathbf{x}_{i} ; \boldsymbol{\theta}\right) v(d \mathbf{y})$. Putting in $\boldsymbol{\theta}_{\mathrm{o}}$ for $\boldsymbol{\theta}$ and transposing yields condition (13.20).

Example 13.1 (continued): Define $u_{i} \equiv y_{i}-\Phi\left(\mathbf{x}_{i} \boldsymbol{\theta}_{\mathrm{o}}\right)=y_{i}-\mathrm{E}\left(y_{i} \mid \mathbf{x}_{i}\right)$. Then\\
$\mathbf{s}_{i}\left(\boldsymbol{\theta}_{\mathrm{o}}\right)=\frac{\phi\left(\mathbf{x}_{i} \boldsymbol{\theta}_{\mathrm{o}}\right) \mathbf{x}_{i}^{\prime} u_{i}}{\Phi\left(\mathbf{x}_{i} \boldsymbol{\theta}_{\mathrm{o}}\right)\left[1-\Phi\left(\mathbf{x}_{i} \boldsymbol{\theta}_{\mathrm{o}}\right)\right]}$\\
and, since $\mathrm{E}\left(u_{i} \mid \mathbf{x}_{i}\right)=0$, it follows that $\mathrm{E}\left[\mathbf{s}_{i}\left(\boldsymbol{\theta}_{\mathrm{o}}\right) \mid \mathbf{x}_{i}\right]=\mathbf{0}$.\\
Example 13.2 (continued): Define $u_{i} \equiv y_{i}-\exp \left(\mathbf{x}_{i} \boldsymbol{\theta}_{\mathrm{o}}\right)$. Then $\mathbf{s}_{i}\left(\boldsymbol{\theta}_{\mathrm{o}}\right)=\mathbf{x}_{i}^{\prime} u_{i}$ and so $\mathrm{E}\left[\mathbf{s}_{i}\left(\boldsymbol{\theta}_{\mathrm{o}}\right) \mid \mathbf{x}_{i}\right]=\mathbf{0}$.

Assuming that $\ell_{i}(\boldsymbol{\theta})$ is twice continuously differentiable on the interior of $\boldsymbol{\Theta}$, let the Hessian for observation $i$ be the $P \times P$ matrix of second partial derivatives of $\ell_{i}(\boldsymbol{\theta})$ :\\
$\mathbf{H}_{i}(\boldsymbol{\theta}) \equiv \nabla_{\theta} \mathbf{s}_{i}(\boldsymbol{\theta})=\nabla_{\theta}^{2} \ell_{i}(\boldsymbol{\theta})$.

The Hessian is a symmetric matrix that generally depends on $\left(\mathbf{x}_{i}, \mathbf{y}_{i}\right)$. Since MLE is a maximization problem, the expected value of $\mathbf{H}_{i}\left(\boldsymbol{\theta}_{\mathrm{o}}\right)$ is negative definite. Thus, to apply the theory in Chapter 12, we define\\
$\mathbf{A}_{\mathrm{o}} \equiv-\mathrm{E}\left[\mathbf{H}_{i}\left(\boldsymbol{\theta}_{\mathrm{o}}\right)\right]$,\\
which is generally a positive definite matrix when $\boldsymbol{\theta}_{\mathrm{o}}$ is identified. Under standard regularity conditions, the asymptotic normality of the CMLE follows from Theorem 12.3: $\sqrt{N}\left(\hat{\boldsymbol{\theta}}-\boldsymbol{\theta}_{\mathrm{o}}\right) \stackrel{a}{\sim} \operatorname{Normal}\left(\mathbf{0}, \mathbf{A}_{\mathrm{o}}^{-1} \mathbf{B}_{\mathrm{o}} \mathbf{A}_{\mathrm{o}}^{-1}\right)$, where $\mathbf{B}_{\mathrm{o}} \equiv \operatorname{Var}\left[\mathbf{s}_{i}\left(\boldsymbol{\theta}_{\mathrm{o}}\right)\right] \equiv \mathrm{E}\left[\mathbf{s}_{i}\left(\boldsymbol{\theta}_{\mathrm{o}}\right) \mathbf{s}_{i}\left(\boldsymbol{\theta}_{\mathrm{o}}\right)^{\prime}\right]$. It turns out that this general form of the asymptotic variance matrix is too complicated. We now show that $\mathbf{B}_{\mathrm{o}}=\mathbf{A}_{\mathrm{o}}$.

We must assume enough smoothness such that the following interchange of integral and derivative is valid (see Newey and McFadden, 1994, Sect. 5.1, for the case of unconditional MLE):\\
$\nabla_{\theta}\left(\int_{\mathscr{Y}} \mathbf{s}_{i}(\boldsymbol{\theta}) f\left(\mathbf{y} \mid \mathbf{x}_{i} ; \boldsymbol{\theta}\right) v(d \mathbf{y})\right)=\int_{\mathscr{Y}} \nabla_{\theta}\left[\mathbf{s}_{i}(\boldsymbol{\theta}) f\left(\mathbf{y} \mid \mathbf{x}_{i} ; \boldsymbol{\theta}\right)\right] v(d \mathbf{y})$.

Then, taking the derivative of the identity\\
$\int_{\mathscr{Y}} \mathbf{s}_{i}(\boldsymbol{\theta}) f\left(\mathbf{y} \mid \mathbf{x}_{i} ; \boldsymbol{\theta}\right) v(d \mathbf{y}) \equiv \mathrm{E}_{\theta}\left[\mathbf{s}_{i}(\boldsymbol{\theta}) \mid \mathbf{x}_{i}\right]=\mathbf{0}, \quad \boldsymbol{\theta} \in \operatorname{int}(\boldsymbol{\Theta})$,\\
and using equation (13.25), gives, for all $\boldsymbol{\theta} \in \operatorname{int}(\boldsymbol{\Theta})$,\\
$-\mathrm{E}_{\theta}\left[\mathbf{H}_{i}(\boldsymbol{\theta}) \mid \mathbf{x}_{i}\right]=\operatorname{Var}_{\theta}\left[\mathbf{s}_{i}(\boldsymbol{\theta}) \mid \mathbf{x}_{i}\right]$,\\
where the indexing by $\boldsymbol{\theta}$ denotes expectation and variance when $f\left(\cdot \mid \mathbf{x}_{i} ; \boldsymbol{\theta}\right)$ is the density of $\mathbf{y}_{i}$ given $\mathbf{x}_{i}$. When evaluated at $\boldsymbol{\theta}=\boldsymbol{\theta}_{\mathrm{o}}$ we get a very important equality:


\begin{equation*}
-\mathrm{E}\left[\mathbf{H}_{i}\left(\boldsymbol{\theta}_{\mathrm{o}}\right) \mid \mathbf{x}_{i}\right]=\mathrm{E}\left[\mathbf{s}_{i}\left(\boldsymbol{\theta}_{\mathrm{o}}\right) \mathbf{s}_{i}\left(\boldsymbol{\theta}_{\mathrm{o}}\right)^{\prime} \mid \mathbf{x}_{i}\right] \tag{13.26}
\end{equation*}


where the expectation and variance are with respect to the true conditional distribution of $\mathbf{y}_{i}$ given $\mathbf{x}_{i}$. Equation (13.26) is called the conditional information matrix equality (CIME). Taking the expectation of equation (13.26) (with respect to the distribution of $\mathbf{x}_{i}$ ) and using the law of iterated expectations gives


\begin{equation*}
-\mathrm{E}\left[\mathbf{H}_{i}\left(\boldsymbol{\theta}_{\mathrm{o}}\right)\right]=\mathrm{E}\left[\mathbf{s}_{i}\left(\boldsymbol{\theta}_{\mathrm{o}}\right) \mathbf{s}_{i}\left(\boldsymbol{\theta}_{\mathrm{o}}\right)^{\prime}\right], \tag{13.27}
\end{equation*}


or $\mathbf{A}_{\mathrm{o}}=\mathbf{B}_{\mathrm{o}}$. This relationship is best thought of as the unconditional information matrix equality (UIME).\\
theorem 13.2 (Asymptotic Normality of CMLE): Let the conditions of Theorem 13.1 hold. In addition, assume that (a) $\boldsymbol{\theta}_{\mathrm{o}} \in \operatorname{int}(\boldsymbol{\Theta})$; (b) for each $(\mathbf{y}, \mathbf{x}) \in \mathscr{Y} \times \mathscr{X}$, $\ell(\mathbf{y}, \mathbf{x}, \cdot)$ is twice continuously differentiable on $\operatorname{int}(\boldsymbol{\Theta})$; (c) the interchanges of derivative and integral in equations (13.21) and (13.25) hold for all $\boldsymbol{\theta} \in \operatorname{int}(\boldsymbol{\Theta})$; (d) the elements of $\nabla_{\theta}^{2} \ell(\mathbf{y}, \mathbf{x}, \theta)$ are bounded in absolute value by a function $b(\mathbf{y}, \mathbf{x})$ with finite expectation; and (e) $\mathbf{A}_{\mathrm{o}}$ defined by expression (13.24) is positive definite. Then


\begin{equation*}
\sqrt{N}\left(\hat{\boldsymbol{\theta}}-\boldsymbol{\theta}_{\mathrm{o}}\right) \xrightarrow{d} \operatorname{Normal}\left(\mathbf{0}, \mathbf{A}_{\mathrm{o}}^{-1}\right) \tag{13.28}
\end{equation*}


and therefore


\begin{equation*}
\operatorname{Avar}(\hat{\boldsymbol{\theta}})=\mathbf{A}_{\mathrm{o}}^{-1} / N . \tag{13.29}
\end{equation*}


In standard applications, the log likelihood has many continuous partial derivatives, although there are examples where it does not. Some examples also violate the interchange of the integral and derivative in equation (13.21) or (13.25), such as when the conditional support of $\mathbf{y}_{i}$ depends on the parameters $\boldsymbol{\theta}_{\mathrm{o}}$. In such cases we cannot expect the CMLE to have a limiting normal distribution; it may not even converge at the rate $\sqrt{N}$. Some progress has been made for specific models when the support of the distribution depends on unknown parameters; see, for example, Donald and Paarsch (1996).

\subsection*{13.5.2 Estimating the Asymptotic Variance}
Estimating $\operatorname{Avar}(\hat{\boldsymbol{\theta}})$ requires estimating $\mathbf{A}_{\mathrm{o}}$. From the equalities derived previously, there are at least three possible estimators of $\mathbf{A}_{\mathrm{o}}$ in the CMLE context. In fact, under slight extensions of the regularity conditions in Theorem 13.2, each of the matrices


\begin{equation*}
N^{-1} \sum_{i=1}^{N}-\mathbf{H}_{i}(\hat{\boldsymbol{\theta}}), \quad N^{-1} \sum_{i=1}^{N} \mathbf{s}_{i}(\hat{\boldsymbol{\theta}}) \mathbf{s}_{i}(\hat{\boldsymbol{\theta}})^{\prime}, \quad \text { and } \quad N^{-1} \sum_{i=1}^{N} \mathbf{A}\left(\mathbf{x}_{i}, \hat{\boldsymbol{\theta}}\right) \tag{13.30}
\end{equation*}


converges to $\mathbf{A}_{\mathrm{o}}=\mathbf{B}_{\mathrm{o}}$, where\\
$\mathbf{A}\left(\mathbf{x}_{i}, \boldsymbol{\theta}_{\mathrm{o}}\right) \equiv-\mathrm{E}\left[\mathbf{H}\left(\mathbf{y}_{i}, \mathbf{x}_{i}, \boldsymbol{\theta}_{\mathrm{o}}\right) \mid \mathbf{x}_{i}\right]$.

Thus, $\widehat{\operatorname{Avar}}(\hat{\boldsymbol{\theta}})$ can be taken to be any of the three matrices\\
$\left[-\sum_{i=1}^{N} \mathbf{H}_{i}(\hat{\boldsymbol{\theta}})\right]^{-1}, \quad\left[\sum_{i=1}^{N} \mathbf{s}_{i}(\hat{\boldsymbol{\theta}}) \mathbf{s}_{i}(\hat{\boldsymbol{\theta}})^{\prime}\right]^{-1}, \quad$ or $\quad\left[\sum_{i=1}^{N} \mathbf{A}\left(\mathbf{x}_{i}, \hat{\boldsymbol{\theta}}\right)\right]^{-1}$\\
and the asymptotic standard errors are the square roots of the diagonal elements of any of the matrices. We discussed each of these estimators in the general M-estimator case in Chapter 12, but a brief review is in order. The first estimator, based on the Hessian of the log likelihood, requires computing second derivatives. When the inverse exists, the estimate is positive definite because $\hat{\boldsymbol{\theta}}$ maximizes the objective function.

The second estimator in equation (13.32), based on the outer product of the score, depends only on first derivatives of the log-likelihood function. This simple estimator was proposed by Berndt, Hall, Hall, and Hausman (1974). Its primary drawback is that it can be poorly behaved in even moderate sample sizes, as we discussed in Section 12.6.2.

If the conditional expectation $\mathbf{A}\left(\mathbf{x}_{i}, \boldsymbol{\theta}_{\mathrm{o}}\right)$ is in closed form (as it is in some leading cases) or can be simulated-as discussed in Porter (2002)-then the estimator based on $\mathbf{A}\left(\mathbf{x}_{i}, \hat{\boldsymbol{\theta}}\right)$ has some attractive features. First, it often depends only on first derivatives of a conditional mean or conditional variance function. Second, it is also positive definite when it exists because of the conditional information matrix equality (13.26). Third, this estimator has been found to have significantly better finite-sample properties than the outer product of the score estimator in some situations where $\mathbf{A}\left(\mathbf{x}_{i}, \boldsymbol{\theta}_{\mathrm{o}}\right)$ can be obtained in closed form.

Example 13.1 (continued): The Hessian for the probit log likelihood is a mess. Fortunately, $\mathrm{E}\left[\mathbf{H}_{i}\left(\boldsymbol{\theta}_{\mathrm{o}}\right) \mid \mathbf{x}_{i}\right]$ has a fairly simple form. Taking the derivative of equation (13.18) and using the product rule gives\\
$\mathbf{H}_{i}(\boldsymbol{\theta})=-\frac{\left\{\phi\left(\mathbf{x}_{i} \boldsymbol{\theta}\right)\right\}^{2} \mathbf{x}_{i}^{\prime} \mathbf{x}_{i}}{\Phi\left(\mathbf{x}_{i} \boldsymbol{\theta}\right)\left[1-\Phi\left(\mathbf{x}_{i} \boldsymbol{\theta}\right)\right]}+\left[y_{i}-\Phi\left(\mathbf{x}_{i} \boldsymbol{\theta}\right)\right] \mathbf{L}\left(\mathbf{x}_{i} \boldsymbol{\theta}\right)$,\\
where $\mathbf{L}\left(\mathbf{x}_{i} \boldsymbol{\theta}\right)$ is a $K \times K$ complicated function of $\mathbf{x}_{i} \boldsymbol{\theta}$ that we need not find explicitly. Now, when we evaluate this expression at $\boldsymbol{\theta}_{\mathrm{o}}$ and note that $E\left\{\left[y_{i}-\Phi\left(\mathbf{x}_{i} \boldsymbol{\theta}_{\mathrm{o}}\right)\right] \mathbf{L}\left(\mathbf{x}_{i} \boldsymbol{\theta}_{\mathrm{o}}\right) \mid\right. \left.\mathbf{x}_{i}\right\}=\left[\mathrm{E}\left(y_{i} \mid \mathbf{x}_{i}\right)-\Phi\left(\mathbf{x}_{i} \boldsymbol{\theta}_{\mathrm{o}}\right)\right] \mathbf{L}\left(\mathbf{x}_{i} \boldsymbol{\theta}_{\mathrm{o}}\right)=\mathbf{0}$, we have\\
$-\mathrm{E}\left[\mathbf{H}_{i}\left(\boldsymbol{\theta}_{\mathrm{o}}\right) \mid \mathbf{x}_{i}\right]=\mathbf{A}_{i}\left(\boldsymbol{\theta}_{\mathrm{o}}\right)=\frac{\left\{\phi\left(\mathbf{x}_{i} \boldsymbol{\theta}_{\mathrm{o}}\right)\right\}^{2} \mathbf{x}_{i}^{\prime} \mathbf{x}_{i}}{\Phi\left(\mathbf{x}_{i} \boldsymbol{\theta}_{\mathrm{o}}\right)\left[1-\Phi\left(\mathbf{x}_{i} \boldsymbol{\theta}_{\mathrm{o}}\right)\right]}$.

Thus, $\widehat{\operatorname{Avar}(\hat{\boldsymbol{\theta}})}$ in probit analysis is\\
$\left(\sum_{i=1}^{N} \frac{\left\{\phi\left(\mathbf{x}_{i} \hat{\boldsymbol{\theta}}\right)\right\}^{2} \mathbf{x}_{i}^{\prime} \mathbf{x}_{i}}{\Phi\left(\mathbf{x}_{i} \hat{\boldsymbol{\theta}}\right)\left[1-\Phi\left(\mathbf{x}_{i} \hat{\boldsymbol{\theta}}\right)\right]}\right)^{-1}$,\\
which is always positive definite when the inverse exists. Note that $\mathbf{x}_{i}^{\prime} \mathbf{x}_{i}$ is a $K \times K$ matrix for each $i$.

Example 13.2 (continued): For the Poisson model with exponential conditional mean, $\mathbf{H}_{i}(\boldsymbol{\theta})=-\exp \left(\mathbf{x}_{i} \boldsymbol{\theta}\right) \mathbf{x}_{i}^{\prime} \mathbf{x}_{i}$. In this example, the Hessian does not depend on $y_{i}$, so there is no distinction between $\mathbf{H}_{i}\left(\boldsymbol{\theta}_{\mathrm{O}}\right)$ and $\mathrm{E}\left[\mathbf{H}_{i}\left(\boldsymbol{\theta}_{\mathrm{O}}\right) \mid \mathbf{x}_{i}\right]$. The positive definite estimate of $\operatorname{Avar}(\hat{\boldsymbol{\theta}})$ is simply\\
$\left[\sum_{i=1}^{N} \exp \left(\mathbf{x}_{i} \hat{\boldsymbol{\theta}}\right) \mathbf{x}_{i}^{\prime} \mathbf{x}_{i}\right]^{-1}$.

\subsection*{13.6 Hypothesis Testing}
Given the asymptotic standard errors, it is easy to form asymptotic $t$ statistics for testing single hypotheses. These $t$ statistics are asymptotically distributed as standard normal.

The three tests covered in Chapter 12 are immediately applicable to the MLE case. Since the information matrix equality holds when the density is correctly specified, we need only consider the simplest forms of the test statistics. The Wald statistic is given in equation (12.63), and the conditions sufficient for it to have a limiting chi-square distribution are discussed in Section 12.6.1.

Define the log-likelihood function for the entire sample by $\mathscr{L}(\boldsymbol{\theta}) \equiv \sum_{i=1}^{N} \ell_{i}(\boldsymbol{\theta})$. Let $\hat{\boldsymbol{\theta}}$ be the unrestricted estimator, and let $\tilde{\boldsymbol{\theta}}$ be the estimator with the $Q$ nonredundant constraints imposed. Then, under the regularity conditions discussed in Section 12.6.3, the likelihood ratio (LR) statistic,


\begin{equation*}
L R \equiv 2[\mathscr{L}(\hat{\boldsymbol{\theta}})-\mathscr{L}(\tilde{\boldsymbol{\theta}})] \tag{13.35}
\end{equation*}


is distributed asymptotically as $\chi_{Q}^{2}$ under $\mathrm{H}_{0}$. As with the Wald statistic, we cannot use $L R$ as approximately $\chi_{Q}^{2}$ when $\boldsymbol{\theta}_{\mathrm{o}}$ is on the boundary of the parameter set. The LR statistic is very easy to compute once the restricted and unrestricted models have been estimated, and the LR statistic is invariant to reparameterizing the conditional density.

The score or LM test is based on the restricted estimation only. Let $\mathbf{s}_{i}(\tilde{\boldsymbol{\theta}})$ be the $P \times 1$ score of $\ell_{i}(\boldsymbol{\theta})$ evaluated at the restricted estimates $\tilde{\boldsymbol{\theta}}$. That is, we compute the partial derivatives of $\ell_{i}(\boldsymbol{\theta})$ with respect to each of the $P$ parameters, but then we evaluate this vector of partials at the restricted estimates. Then, from Section 12.6.2 and the information matrix equality, the statistics


\begin{align*}
& \left(\sum_{i=1}^{N} \tilde{\mathbf{s}}_{i}\right)^{\prime}\left(-\sum_{i=1}^{N} \tilde{\mathbf{H}}_{i}\right)^{-1}\left(\sum_{i=1}^{N} \tilde{\mathbf{s}}_{i}\right), \quad\left(\sum_{i=1}^{N} \tilde{\mathbf{s}}_{i}\right)^{\prime}\left(\sum_{i=1}^{N} \tilde{\mathbf{A}}_{i}\right)^{-1}\left(\sum_{i=1}^{N} \tilde{\mathbf{s}}_{i}\right), \quad \text { and } \\
& \left(\sum_{i=1}^{N} \tilde{\mathbf{s}}_{i}\right)^{\prime}\left(\sum_{i=1}^{N} \tilde{\mathbf{s}}_{i} \tilde{\mathbf{s}}_{i}^{\prime}\right)^{-1}\left(\sum_{i=1}^{N} \tilde{\mathbf{s}}_{i}\right) \tag{13.36}
\end{align*}


have limiting $\chi_{Q}^{2}$ distributions under $\mathrm{H}_{0}$. As we know from Section 12.6.2, the first statistic is not guaranteed to be nonnegative (because the matrix in the middle is not necessarily positive definite) and is not invariant to reparameterizations, but the outer product statistic is. In addition, using the conditional information matrix equality, it can be shown that the LM statistic based on $\tilde{\mathbf{A}}_{i}$ is invariant to reparameterization. Davidson and MacKinnon (1993, Sect. 13.6) show invariance in the case of unconditional maximum likelihood. Invariance holds in the more general conditional ML setup, with $\mathbf{x}_{i}$ containing any conditioning variables (see Problem 13.5). We have already used the expected Hessian form of the LM statistic for nonlinear regression in Section 12.6.2. We will use it in several applications in Part IV, including binary response models and Poisson regression models. In these examples, the statistic can be computed conveniently using auxiliary regressions based on weighted residuals.

Because the unconditional information matrix equality holds, we know from Section 12.6.4 that the three classical statistics have the same limiting distribution under local alternatives. Therefore, either small-sample considerations, invariance, or computational issues must be used to choose among the statistics.

\subsection*{13.7 Specification Testing}
Because MLE generally relies on its distributional assumptions, it is useful to have available a general class of specification tests that are simple to compute. One general approach is to nest the model of interest within a more general model (which may be much harder to estimate) and obtain the score test against the more general alternative. RESET in a linear model and its extension to exponential regression models in Section 12.6.2 are examples of this approach, albeit in a non-maximum-likelihood setting.

In the context of MLE, it makes sense to test moment conditions implied by the conditional density specification. Let $\mathbf{w}_{i}=\left(\mathbf{x}_{i}, \mathbf{y}_{i}\right)$ and suppose that, when $f(\cdot \mid \mathbf{x} ; \boldsymbol{\theta})$ is correctly specified,


\begin{equation*}
\mathrm{H}_{0}: \mathrm{E}\left[\mathbf{g}\left(\mathbf{w}_{i}, \boldsymbol{\theta}_{\mathrm{o}}\right)\right]=\mathbf{0} \tag{13.37}
\end{equation*}


where $\mathbf{g}(\mathbf{w}, \boldsymbol{\theta})$ is a $Q \times 1$ vector. Any application implies innumerable choices for the function $\mathbf{g}$. Since the MLE $\hat{\boldsymbol{\theta}}$ sets the sum of the score to zero, $\mathbf{g}(\mathbf{w}, \boldsymbol{\theta})$ cannot contain elements of $\mathbf{s}(\mathbf{w}, \boldsymbol{\theta})$. Generally, $\mathbf{g}$ should be chosen to test features of a model that are of primary interest, such as first and second conditional moments, or various conditional probabilities.

A test of hypothesis (13.37) is based on how far the sample average of $\mathbf{g}\left(\mathbf{w}_{i}, \hat{\boldsymbol{\theta}}\right)$ is from zero. To derive the asymptotic distribution, note that

$$
N^{-1 / 2} \sum_{i=1}^{N} \mathbf{g}_{i}(\hat{\boldsymbol{\theta}})=N^{-1 / 2} \sum_{i=1}^{N}\left[\mathbf{g}_{i}(\hat{\boldsymbol{\theta}})-\Pi_{\mathrm{o}}^{\prime} \mathbf{s}_{i}(\hat{\boldsymbol{\theta}})\right]
$$

holds trivially because $\sum_{i=1}^{N} \mathbf{s}_{i}(\hat{\boldsymbol{\theta}})=\mathbf{0}$, where

$$
\boldsymbol{\Pi}_{\mathrm{o}} \equiv\left\{\mathrm{E}\left[\mathbf{s}_{i}\left(\boldsymbol{\theta}_{\mathrm{o}}\right) \mathbf{s}_{i}\left(\boldsymbol{\theta}_{\mathrm{o}}\right)^{\prime}\right]\right\}^{-1}\left\{\mathrm{E}\left[\mathbf{s}_{i}\left(\boldsymbol{\theta}_{\mathrm{o}}\right) \mathbf{g}_{i}\left(\boldsymbol{\theta}_{\mathrm{o}}\right)^{\prime}\right]\right\}
$$

is the $P \times Q$ matrix of population regression coefficients from regressing $\mathbf{g}_{i}\left(\boldsymbol{\theta}_{\mathrm{o}}\right)^{\prime}$ on $\mathbf{s}_{i}\left(\boldsymbol{\theta}_{\mathrm{o}}\right)^{\prime}$. Using a mean-value expansion about $\boldsymbol{\theta}_{\mathrm{o}}$ and algebra similar to that in Chapter 12, we can write


\begin{align*}
N^{-1 / 2} \sum_{i=1}^{N}\left[\mathbf{g}_{i}(\hat{\boldsymbol{\theta}})-\Pi_{\mathrm{o}}^{\prime} \mathbf{s}_{i}(\hat{\boldsymbol{\theta}})\right]= & N^{-1 / 2} \sum_{i=1}^{N}\left[\mathbf{g}_{i}\left(\boldsymbol{\theta}_{\mathrm{o}}\right)-\Pi_{\mathrm{o}}^{\prime} \mathbf{s}_{i}\left(\boldsymbol{\theta}_{\mathrm{o}}\right)\right] \\
& +\mathrm{E}\left[\nabla_{\theta} \mathbf{g}_{i}\left(\boldsymbol{\theta}_{\mathrm{o}}\right)-\Pi_{\mathrm{o}}^{\prime} \nabla_{\theta} \mathbf{s}_{i}\left(\boldsymbol{\theta}_{\mathrm{o}}\right)\right] \sqrt{N}\left(\hat{\boldsymbol{\theta}}-\boldsymbol{\theta}_{\mathrm{o}}\right)+\mathrm{o}_{p}(1) \tag{13.38}
\end{align*}


The key is that, when the density is correctly specified, the second term on the righthand side of equation (13.38) is identically zero. Here is the reason: First, equation (13.27) implies that $\left[\mathrm{E} \nabla_{\theta} \mathbf{s}_{i}\left(\boldsymbol{\theta}_{\mathrm{o}}\right)\right]\left\{\mathrm{E}\left[\mathbf{s}_{i}\left(\boldsymbol{\theta}_{\mathrm{o}}\right) \mathbf{s}_{i}\left(\boldsymbol{\theta}_{\mathrm{o}}\right)^{\prime}\right]\right\}^{-1}=-\mathbf{I}_{P}$. Second, an extension of the conditional information matrix equality (Newey, 1985; Tauchen, 1985) implies that


\begin{equation*}
-\mathrm{E}\left[\nabla_{\theta} \mathbf{g}_{i}\left(\boldsymbol{\theta}_{\mathrm{o}}\right) \mid \mathbf{x}_{i}\right]=\mathrm{E}\left[\mathbf{g}_{i}\left(\boldsymbol{\theta}_{\mathrm{o}}\right) \mathbf{s}_{i}\left(\boldsymbol{\theta}_{\mathrm{o}}\right)^{\prime} \mid \mathbf{x}_{i}\right] \tag{13.39}
\end{equation*}


To show equation (13.39), write


\begin{equation*}
\mathrm{E}_{\theta}\left[\mathbf{g}_{i}(\boldsymbol{\theta}) \mid \mathbf{x}_{i}\right]=\int_{\mathscr{Y}} \mathbf{g}\left(\mathbf{y}, \mathbf{x}_{i}, \boldsymbol{\theta}\right) f\left(\mathbf{y} \mid \mathbf{x}_{i} ; \boldsymbol{\theta}\right) v(d \mathbf{y})=\mathbf{0} \tag{13.40}
\end{equation*}


for all $\boldsymbol{\theta}$. Now, if we take the derivative with respect to $\boldsymbol{\theta}$ and assume that the integrals and derivative can be interchanged, equation (13.40) implies that\\
$\int_{\mathscr{Y}} \nabla_{\theta} \mathbf{g}\left(\mathbf{y}, \mathbf{x}_{i}, \boldsymbol{\theta}\right) f\left(\mathbf{y} \mid \mathbf{x}_{i} ; \boldsymbol{\theta}\right) v(d \mathbf{y})+\int_{\mathscr{Y}} \mathbf{g}\left(\mathbf{y}, \mathbf{x}_{i}, \boldsymbol{\theta}\right) \nabla_{\theta} f\left(\mathbf{y} \mid \mathbf{x}_{i} ; \boldsymbol{\theta}\right) v(d \mathbf{y})=\mathbf{0}$\\
or $\mathrm{E}_{\theta}\left[\nabla_{\theta} \mathbf{g}_{i}(\boldsymbol{\theta}) \mid \mathbf{x}_{i}\right]+\mathrm{E}_{\theta}\left[\mathbf{g}_{i}(\boldsymbol{\theta}) \mathbf{s}_{i}(\boldsymbol{\theta})^{\prime} \mid \mathbf{x}_{i}\right]=\mathbf{0}$, where we use the fact that $\nabla_{\theta} f(\mathbf{y} \mid \mathbf{x} ; \boldsymbol{\theta})= \mathbf{s}(\mathbf{y}, \mathbf{x}, \boldsymbol{\theta})^{\prime} f(\mathbf{y} \mid \mathbf{x} ; \boldsymbol{\theta})$. Plugging in $\boldsymbol{\theta}=\boldsymbol{\theta}_{\mathrm{o}}$ and rearranging gives equation (13.39).

What we have shown is that\\
$N^{-1 / 2} \sum_{i=1}^{N}\left[\mathbf{g}_{i}(\hat{\boldsymbol{\theta}})-\Pi_{\mathrm{o}}^{\prime} \mathbf{s}_{i}(\hat{\boldsymbol{\theta}})\right]=N^{-1 / 2} \sum_{i=1}^{N}\left[\mathbf{g}_{i}\left(\boldsymbol{\theta}_{\mathrm{o}}\right)-\Pi_{\mathrm{o}}^{\prime} \mathbf{s}_{i}\left(\boldsymbol{\theta}_{\mathrm{o}}\right)\right]+\mathrm{o}_{p}(1)$,\\
which means these standardized partial sums have the same asymptotic distribution. Letting\\
$\hat{\boldsymbol{\Pi}} \equiv\left(\sum_{i=1}^{N} \hat{\mathbf{s}}_{i} \hat{\mathbf{s}}_{i}^{\prime}\right)^{-1}\left(\sum_{i=1}^{N} \hat{\mathbf{s}}_{i} \hat{\mathbf{g}}_{i}^{\prime}\right)$,\\
it is easily seen that plim $\mathbf{\boldsymbol { \Pi }}=\boldsymbol{\Pi}_{\mathrm{o}}$ under standard regularity conditions. Therefore, the asymptotic variance of $N^{-1 / 2} \sum_{i=1}^{N}\left[\mathbf{g}_{i}(\hat{\boldsymbol{\theta}})-\boldsymbol{\Pi}_{\mathrm{o}}^{\prime} \mathbf{s}_{i}(\hat{\boldsymbol{\theta}})\right]=N^{-1 / 2} \sum_{i=1}^{N} \mathbf{g}_{i}(\hat{\boldsymbol{\theta}})$ is consistently estimated by $N^{-1} \sum_{i=1}^{N}\left(\hat{\mathbf{g}}_{i}-\hat{\boldsymbol{\Pi}}^{\prime} \hat{\mathbf{s}}_{i}\right)\left(\hat{\mathbf{g}}_{i}-\hat{\boldsymbol{\Pi}}^{\prime} \hat{\mathbf{s}}_{i}\right)^{\prime}$. When we construct the quadratic form, we get the Newey-Tauchen-White (NTW) statistic,\\
$N T W=\left[\sum_{i=1}^{N} \mathbf{g}_{i}(\hat{\boldsymbol{\theta}})\right]^{\prime}\left[\sum_{i=1}^{N}\left(\hat{\mathbf{g}}_{i}-\hat{\boldsymbol{\Pi}}^{\prime} \hat{\mathbf{s}}_{i}\right)\left(\hat{\mathbf{g}}_{i}-\hat{\boldsymbol{\Pi}}^{\prime} \hat{\mathbf{s}}_{i}\right)^{\prime}\right]^{-1}\left[\sum_{i=1}^{N} \mathbf{g}_{i}(\hat{\boldsymbol{\theta}})\right]$.

This statistic was proposed independently by Newey (1985) and Tauchen (1985), and is an extension of White's (1982a) information matrix (IM) test statistic.

For computational purposes it is useful to note that equation (13.41) is identical to $N-\mathrm{SSR}_{0}=N R_{0}^{2}$ from the regression

1 on $\hat{\mathbf{s}}_{i}^{\prime}, \hat{\mathbf{g}}_{i}^{\prime}, \quad i=1,2, \ldots, N$,\\
where $\mathrm{SSR}_{0}$ is the usual sum of squared residuals. Under the null that the density is correctly specified, $N T W$ is distributed asymptotically as $\chi_{Q}^{2}$, assuming that $\mathbf{g}(\mathbf{w}, \boldsymbol{\theta})$ contains $Q$ nonredundant moment conditions. Unfortunately, the outer product form of regression (13.42) means that the statistic can have poor finite-sample properties. In particular applications-such as nonlinear least squares, binary response analysis, and Poisson regression, to name a few-it is best to use forms of test statistics based on the expected Hessian. We gave the regression-based test for NLS in equation\\
(12.72), and we will see other examples in later chapters. For the IM test statistic, Davidson and MacKinnon (1992) have suggested an alternative form of the IM statistic that appears to have better finite-sample properties.

Example 13.2 (continued): To test the specification of the conditional mean for Poission regression, we might take $\mathbf{g}(\mathbf{w}, \boldsymbol{\theta})=\exp (\mathbf{x} \boldsymbol{\theta}) \mathbf{x}^{\prime}[y-\exp (\mathbf{x} \boldsymbol{\theta})]=\exp (\mathbf{x} \boldsymbol{\theta}) \mathbf{s}(\mathbf{w}, \boldsymbol{\theta})$, where the score is given by equation (13.19). If $\mathrm{E}(y \mid \mathbf{x})=\exp \left(\mathbf{x} \boldsymbol{\theta}_{\mathrm{o}}\right)$ then $\mathrm{E}\left[\mathbf{g}\left(\mathbf{w}, \boldsymbol{\theta}_{\mathrm{o}}\right) \mid \mathbf{x}\right] =\exp \left(\mathbf{x} \boldsymbol{\theta}_{\mathrm{o}}\right) \mathrm{E}\left[\mathbf{s}\left(\mathbf{w}, \boldsymbol{\theta}_{\mathrm{o}}\right) \mid \mathbf{x}\right]=\mathbf{0}$. To test the Poisson variance assumption, $\operatorname{Var}(y \mid \mathbf{x})= \mathrm{E}(y \mid \mathbf{x})=\exp \left(\mathbf{x} \boldsymbol{\theta}_{\mathrm{o}}\right), \mathbf{g}$ can be of the form $\mathbf{g}(\mathbf{w}, \boldsymbol{\theta})=\mathbf{a}(\mathbf{x}, \boldsymbol{\theta})\left\{[y-\exp (\mathbf{x} \boldsymbol{\theta})]^{2}-\exp (\mathbf{x} \boldsymbol{\theta})\right\}$, where $\mathbf{a}(\mathbf{x}, \boldsymbol{\theta})$ is a $Q \times 1$ vector. If the Poisson assumption is true, then $u=y- \exp \left(\mathbf{x} \boldsymbol{\theta}_{\mathrm{o}}\right)$ has a zero conditional mean and $\mathrm{E}\left(u^{2} \mid \mathbf{x}\right)=\operatorname{Var}(y \mid \mathbf{x})=\exp \left(\mathbf{x} \boldsymbol{\theta}_{\mathrm{o}}\right)$. It follows that $\mathrm{E}\left[\mathbf{g}\left(\mathbf{w}, \boldsymbol{\theta}_{\mathrm{o}}\right) \mid \mathbf{x}\right]=\mathbf{0}$.

Example 13.2 contains examples of what are known as conditional moment tests. As the name suggests, the idea is to form orthogonality conditions based on some key conditional moments, usually the conditional mean or conditional variance, but sometimes conditional probabilities or higher order moments. The tests for nonlinear regression in Chapter 12 can be viewed as conditional moment tests, and we will see several other examples in Part IV. For reasons discussed earlier, we will avoid computing the tests using regression (13.42) whenever possible. See Newey (1985), Tauchen (1985), and Pagan and Vella (1989) for general treatments and applications of conditional moment tests. White's (1982a) IM test can often be viewed as a conditional moment test; see Hall (1987) for the linear regression model and White (1994) for a general treatment. White (1994, Chap. 10) shows how to allow the moment function to depend on parameters other than $\boldsymbol{\theta}_{\mathrm{o}}$.

\subsection*{13.8 Partial (or Pooled) Likelihood Methods for Panel Data}
Up to this point we have assumed that the parametric model for the density of $\mathbf{y}$ given $\mathbf{x}$ is correctly specified. This assumption is fairly general because $\mathbf{x}$ can contain any observable variable. The leading case occurs when $\mathbf{x}$ contains variables we view as exogenous in a structural model. In other cases, $\mathbf{x}$ will contain variables that are endogenous in a structural model, but putting them in the conditioning set and finding the new conditional density makes estimation of the structural parameters easier.

For studying various panel data models, for estimation using cluster samples, and for various other applications, we need to relax the assumption that the full conditional density of $\mathbf{y}$ given $\mathbf{x}$ is correctly specified. In some examples, such a model is too complicated. Or, for robustness reasons, we do not wish to fully specify the density of $\mathbf{y}$ given $\mathbf{x}$.

\subsection*{13.8.1 Setup for Panel Data}
For panel data applications we let $\mathbf{y}$ denote a $T \times 1$ vector, with generic element $y_{t}$. Thus, $\mathbf{y}_{i}$ is a $T \times 1$ random draw vector from the cross section, with $t$ th element $y_{i t}$. As always, we are thinking of $T$ small relative to the cross section sample size. With a slight notational change we can replace $y_{i t}$ with, say, a $G$-vector for each $t$, an extension that allows us to cover general systems of equations with panel data.

For some vector $\mathbf{x}_{t}$ containing any set of observable variables, let $\mathrm{D}\left(y_{t} \mid \mathbf{x}_{t}\right)$ denote the distribution of $y_{t}$ given $\mathbf{x}_{t}$. The key assumption is that we have a correctly specified model for the density of $y_{t}$ given $\mathbf{x}_{t}$; call it $f_{t}\left(y_{t} \mid \mathbf{x}_{t} ; \boldsymbol{\theta}\right), t=1,2 \ldots, T$. The vector $\mathbf{x}_{t}$ can contain anything, including conditioning variables $\mathbf{z}_{t}$, lags of these, and lagged values of $y$. The vector $\boldsymbol{\theta}$ consists of all parameters appearing in $f_{t}$ for any $t$; some or all of these may appear in the density for every $t$, and some may appear only in the density for a single time period.

What distinguishes partial likelihood from maximum likelihood is that we do not assume that


\begin{equation*}
\prod_{t=1}^{T} \mathrm{D}\left(y_{i t} \mid \mathbf{x}_{i t}\right) \tag{13.43}
\end{equation*}


is a conditional distribution of the vector $\mathbf{y}_{i}$ given some set of conditioning variables. In other words, even though $f_{t}\left(y_{t} \mid \mathbf{x}_{t} ; \boldsymbol{\theta}_{\mathrm{o}}\right)$ is the correct density for $y_{i t}$ given $\mathbf{x}_{i t}=\mathbf{x}_{t}$ for each $t$, the product of these is not (necessarily) the density of $\mathbf{y}_{i}$ given some conditioning variables. Usually, we specify $f_{t}\left(y_{t} \mid \mathbf{x}_{t} ; \boldsymbol{\theta}\right)$ because it is the density of interest for each $t$.

We define the partial log likelihood for each observation $i$ as\\
$\ell_{i}(\boldsymbol{\theta}) \equiv \sum_{t=1}^{T} \log f_{t}\left(y_{i t} \mid \mathbf{x}_{i t} ; \boldsymbol{\theta}\right)$,\\
which is the sum of the log likelihoods across $t$. What makes partial likelihood methods work is that $\boldsymbol{\theta}_{\mathrm{o}}$ maximizes the expected value of equation (13.44) provided we have the densities $f_{t}\left(y_{t} \mid \mathbf{x}_{t} ; \boldsymbol{\theta}\right)$ correctly specified. We also refer to equation (13.44) as a pooled log likelihood.

By the Kullback-Leibler information inequality, $\boldsymbol{\theta}_{\mathrm{o}}$ maximizes $\mathrm{E}\left[\log f_{t}\left(y_{i t} \mid \mathbf{x}_{i t} ; \boldsymbol{\theta}\right)\right]$ over $\boldsymbol{\Theta}$ for each $t$, so $\boldsymbol{\theta}_{\mathrm{o}}$ also maximizes the sum of these over $t$. As usual, identification requires that $\boldsymbol{\theta}_{\mathrm{o}}$ be the unique maximizer of the expected value of equation (13.44). It is sufficient that $\boldsymbol{\theta}_{\mathrm{o}}$ uniquely maximizes $\mathrm{E}\left[\log f_{t}\left(y_{i t} \mid \mathbf{x}_{i t} ; \boldsymbol{\theta}\right)\right]$ for each $t$, but this assumption is not necessary.

The partial (pooled) maximum likelihood estimator (PMLE) $\hat{\boldsymbol{\theta}}$ solves


\begin{equation*}
\max _{\boldsymbol{\theta} \in \boldsymbol{\Theta}} \sum_{i=1}^{N} \sum_{t=1}^{T} \log f_{t}\left(y_{i t} \mid \mathbf{x}_{i t} ; \boldsymbol{\theta}\right) \tag{13.45}
\end{equation*}


which is clearly an M-estimator problem (where the asymptotics are with fixed $T$ and $N \rightarrow \infty)$. Therefore, from Theorem 12.2, the partial MLE is generally consistent provided $\boldsymbol{\theta}_{\mathrm{o}}$ is identified.

It is also clear that the partial MLE will be asymptotically normal by Theorem 12.3 in Section 12.3. However, unless


\begin{equation*}
p_{\mathrm{o}}(\mathbf{y} \mid \mathbf{z})=\prod_{t=1}^{T} f_{t}\left(y_{t} \mid \mathbf{x}_{t} ; \boldsymbol{\theta}_{\mathrm{o}}\right) \tag{13.46}
\end{equation*}


for some subvector $\mathbf{z}$ of $\mathbf{x}$, we cannot apply the CIME. A more general asymptotic variance estimator of the type covered in Section 12.5.1 is needed, and we provide such estimators in the next two subsections.

It is useful to discuss at a general level why equation (13.46) does not necessarily hold in a panel data setting. First, suppose $\mathbf{x}_{t}$ contains only contemporaneous conditioning variables, $\mathbf{z}_{t}$; in particular, $\mathbf{x}_{t}$ contains no lagged dependent variables. Then we can always write

$$
\begin{aligned}
p_{\mathrm{o}}(\mathbf{y} \mid \mathbf{z})= & p_{1}^{\mathrm{o}}\left(y_{1} \mid \mathbf{z}\right) \cdot p_{2}^{\mathrm{o}}\left(y_{2} \mid y_{1}, \mathbf{z}\right) \cdots p_{t}^{\mathrm{o}}\left(y_{t} \mid y_{t-1}, y_{t-2}, \ldots, y_{1}, \mathbf{z}\right) \cdots \\
& p_{T}^{\mathrm{o}}\left(y_{T} \mid y_{T-1}, y_{T-2}, \ldots, y_{1}, \mathbf{z}\right)
\end{aligned}
$$

where $p_{t}^{\mathrm{o}}\left(y_{t} \mid y_{t-1}, y_{t-2}, \ldots, y_{1}, \mathbf{z}\right)$ is the true conditional density of $y_{t}$ given $y_{t-1}$, $y_{t-2}, \ldots, y_{1}$ and $\mathbf{z} \equiv\left(\mathbf{z}_{1}, \ldots, \mathbf{z}_{T}\right)$. (For $t=1, p_{1}^{o}$ is the density of $y_{1}$ given $\mathbf{z}$.) For equation (13.46) to hold, we should have\\
$p_{t}^{\mathrm{o}}\left(y_{t} \mid y_{t-1}, y_{t-2}, \ldots, y_{1}, \mathbf{z}\right)=f_{t}\left(y_{t} \mid \mathbf{z}_{t} ; \boldsymbol{\theta}_{\mathrm{o}}\right), \quad t=1, \ldots, T$,\\
which requires that, once $\mathbf{z}_{t}$ is conditioned on, neither past lags of $y_{t}$ nor elements of $\mathbf{z}$ from any other time period-past or future-appear in the conditional density $p_{t}^{\mathrm{o}}\left(y_{t} \mid y_{t-1}, y_{t-2}, \ldots, y_{1}, \mathbf{z}\right)$. Generally, this requirement is very strong, as it requires a combination of strict exogeneity of $\mathbf{z}_{t}$ and the absense of dynamics in $p_{t}^{\circ}$.

Equation (13.46) is more likely to hold when $\mathbf{x}_{t}$ contains lagged dependent variables. In fact, if $\mathbf{x}_{t}$ contains only lagged values of $y_{t}$, then\\
$p_{\mathrm{o}}(\mathbf{y})=\prod_{t=1}^{T} f_{t}\left(y_{t} \mid \mathbf{x}_{t} ; \boldsymbol{\theta}_{\mathrm{o}}\right)$\\
holds if $f_{t}\left(y_{t} \mid \mathbf{x}_{t} ; \boldsymbol{\theta}_{\mathrm{o}}\right)=p_{t}^{\mathrm{o}}\left(y_{t} \mid y_{t-1}, y_{t-2}, \ldots, y_{1}\right)$ for all $t$ (where $p_{1}^{\mathrm{o}}$ is the unconditional density of $y_{1}$ ), so that all dynamics are captured by $f_{t}$. When $\mathbf{x}_{t}$ contains some variables $\mathbf{z}_{t}$ in addition to lagged $y_{t}$, equation (13.46) requires that the parametric density captures all of the dynamics - that is, that all lags of $y_{t}$ and $\mathbf{z}_{t}$ have been properly accounted for in $f\left(y_{t} \mid \mathbf{x}_{t}, \boldsymbol{\theta}_{\mathrm{o}}\right)$-and strict exogeneity of $\mathbf{z}_{t}$.

In most treatments of MLE of dynamic models containing additional exogenous variables, the strict exogeneity assumption is maintained, often implicitly by taking $\mathbf{z}_{t}$ to be nonrandom. In Chapter 7 we saw that strict exogeneity played no role in getting consistent, asymptotically normal estimators in linear panel data models using pooled OLS, and the same is true here. We also allow models where the dynamics have been incompletely specified.

With partial MLE we are interested in fully specifying a density for the conditional distribution $\mathrm{D}\left(y_{i t} \mid \mathbf{x}_{i t}\right)$, and it is useful to have a general yet simple way to define strict exogeneity of $\left\{\mathbf{x}_{i t}: t=1, \ldots, T\right\}$. The definition is simply stated: the conditioning variables are strictly exogenous if $\mathrm{D}\left(y_{i t} \mid \mathbf{x}_{i 1}, \mathbf{x}_{i 2}, \ldots, \mathbf{x}_{i T}\right)=\mathrm{D}\left(y_{i t} \mid \mathbf{x}_{i t}\right)$ for $t=1, \ldots, T$. Naturally, this condition must fail if $\mathbf{x}_{i t}$ contains lags of $y_{i t}$, and, just as in linear models, it can fail if $\mathbf{x}_{i t}$ contains elements $\mathbf{z}_{i t}$ whose future values react to unpredictable changes in $y_{i t}$.

Example 13.3 (Probit with Panel Data): To illustrate the previous discussion, we consider estimation of a panel data binary choice model. The idea is that, for each unit $i$ in the population (individual, firm, and so on) we have a binary outcome, $y_{i t}$, for each of $T$ time periods. For example, if $t$ represents a year, then $y_{i t}$ might indicate whether a person was arrested for a crime during year $t$.

Consider the model in latent variable form:


\begin{align*}
& y_{i t}^{*}=\mathbf{x}_{i t} \boldsymbol{\theta}_{\mathrm{o}}+e_{i t} \\
& y_{i t}=1\left[y_{i t}^{*}>0\right]  \tag{13.47}\\
& e_{i t} \mid \mathbf{x}_{i t} \sim \operatorname{Normal}(0,1)
\end{align*}


The vector $\mathbf{x}_{i t}$ might contain exogenous variables $\mathbf{z}_{i t}$, lags of these, and even lagged $y_{i t}$ (not lagged $y_{i t}^{*}$ ). Under the assumptions in model (13.47), we have, for each $t, P\left(y_{i t}=1 \mid \mathbf{x}_{i t}\right)=\Phi\left(\mathbf{x}_{i t} \boldsymbol{\theta}_{\mathrm{o}}\right)$, and the density of $y_{i t}$ given $\mathbf{x}_{i t}=\mathbf{x}_{t}$ is $f\left(y_{t} \mid \mathbf{x}_{t}\right)= \left[\Phi\left(\mathbf{x}_{t} \boldsymbol{\theta}_{\mathrm{o}}\right)\right]^{y_{t}}\left[1-\Phi\left(\mathbf{x}_{t} \boldsymbol{\theta}_{\mathrm{o}}\right)\right]^{1-y_{t}}$.

The partial log likelihood for a cross section observation $i$ is


\begin{equation*}
\ell_{i}(\boldsymbol{\theta})=\sum_{t=1}^{T}\left\{y_{i t} \log \Phi\left(\mathbf{x}_{i t} \boldsymbol{\theta}\right)+\left(1-y_{i t}\right) \log \left[1-\Phi\left(\mathbf{x}_{i t} \boldsymbol{\theta}\right)\right]\right\} \tag{13.48}
\end{equation*}


and the partial MLE in this case-which simply maximizes $\ell_{i}(\boldsymbol{\theta})$ summed across all $i$-is the pooled probit estimator. With $T$ fixed and $N \rightarrow \infty$, this estimator is consistent and $\sqrt{N}$-asymptotically normal without any assumptions other than identification and standard regularity conditions.

It is very important to know that the pooled probit estimator works without imposing additional assumptions on $\mathbf{e}_{i}=\left(e_{i 1}, \ldots, e_{i T}\right)^{\prime}$. When $\mathbf{x}_{i t}$ contains only exogenous variables $\mathbf{z}_{i t}$, it would be standard to assume that\\
$e_{i t}$ is independent of $\mathbf{z}_{i} \equiv\left(\mathbf{z}_{i 1}, \mathbf{z}_{i 2}, \ldots, \mathbf{z}_{i T}\right), \quad t=1, \ldots, T$.

This is the natural strict exogeneity assumption (and is much stronger than simply assuming that $e_{i t}$ and $\mathbf{z}_{i t}$ are independent for each $t$ ). The crime example can illustrate how strict exogeneity might fail. For example, suppose that $\mathbf{z}_{i t}$ measures the amount of time the person has spent in prison prior to the current year. An arrest this year $\left(y_{i t}=1\right)$ certainly has an effect on expected future values of $\mathbf{z}_{i t}$, so that assumption (13.49) is almost certainly false. Fortunately, we do not need assumption (13.49) to apply partial likelihood methods.

A second standard assumption is that the $e_{i t}, t=1,2, \ldots, T$ are serially independent. This is especially restrictive in a static model. If we maintain this assumption in addition to assumption (13.49), then equation (13.46) holds (because the $y_{i t}$ are then independent conditional on $\mathbf{z}_{i}$ ) and the partial MLE is a conditional MLE.

To relax the assumption that the $y_{i t}$ are conditionally independent, we can allow the $e_{i t}$ to be correlated across $t$ (still assuming that no lagged dependent variables appear). A common assumption is that $\mathbf{e}_{i}$ has a multivariate normal distribution with a general correlation matrix. Under this assumption, we can write down the joint distribution of $\mathbf{y}_{i}$ given $\mathbf{z}_{i}$, but it is complicated, and estimation is very computationally intensive (for discussions see Keane, 1993, and Hajivassilou and Ruud, 1994). We will cover a special case, the random effects probit model, in Chapter 15.

A nice feature of the partial MLE is that $\hat{\boldsymbol{\theta}}$ will be consistent and asymptotically normal even if the $e_{i t}$ are arbitrarily serially correlated. This result is entirely analogous to using pooled OLS in linear panel data models when the errors have arbitrary serial correlation.

When $\mathbf{x}_{i t}$ contains lagged dependent variables, model (13.47) provides a way of examining dynamic behavior. Or, perhaps $y_{i, t-1}$ is included in $\mathbf{x}_{i t}$ as a proxy for unobserved factors, and our focus is on policy variables in $\mathbf{z}_{i t}$. For example, if $y_{i t}$ is a binary indicator of employment, $y_{i, t-1}$ might be included as a control when studying the effect of a job training program (which may be a binary element of $\mathbf{z}_{i t}$ ) on the employment probability; this method controls for the fact that participation in job training this year might depend on employment last year, and it captures the fact that\\
employment status is persistent. In any case, provided $\mathrm{P}\left(y_{i t}=1 \mid \mathbf{x}_{i t}\right)$ follows a probit, the pooled probit estimator is consistent and asymptotically normal. The dynamics may or may not be correctly specified (more on this topic later), and the $\mathbf{z}_{i t}$ need not be strictly exogenous (so that whether someone participates in job training in year $t$ can depend on the past employment history).

\subsection*{13.8.2 Asymptotic Inference}
The most important practical difference between conditional MLE and partial MLE is in the computation of asymptotic standard errors and test statistics. In many cases, including the pooled probit estimator, the pooled Poisson estimator (see Problem 13.6), and many other pooled procedures, standard econometrics packages can be used to compute the partial MLEs. However, except under certain assumptions, the usual standard errors and test statistics reported from a pooled analysis are not valid. This situation is entirely analogous to the linear model case in Section 7.8 when the errors are serially correlated.

Estimation of the asymptotic variance of the partial MLE is not difficult. In fact, we can combine the M-estimation results from Section 12.5.1 and the results of Section 13.5 to obtain valid estimators.

From Theorem 12.3, we have $\operatorname{Avar} \sqrt{N}\left(\hat{\boldsymbol{\theta}}-\boldsymbol{\theta}_{\mathrm{o}}\right)=\mathbf{A}_{\mathrm{o}}^{-1} \mathbf{B}_{\mathrm{o}} \mathbf{A}_{\mathrm{o}}^{-1}$, where\\
$\mathbf{A}_{\mathrm{o}}=-\mathrm{E}\left[\nabla_{\theta}^{2} \ell_{i}\left(\boldsymbol{\theta}_{\mathrm{o}}\right)\right]=-\sum_{t=1}^{T} \mathrm{E}\left[\nabla_{\theta}^{2} \ell_{i t}\left(\boldsymbol{\theta}_{\mathrm{o}}\right)\right]=\sum_{t=1}^{T} \mathrm{E}\left[\mathbf{A}_{i t}\left(\boldsymbol{\theta}_{\mathrm{o}}\right)\right]$,\\
$\mathbf{B}_{\mathrm{o}}=\mathrm{E}\left[\mathbf{s}_{i}\left(\boldsymbol{\theta}_{\mathrm{o}}\right) \mathbf{s}_{i}\left(\boldsymbol{\theta}_{\mathrm{o}}\right)^{\prime}\right]=\mathrm{E}\left\{\left[\sum_{t=1}^{T} \mathbf{s}_{i t}\left(\boldsymbol{\theta}_{\mathrm{o}}\right)\right]\left[\sum_{t=1}^{T} \mathbf{s}_{i t}\left(\boldsymbol{\theta}_{\mathrm{o}}\right)\right]^{\prime}\right\}$,\\
$\mathbf{A}_{i t}\left(\boldsymbol{\theta}_{\mathrm{o}}\right)=-\mathrm{E}\left[\nabla_{\theta}^{2} \ell_{i t}\left(\boldsymbol{\theta}_{\mathrm{o}}\right) \mid \mathbf{x}_{i t}\right], \quad$ and\\
$\mathbf{s}_{i t}(\boldsymbol{\theta})=\nabla_{\theta} \ell_{i t}(\boldsymbol{\theta})^{\prime}$.\\
Because we assume that $\boldsymbol{\theta}_{\mathrm{o}}$ is in the interior of the parameter space, and $\boldsymbol{\theta}_{\mathrm{o}}$ maximizes $\mathrm{E}\left[\boldsymbol{l}_{i t}(\boldsymbol{\theta}) \mid \mathbf{x}_{i t}\right]$ over $\boldsymbol{\Theta}$, it is generally true that $\mathrm{E}\left[\mathbf{s}_{i t}\left(\boldsymbol{\theta}_{\mathrm{o}}\right) \mid \mathbf{x}_{i t}\right]=\mathbf{0}$ for all $t$. If $\left\{\mathbf{x}_{i t}\right\}$ is strictly exogenous, then $\mathrm{E}\left[\mathbf{s}_{i t}\left(\boldsymbol{\theta}_{\mathrm{o}}\right) \mid \mathbf{x}_{i}\right]=\mathbf{0}$ (because $\mathrm{D}\left(y_{i t} \mid \mathbf{x}_{i}\right)=\mathrm{D}\left(y_{i t} \mid \mathbf{x}_{i t}\right)$ ), but without strict exogeneity we can only say that $\mathbf{s}_{i t}\left(\boldsymbol{\theta}_{\mathrm{o}}\right)$ has zero mean conditional on $\mathbf{x}_{i t}$ (which implies $\mathrm{E}\left[\mathbf{s}_{i t}\left(\boldsymbol{\theta}_{\mathrm{o}}\right)\right]=\mathbf{0}$, of course). The natural definition of sequential exogeneity in this context is $\mathrm{D}\left(y_{i t} \mid \mathbf{x}_{i t}, \mathbf{x}_{i, t-1}, \ldots, \mathbf{x}_{i 1}\right)=\mathrm{D}\left(y_{i t} \mid \mathbf{x}_{i t}\right)$, in which case $\mathrm{E}\left[\mathbf{s}_{i t}\left(\boldsymbol{\theta}_{\mathrm{o}}\right) \mid \mathbf{x}_{i t}, \mathbf{x}_{i, t-1}, \ldots, \mathbf{x}_{i 1}\right]=\mathbf{0}$. As we will see in the next subsection, sequential exogeneity ensures that the scores are serially uncorrelated if $\mathbf{x}_{i t}$ contains $y_{i, t-1}$. If $\mathbf{x}_{i t}$\\
only includes contemporaneous variables $\mathbf{z}_{i t}$, lags of such variables, or both, sequential exogeneity does not imply the scores are serially uncorrelated.

The matrix $\mathbf{A}_{\mathrm{o}}$ is just the sum across $t$ of minus the expected Hessian. The matrix $\mathbf{B}_{\mathrm{o}}$ generally depends on the correlation between the scores at different time periods: $\mathrm{E}\left[\mathbf{s}_{i t}\left(\boldsymbol{\theta}_{\mathrm{o}}\right) \mathbf{s}_{i r}\left(\boldsymbol{\theta}_{\mathrm{o}}\right)^{\prime}\right], t \neq r$. For each $t$, the CIME holds:\\
$\mathbf{A}_{i t}\left(\boldsymbol{\theta}_{\mathrm{o}}\right)=\mathrm{E}\left[\mathbf{s}_{i t}\left(\boldsymbol{\theta}_{\mathrm{o}}\right) \mathbf{s}_{i t}\left(\boldsymbol{\theta}_{\mathrm{o}}\right)^{\prime} \mid \mathbf{x}_{i t}\right]$.\\
If $\left\{\mathbf{x}_{i t}\right\}$ is strictly exogenous then $\mathbf{A}_{i t}\left(\boldsymbol{\theta}_{\mathrm{o}}\right)=\mathrm{E}\left[\mathbf{s}_{i t}\left(\boldsymbol{\theta}_{\mathrm{o}}\right) \mathbf{s}_{i t}\left(\boldsymbol{\theta}_{\mathrm{o}}\right)^{\prime} \mid \mathbf{x}_{i t}\right]$. But strict exogeneity does not help simplify inference because, whether or not strict exogeneity holds, $-\mathrm{E}\left[\mathbf{H}_{i}\left(\boldsymbol{\theta}_{\mathrm{o}}\right) \mid \mathbf{x}_{i}\right]=\mathrm{E}\left[\mathbf{s}_{i}\left(\boldsymbol{\theta}_{\mathrm{o}}\right) \mathbf{s}_{i}\left(\boldsymbol{\theta}_{\mathrm{o}}\right)^{\prime} \mid \mathbf{x}_{i}\right]$ likely fails if the scores are serially correlated. More important, serial correlation in the scores causes $\mathbf{B}_{\mathrm{o}} \neq \mathbf{A}_{\mathrm{o}}$. Thus, to perform inference in the context of partial MLE, we generally need separate estimates of $\mathbf{A}_{\mathrm{o}}$ and $\mathbf{B}_{\mathrm{o}}$. Given the structure of the partial MLE, these are easy to obtain. Three possibilities for $\mathbf{A}_{\circ}$ are


\begin{align*}
& N^{-1} \sum_{i=1}^{N} \sum_{t=1}^{T}-\nabla_{\theta}^{2} \ell_{i t}(\hat{\boldsymbol{\theta}}), \quad N^{-1} \sum_{i=1}^{N} \sum_{t=1}^{T} \mathbf{A}_{i t}(\hat{\boldsymbol{\theta}}), \quad \text { and } \\
& N^{-1} \sum_{i=1}^{N} \sum_{t=1}^{T} \mathbf{s}_{i t}(\hat{\boldsymbol{\theta}}) \mathbf{s}_{i t}(\hat{\boldsymbol{\theta}})^{\prime} . \tag{13.50}
\end{align*}


The validity of the second of these follows from a standard iterated expectations argument, and the last of these follows from the CIME for each $t$. In most cases, the second estimator is preferred when it is easy to compute.

Since $\mathbf{B}_{\mathrm{o}}$ depends on $\mathrm{E}\left[\mathbf{s}_{i t}\left(\boldsymbol{\theta}_{\mathrm{o}}\right) \mathbf{s}_{i t}\left(\boldsymbol{\theta}_{\mathrm{o}}\right)^{\prime}\right]$ as well as on cross product terms, there are also at least three estimators available for $\mathbf{B}_{\mathrm{o}}$. The simplest is


\begin{equation*}
N^{-1} \sum_{i=1}^{N} \hat{\mathbf{s}}_{i} \hat{\mathbf{s}}_{i}^{\prime}=N^{-1} \sum_{i=1}^{N} \sum_{t=1}^{T} \hat{\mathbf{s}}_{i t} \hat{\mathbf{s}}_{i t}^{\prime}+N^{-1} \sum_{i=1}^{N} \sum_{t=1}^{T} \sum_{r \neq t} \hat{\mathbf{s}}_{i r} \hat{\mathbf{s}}_{i t}^{\prime}, \tag{13.51}
\end{equation*}


where the second term on the right-hand side accounts for possible serial correlation in the score. The first term on the right-hand side of equation (13.51) can be replaced by one of the other two estimators in equation (13.50). The asymptotic variance of $\hat{\boldsymbol{\theta}}$ is estimated, as usual, by $\hat{\mathbf{A}}^{-1} \hat{\mathbf{B}} \hat{\mathbf{A}}^{-1} / N$ for the chosen estimators $\hat{\mathbf{A}}$ and $\hat{\mathbf{B}}$. The asymptotic standard errors come directly from this matrix, and Wald tests for linear and nonlinear hypotheses can be obtained directly. The robust score statistic discussed in Section 12.6.2 can also be used. When $\mathbf{B}_{\mathrm{o}} \neq \mathbf{A}_{\mathrm{o}}$, the likelihood ratio statistic computed after pooled estimation is not valid.

Because the CIME holds for each $t, \mathbf{B}_{\mathrm{o}}=\mathbf{A}_{\mathrm{o}}$ when the scores evaluated at $\boldsymbol{\theta}_{\mathrm{o}}$ are serially uncorrelated, that is, when


\begin{equation*}
\mathrm{E}\left[\mathbf{s}_{i t}\left(\boldsymbol{\theta}_{\mathrm{o}}\right) \mathbf{s}_{i r}\left(\boldsymbol{\theta}_{\mathrm{o}}\right)^{\prime}\right]=\mathbf{0}, \quad t \neq r . \tag{13.52}
\end{equation*}


When the score is serially uncorrelated, inference is very easy: the usual MLE statistics computed from the pooled estimation, including likelihood ratio statistics, are asymptotically valid. Effectively, we can ignore the fact that a time dimension is present. The estimator of $\operatorname{Avar}(\hat{\boldsymbol{\theta}})$ is just $\hat{\mathbf{A}}^{-1} / N$, where $\hat{\mathbf{A}}$ is one of the matrices in equation (13.50).

Example 13.3 (continued): For the pooled probit example, a simple, general estimator of the asymptotic variance is


\begin{equation*}
\left[\sum_{i=1}^{N} \sum_{t=1}^{T} \mathbf{A}_{i t}(\hat{\boldsymbol{\theta}})\right]^{-1}\left[\sum_{i=1}^{N} \mathbf{s}_{i}(\hat{\boldsymbol{\theta}}) \mathbf{s}_{i}(\hat{\boldsymbol{\theta}})^{\prime}\right]\left[\sum_{i=1}^{N} \sum_{t=1}^{T} \mathbf{A}_{i t}(\hat{\boldsymbol{\theta}})\right]^{-1}, \tag{13.53}
\end{equation*}


where

$$
\mathbf{A}_{i t}(\hat{\boldsymbol{\theta}})=\frac{\left\{\phi\left(\mathbf{x}_{i t} \hat{\boldsymbol{\theta}}\right)\right\}^{2} \mathbf{x}_{i t}^{\prime} \mathbf{x}_{i t}}{\Phi\left(\mathbf{x}_{i t} \hat{\boldsymbol{\theta}}\right)\left[1-\Phi\left(\mathbf{x}_{i t} \hat{\boldsymbol{\theta}}\right)\right]}
$$

and

$$
\mathbf{s}_{i}(\boldsymbol{\theta})=\sum_{t=1}^{T} \mathbf{s}_{i t}(\boldsymbol{\theta})=\sum_{t=1}^{T} \frac{\phi\left(\mathbf{x}_{i t} \boldsymbol{\theta}\right) \mathbf{x}_{i t}^{\prime}\left[y_{i t}-\Phi\left(\mathbf{x}_{i t} \boldsymbol{\theta}\right)\right]}{\Phi\left(\mathbf{x}_{i t} \boldsymbol{\theta}\right)\left[1-\Phi\left(\mathbf{x}_{i t} \boldsymbol{\theta}\right)\right]} .
$$

The estimator (13.53) contains cross product terms of the form $\mathbf{s}_{i t}(\hat{\boldsymbol{\theta}}) \mathbf{s}_{i r}(\hat{\boldsymbol{\theta}})^{\prime}, t \neq r$, and so it is fully robust. If the score is serially uncorrelated, then the usual probit standard errors and test statistics from the pooled estimation are valid. We will discuss a sufficient condition for the scores to be serially uncorrelated in the next subsection.

\subsection*{13.8.3 Inference with Dynamically Complete Models}
There is a very important case where condition (13.52) holds, in which case all statistics obtained by treating $\ell_{i}(\boldsymbol{\theta})$ as a standard log likelihood are valid. For any definition of $\mathbf{x}_{t}$, we say that $\left\{f_{t}\left(y_{t} \mid \mathbf{x}_{t} ; \boldsymbol{\theta}_{\mathrm{o}}\right): t=1, \ldots, T\right\}$ is a dynamically complete conditional density if


\begin{equation*}
f_{t}\left(y_{t} \mid \mathbf{x}_{t} ; \boldsymbol{\theta}_{\mathrm{o}}\right)=p_{t}^{\mathrm{o}}\left(y_{t} \mid \mathbf{x}_{t}, y_{t-1}, \mathbf{x}_{t-1}, y_{t-2}, \ldots, y_{1}, \mathbf{x}_{1}\right), \quad t=1, \ldots, T . \tag{13.54}
\end{equation*}


In other words, $f_{t}\left(y_{t} \mid \mathbf{x}_{t} ; \boldsymbol{\theta}_{\mathrm{o}}\right)$ must be the conditional density of $y_{t}$ given $\mathbf{x}_{t}$ and the entire past of $\left(\mathbf{x}_{t}, y_{t}\right)$. Equation (13.54) implies that the density is correctly specified. We can state the assumption that $\mathbf{x}_{i t}$ captures all of the distributional dynamics for $y_{i t}$ by writing $\mathrm{D}\left(y_{i t} \mid \mathbf{x}_{i t}, y_{i, t-1}, \mathbf{x}_{i, t-1}, \ldots, y_{i 1}, \mathbf{x}_{i 1}\right)=\mathrm{D}\left(y_{i t} \mid \mathbf{x}_{i t}\right)$, a shorthand that is useful because it is separate from model specification.

When $\mathbf{x}_{t}=\mathbf{z}_{t}$ for contemporaneous exogenous variables, assumption (13.54) is very strong: it means that, once $\mathbf{z}_{t}$ is controlled for, no past values of $\mathbf{z}_{t}$ or $y_{t}$ appear in the conditional density $p_{t}^{\mathrm{o}}\left(y_{t} \mid \mathbf{z}_{t}, y_{t-1}, \mathbf{z}_{t-1}, y_{t-2}, \ldots, y_{1}, \mathbf{z}_{1}\right)$. When $\mathbf{x}_{t}$ contains $\mathbf{z}_{t}$ and some lags-similar to a finite distributed lag model-then equation (13.54) is perhaps more reasonable, but it still assumes that lagged $y_{t}$ has no effect on $y_{t}$ once current and lagged $\mathbf{z}_{t}$ are controlled for. That assumption (13.54) can be false is analogous to the omnipresence of serial correlation in static and finite distributed lag regression models. One important feature of dynamic completeness is that it does not require strict exogeneity of $\mathbf{z}_{t}$ [since only current and lagged $\mathbf{x}_{t}$ appear in equation (13.54)].

Dynamic completeness is more likely to hold when $\mathbf{x}_{t}$ contains lagged dependent variables. The issue, then, is whether enough lags of $y_{t}\left(\right.$ and $\left.\mathbf{z}_{t}\right)$ have been included in $\mathbf{x}_{t}$ to fully capture the dynamics. For example, if $\mathbf{x}_{t} \equiv\left(\mathbf{z}_{t}, y_{t-1}\right)$, then equation (13.54) means that, along with $\mathbf{z}_{t}$, only one lag of $y_{t}$ is needed to capture all of the dynamics.

Showing that condition (13.52) holds under dynamic completeness is easy. First, for each $t, \mathrm{E}\left[\mathbf{s}_{i t}\left(\boldsymbol{\theta}_{\mathrm{o}}\right) \mid \mathbf{x}_{i t}\right]=\mathbf{0}$, since $f_{t}\left(y_{t} \mid \mathbf{x}_{t} ; \boldsymbol{\theta}_{\mathrm{o}}\right)$ is a correctly specified conditional density. But then, under assumption (13.54),\\
$\mathrm{E}\left[\mathbf{s}_{i t}\left(\boldsymbol{\theta}_{\mathrm{o}}\right) \mid \mathbf{x}_{i t}, y_{i, t-1}, \ldots, y_{i 1}, \mathbf{x}_{i 1}\right]=\mathbf{0}$.

Now consider the expected value in condition (13.52) for $r<t$. Since $\mathbf{s}_{i r}\left(\boldsymbol{\theta}_{\mathrm{o}}\right)$ is a function of $\left(\mathbf{x}_{i r}, y_{i r}\right)$, which is in the conditioning set (13.55), the usual iterated expectations argument shows that condition (13.52) holds. It follows that, under dynamic completeness, the usual maximum likelihood statistics from the pooled estimation are asymptotically valid. This result is completely analogous to pooled OLS under dynamic completeness of the conditional mean and homoskedasticity (see Section 7.8).

If the panel data probit model is dynamically complete, any software package that does standard probit can be used to obtain valid standard errors and test statistics, provided the response probability satisfies $\mathrm{P}\left(y_{i t}=1 \mid \mathbf{x}_{i t}\right)=\mathrm{P}\left(y_{i t}=1 \mid \mathbf{x}_{i t}, y_{i, t-1}\right.$, $\left.\mathbf{x}_{i, t-1}, \ldots\right)$. Without dynamic completeness the standard errors and test statistics generally need to be adjusted for serial dependence.

Since dynamic completeness affords nontrivial simplifications, does this fact mean that we should always include lagged values of exogenous and dependent variables until equation (13.54) appears to be satisfied? Not necessarily. Static models are sometimes desirable even if they neglect dynamics. For example, suppose that we have panel data on individuals in an occupation where pay is determined partly by cumulative productivity. (Professional athletes and college professors are two examples.) An equation relating salary to the productivity measures, and possibly demographic variables, is appropriate. Nothing implies that the equation would be dynamically complete; in fact, past salary could help predict current salary, even after controlling for observed productivity. But it does not make much sense to include past salary in the regression equation. As we know from Chapter 10, a reasonable approach is to include an unobserved effect in the equation, and this does not lead to a model with complete dynamics. See also Section 13.9.

We may wish to test the null hypothesis that the density is dynamically complete. White (1994) shows how to test whether the score is serially correlated in a pure time series setting. A similar approach can be used with panel data. A general test for dynamic misspecification can be based on the limiting distribution of (the vectorization of)

$$
N^{-1 / 2} \sum_{i=1}^{N} \sum_{t=2}^{T} \hat{\mathbf{s}}_{i t} \hat{\mathbf{s}}_{i, t-1}^{\prime},
$$

where the scores are evaluated at the pooled MLE. Rather than derive a general statistic here, we will study tests of dynamic completeness in particular applications later (see particularly Chapters 15, 16, and 17).

\subsection*{13.9 Panel Data Models with Unobserved Effects}
As we saw in Chapters 10 and 11, linear unobserved effects panel data models play an important role in modern empirical research. Nonlinear unobserved effects panel data models are becoming increasingly more important. Although we will cover particular models in Chapters 15 through 18, it is useful to have a general treatment.

\subsection*{13.9.1 Models with Strictly Exogenous Explanatory Variables}
For each $i$, let $\left\{\left(\mathbf{y}_{i t}, \mathbf{x}_{i t}\right): t=1,2, \ldots, T\right\}$ be a random draw from the cross section, where $\mathbf{y}_{i t}$ and $\mathbf{x}_{i t}$ can both be vectors. Associated with each cross section unit $i$ is unobserved heterogeneity, $\mathbf{c}_{i}$, which could be a vector. We assume interest lies in the distribution of $\mathbf{y}_{i t}$ given ( $\mathbf{x}_{i t}, \mathbf{c}_{i}$ ). The vector $\mathbf{x}_{i t}$ can contain lags of contemporaneous\\
variables, say $\mathbf{z}_{i t}$ (for example, $\mathbf{x}_{i t}=\left(\mathbf{z}_{i t}, \mathbf{z}_{i, t-1}, \mathbf{z}_{i, t-2}\right)$ ), or even leads of $\mathbf{z}_{i t}$ (for example, $\left.\mathbf{x}_{i t}=\left(\mathbf{z}_{i t}, \mathbf{z}_{i, t+1}\right)\right)$, but not lags of $y_{i t}$. Whatever the lag structure, we let $t=1$ denote the first time period available for estimation.

Let $f_{t}\left(\mathbf{y}_{t} \mid \mathbf{x}_{t}, \mathbf{c} ; \boldsymbol{\theta}\right)$ denote a correctly specified density for each $t$. A key assumption on $\mathbf{x}_{i t}$ is strict exogeneity conditional on the unobserved effects:\\
$\mathrm{D}\left(\mathbf{y}_{i t} \mid \mathbf{x}_{i 1}, \mathbf{x}_{i 2}, \ldots, \mathbf{x}_{i T}, \mathbf{c}_{i}\right)=\mathrm{D}\left(\mathbf{y}_{i t} \mid \mathbf{x}_{i t}, \mathbf{c}_{i}\right), \quad t=1, \ldots, T$,\\
which means that $\mathbf{x}_{i r}, r \neq t$, does not appear in the conditional distribution of $\mathbf{y}_{i t}$ once $\mathbf{x}_{i t}$ and $\mathbf{c}_{i}$ have been counted for. In addition to ruling out lagged dependent variables, (13.56) does not allow for general feedback from unanticipated changes in $\mathbf{y}_{i t}$ to changes in $\mathbf{x}_{i, t+h}$ for $h>1$.

A common but restrictive approach to estimating $\boldsymbol{\theta}_{\mathrm{o}}$ (and other quantities of interest, such as average partial effects) is to assume that $\mathbf{c}_{i}$ is independent of $\mathbf{x}_{i}= \left(\mathbf{x}_{i 1}, \mathbf{x}_{i 2}, \ldots, \mathbf{x}_{i T}\right)$, that is, $\mathrm{D}\left(\mathbf{c}_{i} \mid \mathbf{x}_{i}\right)=\mathrm{D}\left(\mathbf{c}_{i}\right)$, and to model the distribution of $\mathbf{c}_{i}$. Such an approach is very similar to the random effects approach to linear panel data models in Chapter 10. There we did not assume full independence (because linear models do not require such a strong assumption), but we assumed conditional mean independence, $\mathrm{E}\left(\mathbf{c}_{i} \mid \mathbf{x}_{i}\right)=\mathrm{E}\left(\mathbf{c}_{i}\right)$, or, at the very least, zero correlation. In general nonlinear models, it is handy to label $\mathrm{D}\left(\mathbf{c}_{i} \mid \mathbf{x}_{i}\right)=\mathrm{D}\left(\mathbf{c}_{i}\right)$ as a random effects (RE) assumption.

In special cases, which we will address in Chapters 15 and 18, we can consistently estimate $\boldsymbol{\theta}_{\mathrm{o}}$ without imposing any assumptions about $\mathrm{D}\left(\mathbf{c}_{i} \mid \mathbf{x}_{i}\right)$. Such situations are reminiscent of the fixed effects (FE) assumptions from Chapter 10, and we will use that label for nonlinear models, too. As we will see through later examples, in nonlinear models estimation of $\boldsymbol{\theta}_{\mathrm{o}}$ is not always sufficient to calculate the quantities of interest, but it is nevertheless useful when possible.

One way to avoid specifying a model for $\mathrm{D}\left(\mathbf{c}_{i} \mid \mathbf{x}_{i}\right)$ would be to treat $\left\{\mathbf{c}_{i}: i=1, \ldots, N\right\}$ as parameters to estimate along with $\boldsymbol{\theta}_{\mathrm{o}}$. However, because the number of $\mathbf{c}_{i}$ increases with $N$, attempting to estimate them leads to an incidental parameters problem for estimating $\boldsymbol{\theta}_{\mathrm{O}}$. Namely, we cannot consistently estimate $\boldsymbol{\theta}_{\mathrm{O}}$ with fixed $T$ and $N \rightarrow \infty$ (except by fluke in a couple of special cases, including the linear model with additive $c_{i}$ in Chapter 10). For this reason, in this book we reserve the designation "fixed effects" for situations in which we can eliminate the $\mathbf{c}_{i}$ from a particular conditional distribution that still depends on $\boldsymbol{\theta}_{\mathrm{o}}$, and then apply conditional maximum likelihood methods to consistently estimate $\boldsymbol{\theta}_{\mathrm{o}}$ for fixed $T$. We do not give a general treatment because the method applies only in special cases, and we cover those in Part IV.

A middle ground between RE and FE approaches is to allow $\mathrm{D}\left(\mathbf{c}_{i} \mid \mathbf{x}_{i}\right)$ to depend on $\mathbf{x}_{i}$, but then to model this distribution by specifying a parametric model of the conditional density. In Chapter 10, we mentioned that modeling $\mathrm{E}\left(\mathbf{c}_{i} \mid \mathbf{x}_{i}\right)$ as a function of $\mathbf{x}_{i}$ has been called a correlated random effects (CRE) approach, and we will use this label subsequently to refer to situations where we model the distribution $\mathrm{D}\left(\mathbf{c}_{i} \mid \mathbf{x}_{i}\right)$. Often, to impose parsimony we restrict the way in which $\mathrm{D}\left(\mathbf{c}_{i} \mid \mathbf{x}_{i}\right)$ can depend on the time series $\left\{\mathbf{x}_{i t}: t=1,2, \ldots, T\right\}$. A common restriction is $\mathrm{D}\left(\mathbf{c}_{i} \mid \mathbf{x}_{i}\right)=\mathrm{D}\left(\mathbf{c}_{i} \mid \overline{\mathbf{x}}_{i}\right)$, where $\overline{\mathbf{x}}_{i}$ is the time average. With large enough $T$, we can allow $\mathrm{D}\left(\mathbf{c}_{i} \mid \mathbf{x}_{i}\right)$ to depend on other features, such as the individual-specific variance $(T-1)^{-1} \sum_{t=1}^{T}\left(\mathbf{x}_{i t}-\overline{\mathbf{x}}_{i}\right)^{\prime}\left(\mathbf{x}_{i t}-\overline{\mathbf{x}}_{i}\right)$ or average growth rates. For the present treatment, there is no gain in considering special cases. Therefore, let $h(\mathbf{c} \mid \mathbf{x} ; \boldsymbol{\delta})$ be a correctly specified density for $\mathrm{D}\left(\mathbf{c}_{i} \mid \mathbf{x}_{i}\right)$.

Once we have specified $h(\mathbf{c} \mid \mathbf{x} ; \boldsymbol{\delta})$, there are two common ways to proceed. First, we can make the additional assumption that, conditional on $\left(\mathbf{x}_{i}, \mathbf{c}_{i}\right)$, the $\mathbf{y}_{i t}$ are independent. Then, the joint density of $\left(\mathbf{y}_{i 1}, \ldots, \mathbf{y}_{i T}\right)$, given $\left(\mathbf{x}_{i}, \mathbf{c}_{i}\right)$, is\\
$\prod_{t=1}^{T} f_{t}\left(\mathbf{y}_{t} \mid \mathbf{x}_{i t}, \mathbf{c}_{i} ; \boldsymbol{\theta}\right)$.\\
We cannot use this density directly to estimate $\boldsymbol{\theta}_{\mathrm{o}}$ because we do not observe the outcomes $\mathbf{c}_{i}$. Treating the $\mathbf{c}_{i}$ as parameters to estimate with $\boldsymbol{\theta}_{\mathrm{o}}$ leads to the incidental parameters problem mentioned earlier. Instead, we can use the density of $\mathbf{c}_{i}$ given $\mathbf{x}_{i}$ to integrate out the dependence on $\mathbf{c}$. The density of $\mathbf{y}_{i}$ given $\mathbf{x}_{i}$ is\\
$\int_{\mathbb{R}^{J}}\left[\prod_{t=1}^{T} f_{t}\left(\mathbf{y}_{t} \mid \mathbf{x}_{i t}, \mathbf{c} ; \boldsymbol{\theta}_{\mathrm{o}}\right)\right] h\left(\mathbf{c} \mid \mathbf{x}_{i} ; \boldsymbol{\delta}_{\mathrm{o}}\right) d \mathbf{c}$,\\
where $J$ is the dimension of $\mathbf{c}$ and $h(\mathbf{c} \mid \mathbf{x} ; \boldsymbol{\delta})$ is the correctly specified model for the density of $\mathbf{c}_{i}$ given $\mathbf{x}_{i}=\mathbf{x}$. For concreteness, we assume that $\mathbf{c}$ is a continuous random vector. For each $i$, the log-likelihood function is\\
$\log \left\{\int_{\mathbb{R}^{J}}\left[\prod_{t=1}^{T} f_{t}\left(\mathbf{y}_{i t} \mid \mathbf{x}_{i t}, \mathbf{c} ; \boldsymbol{\theta}_{\mathrm{o}}\right)\right] h\left(\mathbf{c} \mid \mathbf{x}_{i} ; \boldsymbol{\delta}_{\mathrm{o}}\right) d \mathbf{c}\right\}$.\\[0pt]
[It is important to see that expression (13.58) does not depend on the $\mathbf{c}_{i} ; \mathbf{c}$ has been integrated out.] Assuming identification and standard regularity conditions, we can consistently estimate $\boldsymbol{\theta}_{\mathrm{O}}$ and $\boldsymbol{\delta}_{\mathrm{O}}$ by conditional MLE, where the asymptotics are for fixed $T$ and $N \rightarrow \infty$. The CMLE is $\sqrt{N}$-asymptotically normal.

A different approach is often simpler and places no restrictions on the joint distribution of the $\mathbf{y}_{i t}$ conditional on $\left(\mathbf{x}_{i}, \mathbf{c}_{i}\right)$. For each $t$, we can obtain the density of $\mathbf{y}_{i t}$\\
given $\mathbf{x}_{i}$ by integrating out $\mathbf{c}_{i}$ :\\
$\int_{\mathbb{R}^{J}}\left[f_{t}\left(\mathbf{y}_{t} \mid \mathbf{x}_{i t}, \mathbf{c} ; \boldsymbol{\theta}_{\mathrm{o}}\right)\right] h\left(\mathbf{c} \mid \mathbf{x}_{i} ; \boldsymbol{\delta}_{\mathrm{o}}\right) d \mathbf{c}$.\\
Now the problem becomes one of partial MLE. We estimate $\boldsymbol{\theta}_{\mathrm{o}}$ and $\boldsymbol{\delta}_{\mathrm{o}}$ by maximizing\\
$\sum_{i=1}^{N} \sum_{t=1}^{T} \log \left\{\int_{\mathbb{R}^{J}}\left[f_{t}\left(\mathbf{y}_{i t} \mid \mathbf{x}_{i t}, \mathbf{c} ; \boldsymbol{\theta}\right)\right] h\left(\mathbf{c} \mid \mathbf{x}_{i} ; \boldsymbol{\delta}\right) d \mathbf{c}\right\}$,\\
where the term in braces is a correctly specified model of the density of $\mathrm{D}\left(\mathbf{y}_{i t} \mid \mathbf{x}_{i}\right)$. (Actually, using PMLE, $\boldsymbol{\theta}_{\mathrm{O}}$ and $\boldsymbol{\delta}_{\mathrm{O}}$ are not always separately identified, although interesting functions of them, such as average partial effects, are. We will see examples in Chapters 15, 16, and 17.) Across time, the scores for each $i$ will necessarily be serially correlated because the $\mathbf{y}_{i t}$ are dependent when we condition only on $\mathbf{x}_{i}$, and not also on $\mathbf{c}_{i}$. Therefore, we must make inference robust to serial dependence, as in Section 13.8.2. In Chapter 15, we will study both the conditional MLE and partial MLE approaches for unobserved effects probit models. We do the same for Tobit models in Chapter 17.

\subsection*{13.9.2 Models with Lagged Dependent Variables}
Now assume that we are interested in modeling $\mathrm{D}\left(\mathbf{y}_{i t} \mid \mathbf{z}_{i t}, \mathbf{y}_{i, t-1}, \mathbf{c}_{i}\right)$ where, for simplicity, we include only contemporaneous conditioning variables, $\mathbf{z}_{i t}$, and only one lag of $\mathbf{y}_{i t}$. Adding lags (or even leads) of $\mathbf{z}_{i t}$ or more lags of $\mathbf{y}_{i t}$ requires only a notational change.

A key assumption is that we have the dynamics correctly specified and that $\mathbf{z}_{i}= \left\{\mathbf{z}_{i 1}, \ldots, \mathbf{z}_{i T}\right\}$ is appropriately strictly exogenous (conditional on $\mathbf{c}_{i}$ ). These assumptions are both captured by\\
$\mathrm{D}\left(\mathbf{y}_{i t} \mid \mathbf{z}_{i t}, \mathbf{y}_{i, t-1}, \mathbf{c}_{i}\right)=\mathrm{D}\left(\mathbf{y}_{i t} \mid \mathbf{z}_{i}, \mathbf{y}_{i, t-1}, \ldots, \mathbf{y}_{i 0}, \mathbf{c}_{i}\right)$.

We assume that $f_{t}\left(\mathbf{y}_{t} \mid \mathbf{z}_{t}, \mathbf{y}_{t-1}, \mathbf{c} ; \boldsymbol{\theta}\right)$ is a correctly specified density for the conditional distribution on the left-hand side of equation (13.60). Given strict exogeneity of $\left\{\mathbf{z}_{i t}: t=1, \ldots, T\right\}$ and dynamic completeness, the density of $\left(\mathbf{y}_{i 1}, \ldots, \mathbf{y}_{i T}\right)$ given $\left(\mathbf{z}_{i}=\mathbf{z}, \mathbf{y}_{i 0}=\mathbf{y}_{0}, \mathbf{c}_{i}=\mathbf{c}\right)$ is\\
$\prod_{t=1}^{T} f_{t}\left(\mathbf{y}_{t} \mid \mathbf{z}_{t}, \mathbf{y}_{t-1}, \mathbf{c} ; \boldsymbol{\theta}_{\mathrm{o}}\right)$.\\
(By convention, $\mathbf{y}_{i 0}$ is the first observation on $\mathbf{y}_{i t}$.) Again, to estimate $\boldsymbol{\theta}_{\mathrm{o}}$, we integrate c out of this density. To do so, we specify a density for $\mathbf{c}_{i}$ given $\mathbf{z}_{i}$ and the initial value\\
$\mathbf{y}_{i 0}$ (sometimes called the initial condition). Let $h\left(\mathbf{c} \mid \mathbf{z}, \mathbf{y}_{0} ; \boldsymbol{\delta}\right)$ denote the model for this conditional density. Then, assuming that we have this model correctly specified, the density of $\left(\mathbf{y}_{i 1}, \ldots, \mathbf{y}_{i T}\right)$ given $\left(\mathbf{z}_{i}=\mathbf{z}, \mathbf{y}_{i 0}=\mathbf{y}_{0}\right)$ is


\begin{equation*}
\int_{\mathbb{R}^{J}}\left[\prod_{t=1}^{T} f_{t}\left(\mathbf{y}_{t} \mid \mathbf{z}_{t}, \mathbf{y}_{t-1}, \mathbf{c} ; \boldsymbol{\theta}_{\mathrm{o}}\right)\right] h\left(\mathbf{c} \mid \mathbf{z}, \mathbf{y}_{0} ; \boldsymbol{\delta}_{\mathrm{o}}\right) d \mathbf{c} \tag{13.62}
\end{equation*}


which, for each $i$, leads to the log-likelihood function conditional on $\left(\mathbf{z}_{i}, \mathbf{y}_{i 0}\right)$ :\\
$\log \left\{\int_{\mathbb{R}^{J}}\left[\prod_{t=1}^{T} f_{t}\left(\mathbf{y}_{i t} \mid \mathbf{z}_{i t}, \mathbf{y}_{i, t-1}, \mathbf{c} ; \boldsymbol{\theta}\right)\right] h\left(\mathbf{c} \mid \mathbf{z}_{i}, \mathbf{y}_{i 0} ; \boldsymbol{\delta}\right) d \mathbf{c}\right\}$.

We sum expression (13.63) across $i=1, \ldots, N$ and maximize with respect to $\boldsymbol{\theta}$ and $\boldsymbol{\delta}$ to obtain the CMLEs. Provided all functions are sufficiently differentiable and identification holds, the CMLEs are consistent and $\sqrt{N}$-asymptotically normal, as usual. Because we have fully specified the conditional density of $\left(\mathbf{y}_{i 1}, \ldots, \mathbf{y}_{i T}\right)$ given $\left(\mathbf{z}_{i}, \mathbf{y}_{i 0}\right)$, the general theory of CMLE applies directly. (The fact that the distribution of $\mathbf{y}_{i 0}$ given $\mathbf{z}_{i}$ would typically depend on $\boldsymbol{\theta}_{\mathrm{o}}$ has no bearing on the consistency of the CMLE. The fact that we are conditioning on $\mathbf{y}_{i 0}$, rather than basing the analysis on $\mathrm{D}\left(\mathbf{y}_{i 0}, \mathbf{y}_{i 1}, \ldots, \mathbf{y}_{i T} \mid \mathbf{z}_{i}\right)$, means that we are generally sacrificing efficiency. But by conditioning on $\mathbf{y}_{i 0}$ we do not have to find $\mathrm{D}\left(\mathbf{y}_{i 0} \mid \mathbf{z}_{i}\right)$, a feat that can be very difficult, if not impossible.) The asymptotic variance of $\left(\hat{\boldsymbol{\theta}}^{\prime}, \hat{\boldsymbol{\delta}}^{\prime}\right)^{\prime}$ can be estimated by any of the formulas in equation (13.32) (properly modified to account for estimation of $\boldsymbol{\theta}_{\mathrm{o}}$ and $\boldsymbol{\delta}_{\mathrm{O}}$ ).

A weakness of the CMLE approach is that we must specify a density for $\mathbf{c}_{i}$ given $\left(\mathbf{z}_{i}, \mathbf{y}_{i 0}\right)$, but this is a price we pay for estimating dynamic, nonlinear models with unobserved effects. The alternative of treating the $\mathbf{c}_{i}$ as parameters to estimatewhich is, unfortunately, often labeled the "fixed effects" approach-does not lead to consistent estimation of $\boldsymbol{\theta}_{\mathrm{O}}$ because of the incidental parameters problem.

In any application, several issues need to be addressed. First, when are the parameters identified? Second, what quantities are we interested in? As we cannot observe $\mathbf{c}_{i}$, we typically want to average out $\mathbf{c}_{i}$ when obtaining partial effects. Wooldridge (2005b) shows that average partial effects are generally identified under the assumptions that we have made. Finally, obtaining the CMLE can be very difficult computationally, as can be obtaining the asymptotic variance estimates in equation (13.32). If $\mathbf{c}_{i}$ is a scalar, estimation is easier, but there is still a one-dimensional integral to approximate for each $i$. In Chapters 15 through 18 we will see that, under reasonable assumptions, standard software can be used to estimate dynamic models with unob-\\
served effects, including effects that are averaged across the distribution of heterogeneity. See also Problem 13.11 for application to a dynamic linear model.

\subsection*{13.10 Two-Step Estimators Involving Maximum Likelihood}
From the results in Chapter 12 on two-step estimation, we know that consistency and asymptotic normality generally apply to two-step estimators that involve maximum likelihood. We provide a brief treatment here, covering two cases. The first is when the second-step estimation is MLE (and the first step may or may not be). Then we cover an interesting situation that can arise when the first-step estimator is MLE.

\subsection*{13.10.1 Second-Step Estimator Is Maximum Likelihood Estimator}
Assume that we have a correctly specified density for the conditional distribution $\mathrm{D}\left(\mathbf{y}_{i} \mid \mathbf{x}_{i}\right)$. Write the model as $f(\mathbf{y} \mid \mathbf{x} ; \boldsymbol{\theta}, \boldsymbol{\gamma})$ for $\boldsymbol{\theta}$ a $P \times 1$ vector and $\boldsymbol{\gamma}$ a $J \times 1$ vector. The true density is $f\left(\mathbf{y} \mid \mathbf{x} ; \boldsymbol{\theta}_{\mathrm{o}}, \boldsymbol{\gamma}_{\mathrm{o}}\right)$. A preliminary estimator of $\boldsymbol{\gamma}_{\mathrm{o}}$, say $\hat{\boldsymbol{\gamma}}$, is plugged into the log-likelihood function, and $\hat{\boldsymbol{\theta}}$ solves\\
$\max _{\boldsymbol{\theta} \in \boldsymbol{\Theta}} \sum_{i=1}^{N} \log f\left(\mathbf{y}_{i} \mid \mathbf{x}_{i} ; \boldsymbol{\theta}, \hat{\boldsymbol{\gamma}}\right)$.\\
We call $\hat{\boldsymbol{\theta}}$ a two-step maximum likelihood estimator. Consistency of $\hat{\boldsymbol{\theta}}$ follows from results for two-step M-estimators. The practical limitation is that $\log f\left(\mathbf{y}_{i} \mid \mathbf{x}_{i} ; \boldsymbol{\theta}, \gamma\right)$ is continuous on $\boldsymbol{\Theta} \times \boldsymbol{\Gamma}$ and that $\boldsymbol{\theta}_{\mathrm{o}}$ and $\gamma_{\mathrm{o}}$ are identified.

Asymptotic normality of the two-step MLE follows directly from the results on two-step M-estimation in Chapter 12. As we saw there, in general the asymptotic variance of $\sqrt{N}\left(\hat{\boldsymbol{\theta}}-\boldsymbol{\theta}_{\mathrm{o}}\right)$ depends on the asymptotic variance of $\sqrt{N}\left(\hat{\gamma}-\boldsymbol{\gamma}_{\mathrm{o}}\right)$ [see equation (12.41)], so we need to know the estimation problem solved by $\hat{\gamma}$. In some cases estimation of $\gamma_{\mathrm{o}}$ can be ignored. An important case is where the expected Hessian, defined with respect to $\boldsymbol{\theta}$ and $\boldsymbol{\gamma}$, is block diagonal (the matrix $\mathbf{F}_{\mathrm{O}}$ in equation (12.36) is zero in this case). It can also hold for some values of $\boldsymbol{\theta}_{\mathrm{o}}$, which is important for testing certain hypotheses. We will encounter several examples in Part IV.

Of course, we can also apply two-step methods when the second step is a partial MLE, resulting in a two-step partial MLE. As usual, consistency and asymptotic normality hold under correct specification of the marginal (conditional) density for each $t$ and standard regularity conditions. The practical issue is in computing an appropriate asymptotic variance. We will not give a general treatment here, as the general results on M-estimation can be applied directly when we need them in Part IV.

\subsection*{13.10.2 Surprising Efficiency Result When the First-Step Estimator Is Conditional Maximum Likelihood Estimator}
We now turn to a case where the conclusion seems counterintuitive: under certain assumptions, estimating parameters in a first stage, by CMLE, can improve the asymptotic efficiency of a second-step M-estimator. To study this phenomenon, assume that the second-step M-estimator, $\hat{\boldsymbol{\theta}}$, solves the problem


\begin{equation*}
\min _{\boldsymbol{\theta} \in \boldsymbol{\Theta}} \sum_{i=1}^{N} q\left(\mathbf{v}_{i}, \mathbf{w}_{i}, \mathbf{z}_{i}, \boldsymbol{\theta}, \hat{\boldsymbol{\gamma}}\right), \tag{13.64}
\end{equation*}


where we maintain the regularity conditions used in Chapter 12 to apply mean value expansions, the uniform law of large numbers, and the central limit theorem. For reasons we will see, we have separated the data vector in (13.64) into three vectors, $\mathbf{v}_{i}$, $\mathbf{w}_{i}$, and $\mathbf{z}_{i}$. We assume the first-step estimator, $\hat{\gamma}$, comes from a (conditional) MLE problem (that satisfies the appropriate regularity conditions):

$$
\max _{\gamma \in \Gamma} \sum_{i=1}^{N} \log h\left(\mathbf{v}_{i} \mid \mathbf{z}_{i} ; \gamma\right)
$$

where $h(\cdot \mid \mathbf{z} ; \gamma)$ is a model of the density underlying $\mathrm{D}\left(\mathbf{v}_{i} \mid \mathbf{z}_{i}\right)$, and we assume this density is correctly specified with population value $\gamma_{\mathrm{o}}$. By the information matrix equality, we can assume


\begin{equation*}
\sqrt{N}\left(\hat{\gamma}-\gamma_{\mathrm{o}}\right)=\left\{\mathrm{E}\left[\mathbf{d}_{i}\left(\gamma_{\mathrm{o}}\right) \mathbf{d}_{i}\left(\gamma_{\mathrm{o}}\right)^{\prime}\right]\right\}^{-1} N^{-1 / 2} \sum_{i=1}^{N} \mathbf{d}_{i}\left(\gamma_{\mathrm{o}}\right)+o_{p}(1), \tag{13.65}
\end{equation*}


where $\mathbf{d}_{i}(\gamma) \equiv \nabla_{\gamma} \log h\left(\mathbf{v}_{i} \mid \mathbf{z}_{i} ; \gamma\right)^{\prime}$ is the $J \times 1$ score of the first-step log likelihood. If we assume nothing further, we would simply have to derive, and estimate, the asymptotic variance of $\sqrt{N}\left(\hat{\boldsymbol{\theta}}-\boldsymbol{\theta}_{\mathrm{o}}\right)$ using the methods in Section 12.5.2. But suppose we add the assumption


\begin{equation*}
\mathrm{D}\left(\mathbf{v}_{i} \mid \mathbf{w}_{i}, \mathbf{z}_{i}\right)=\mathrm{D}\left(\mathbf{v}_{i} \mid \mathbf{z}_{i}\right), \tag{13.66}
\end{equation*}


which is often called a conditional independence assumption: conditional on $\mathbf{z}_{i}, \mathbf{v}_{i}$ and $\mathbf{w}_{i}$ are independent. In a sense, assumption (13.66) means that $\mathbf{z}_{i}$ is such a good explainor of $\mathbf{v}_{i}$-at least relative to $\mathbf{w}_{i}$-that once we know $\mathbf{z}_{i}, \mathbf{w}_{i}$ tells us nothing about the likelihood of outcomes on $\mathbf{v}_{i}$. This assumption is special, but, as we will see in Part IV, it holds under certain stratified sampling schemes as well as in the context of estimating treatment effects under so-called ignorability of treatment assumptions.

Now, let $\mathbf{s}_{i}(\boldsymbol{\theta}, \boldsymbol{\gamma}) \equiv \mathbf{s}\left(\mathbf{v}_{i}, \mathbf{w}_{i}, \mathbf{z}_{i}, \boldsymbol{\theta}, \boldsymbol{\gamma}\right) \equiv \nabla_{\boldsymbol{\theta}} q\left(\mathbf{v}_{i}, \mathbf{w}_{i}, \mathbf{z}_{i}, \boldsymbol{\theta}, \boldsymbol{\gamma}\right)^{\prime}$ be the $P \times 1$ score of the second-step objective function, but only with respect to $\boldsymbol{\theta}$. In order to find Avar $\sqrt{N}\left(\hat{\boldsymbol{\theta}}-\boldsymbol{\theta}_{\mathrm{o}}\right)$, we need to find the $P \times J$ matrix $\mathbf{F}_{\mathrm{o}}=\mathrm{E}\left[\nabla_{\gamma} \mathbf{s}_{i}\left(\boldsymbol{\theta}_{\mathrm{o}}, \boldsymbol{\gamma}_{\mathrm{o}}\right)\right]$. The key step is to use the generalized CIME in equation (13.39), where it is important to note that, under (13.66), $\mathbf{d}_{i}(\boldsymbol{\gamma})$ is the score for the density of $\mathbf{v}_{i}$ given ( $\mathbf{w}_{i}, \mathbf{z}_{i}$ ), even though it depends only on $\mathbf{z}_{i}$. Therefore, because $\mathbf{s}_{i}\left(\boldsymbol{\theta}_{\mathrm{o}}, \boldsymbol{\gamma}\right)$ is a function of $\left(\mathbf{v}_{i}, \mathbf{w}_{i}, \mathbf{z}_{i}\right)$, we have


\begin{equation*}
-\mathrm{E}\left[\nabla_{\gamma} \mathbf{s}_{i}\left(\boldsymbol{\theta}_{\mathrm{o}}, \gamma_{\mathrm{o}}\right) \mid \mathbf{w}_{i}, \mathbf{z}_{i}\right]=\mathrm{E}\left[\mathbf{s}_{i}\left(\boldsymbol{\theta}_{\mathrm{o}}, \gamma_{\mathrm{o}}\right) \mathbf{d}_{\mathbf{i}}(\gamma)^{\prime} \mid \mathbf{w}_{i}, \mathbf{z}_{i}\right], \tag{13.67}
\end{equation*}


so, using iterated expectations, we conclude that $\mathbf{F}_{\mathrm{o}}=-\mathrm{E}\left[\mathbf{s}_{i}\left(\boldsymbol{\theta}_{\mathrm{o}}, \boldsymbol{\gamma}_{\mathrm{o}}\right) \mathbf{d}_{\mathrm{i}}(\boldsymbol{\gamma})^{\prime}\right]$. But then by equation (12.39) and (13.65),

$$
\begin{aligned}
\sqrt{N}\left(\hat{\boldsymbol{\theta}}-\boldsymbol{\theta}_{\mathrm{o}}\right) & =-\mathbf{A}_{\mathrm{o}}^{-1} N^{-1 / 2} \sum_{i=1}^{N}\left\{\mathbf{s}_{i}^{\mathrm{o}}-\mathrm{E}\left(\mathbf{s}_{i}^{\mathrm{o}} \mathbf{d}_{i}^{\mathrm{o} \prime}\right)\left[\mathrm{E}\left(\mathbf{d}_{i}^{\mathrm{o}} \mathbf{d}_{i}^{\mathrm{o} \prime}\right)\right]^{-1} \mathbf{d}_{i}^{\mathrm{o}}\right\}+o_{p}(1) \\
& \equiv-\mathbf{A}_{\mathrm{o}}^{-1} N^{-1 / 2} \sum_{i=1}^{N} \mathbf{g}_{i}^{\mathrm{o}}+o_{p}(1)
\end{aligned}
$$

where $\mathbf{A}_{\mathrm{o}} \equiv \mathrm{E}\left[\nabla_{\theta} \mathbf{s}_{i}\left(\boldsymbol{\theta}_{\mathrm{o}}, \boldsymbol{\gamma}_{\mathrm{o}}\right)\right]$ is the $P \times P$ Hessian of the objective function with respect to $\boldsymbol{\theta}, \mathbf{g}_{i}^{\mathrm{o}} \equiv \mathbf{s}_{i}^{\mathrm{o}}-\mathrm{E}\left(\mathbf{s}_{i}^{\mathrm{o}} \mathbf{d}_{i}^{\mathrm{o} \prime}\right)\left[\mathrm{E}\left(\mathbf{d}_{i}^{\mathrm{o}} \mathbf{d}_{i}^{\mathrm{o} \prime}\right)\right]^{-1} \mathbf{d}_{i}^{\mathrm{o}}$ are the population residuals from the population system regression of $\mathbf{s}_{i}^{\mathrm{o}}$ on $\mathbf{d}_{i}^{\mathrm{o}}$, and the " $o$ " superscript denotes evaluation at $\boldsymbol{\theta}_{\mathrm{o}}$ and $\gamma_{\mathrm{o}}$ or just $\gamma_{\mathrm{o}}$. Therefore,\\
Avar $\sqrt{N}\left(\hat{\boldsymbol{\theta}}-\boldsymbol{\theta}_{\mathrm{o}}\right)=\mathbf{A}_{\mathrm{o}}^{-1} \mathbf{D}_{\mathrm{o}} \mathbf{A}_{\mathrm{o}}^{-1}$,\\
where\\
$\mathbf{D}_{\mathrm{o}}=\mathrm{E}\left(\mathbf{g}_{i}^{\mathrm{o}} \mathbf{g}_{i}^{\mathrm{o} \prime}\right)=\operatorname{Var}\left(\mathbf{g}_{i}^{\mathrm{o}}\right)$.

If we knew $\gamma_{\mathrm{o}}$ rather than estimating it by CMLE, the asymptotic variance of the estimator, say $\tilde{\boldsymbol{\theta}}$, would be $\operatorname{Avar} \sqrt{N}\left(\tilde{\boldsymbol{\theta}}-\boldsymbol{\theta}_{\mathrm{o}}\right)=\mathbf{A}_{\mathrm{o}}^{-1} \mathbf{B}_{\mathrm{o}} \mathbf{A}_{\mathrm{o}}^{-1}$ where $\mathbf{B}_{\mathrm{o}}=\mathrm{E}\left(\mathbf{s}_{i}^{\mathrm{o}} \mathbf{s}_{i}^{\mathrm{o} /}\right)$ the usual expected outer product of the score without accounting for the first-step estimation (because there is none). But $\mathbf{B}_{\mathrm{o}}-\mathbf{D}_{\mathrm{o}}$ is positive semi-definite (p.s.d.), and so the two-step M-estimator is generally more (asymptotically) efficient than the onestep M-estimator that uses knowledge of $\boldsymbol{\gamma}_{\mathrm{o}}$. In one case, the two estimators are asymptotically equivalent, namely, when $\mathrm{E}\left(\mathbf{s}_{i}^{\mathrm{o}} \mathbf{d}_{i}^{\mathrm{o} /}\right)=\mathbf{0}$ (which implies $\mathbf{F}_{\mathrm{o}}=\mathbf{0}$ ).

An immediate implication of the improvement in efficiency in estimating $\gamma_{\mathrm{o}}$ is that if we do use $\hat{\gamma}$ but then ignore the estimation in the second stage, our inference will be conservative. In particular, the standard errors computed from $\hat{\mathbf{A}}^{-1} \hat{\mathbf{B}} \hat{\mathbf{A}}^{-1} / N$ are larger, or no smaller, than they could be. (As usual, the "\^{}"" means that we replace all parameters with their estimates.) Consequently, using a standard econometrics\\
package that computes sandwich standard errors but does not easily account for firststep estimation produces conservative inference under assumption (13.66).

It is also easy to compute standard errors that reflect the increased precision from using the MLE, $\hat{\boldsymbol{\gamma}}$, in place of $\boldsymbol{\gamma}_{\mathrm{o}}$. Namely, let $\hat{\mathbf{s}}_{i}=\mathbf{s}_{i}(\hat{\boldsymbol{\theta}}, \hat{\boldsymbol{\gamma}})$ and $\hat{\mathbf{d}}_{i}=\mathbf{d}_{i}(\hat{\boldsymbol{\gamma}})$ be the scores from the second- and first-step estimations, respectively. (Remember, $\mathbf{s}_{i}(\boldsymbol{\theta}, \boldsymbol{\gamma})$ is the score with respect to $\boldsymbol{\theta}$ only.) Then, let $\hat{\mathbf{g}}_{i}^{\prime}$ be the $1 \times P$ residuals from the multivariate regression of $\hat{\mathbf{s}}_{i}^{\prime}$ on $\hat{\mathbf{d}}_{i}^{\prime}, i=1, \ldots, N$. Then, we obtain


\begin{equation*}
\hat{\mathbf{D}}=N^{-1} \sum_{i=1}^{N} \hat{\mathbf{g}}_{i} \hat{\mathbf{g}}_{i}^{\prime} \tag{13.70}
\end{equation*}


and form the sandwich $\hat{\mathbf{A}}^{-1} \hat{\mathbf{D}} \hat{\mathbf{A}}^{-1} / N$ as $\widehat{\operatorname{Avar}(\hat{\boldsymbol{\theta}})}$.

\subsection*{13.11 Quasi-Maximum Likelihood Estimation}
Until now, we have assumed that a (conditional) density function has been correctly specified, although in the panel data case we covered situations where we might only specify a density separately for each time period (see Sections 13.8 and 13.9). As we have seen, if we have a correctly specified density for the conditional distributions $\mathrm{D}\left(\mathbf{y}_{i} \mid \mathbf{x}_{i}\right)$ or $\mathrm{D}\left(\mathbf{y}_{i t} \mid \mathbf{x}_{i t}\right), t=1, \ldots, T$, then MLE or partial MLE has desirable largesample properties for estimating the population parameters.

There are situations in which it is important to understand the properties of MLE methods when densities are misspecified, or partly misspecified. Some authors, notably White (1982a, 1994), take the view that all models should be viewed, at least initially, as being misspecified, and inference and interpretation should proceed accordingly. White (1994) also considers an intermediate view in which part of the distribution is correctly specified-the leading case being a conditional mean function-while other aspects might be misspecified. As shown by Gourieroux, Monfort, and Trognon (1984a), this posture is especially useful for the class of densities in the linear exponential family, where often we either suspect or know for certain that the full distribution is not correctly specified, but we pay careful attention to specification of the conditional mean. In the next few subsections, we study the properties of estimators that maximize a log likelihood under various degrees of misspecification.

\subsection*{13.11.1 General Misspecification}
In general analyses of maximum likelihood estimation of misspecified models, one typically posits the same setup as in Section 13.3, less the assumption that the con-\\
ditional density is correctly specified, stated in equation (13.11). Thus, for MLE with a generally misspecified density, there is no "true" value of theta, which we called $\boldsymbol{\theta}_{\mathrm{o}}$. Instead, it is standard to postulate the existence of a unique solution to the population problem (13.14). Following White (1994), we denote this value by $\boldsymbol{\theta}^{*}$. White $(1982 a, 1994)$ also discusses the interpretation of $\boldsymbol{\theta}^{*}$ as providing the best approximation to the true density in the parametric class $f(\mathbf{y} \mid \mathbf{x} ; \boldsymbol{\theta})$, where closeness is measured in terms of the Kullback-Leibler information criterion. This interpretation can be inferred from our discussion in Section 13.3.

When the density is misspecified, we typically call the solution to (13.15), which we still denote by $\hat{\boldsymbol{\theta}}$, a quasi-maximum likelihood estimator (QMLE). Some authors prefer the name pseudo-maximum likelihood estimator. We also refer to the loglikelihood function as the quasi-log-likelihood (pseudo-log-likelihood) function.

Consistency of $\hat{\boldsymbol{\theta}}$ for $\boldsymbol{\theta}^{*}$ follows in the same way as when the model is correctly specified: by assumption, $\boldsymbol{\theta}^{*}$ maximizes $\mathrm{E}\left[\log f\left(y_{i} \mid \mathbf{x}_{i} ; \boldsymbol{\theta}\right)\right]$ over the parameter space $\Theta$, and then we simply assume enough regularity conditions so that the quasilog likelihood converges to its expectation uniformly over $\boldsymbol{\Theta}$, just as in Theorem 13.1.

Asymptotic inference concerning $\boldsymbol{\theta}^{*}$ is more interesting. First, one might legitimately ask: If the model is misspecified, what does it mean to test hypotheses about $\boldsymbol{\theta}^{*}$ ? After all, $\boldsymbol{\theta}^{*}$ does not generally index conditional probabilities or conditional moments. Nevertheless, if we take the realistic stance that models of conditional densities are probably misspecified, the best we can do is to test hypotheses about our best approximation to the true density. And sometimes we assume that the main model we are interested in is correctly specified, but we estimate an auxiliary model as a way to obtain, say, instrumental variables. In such cases, we often do not want to assume that the auxiliary model-which is often chosen for computational convenience-is correctly specified in any sense.

It is fairly straightforward to conduct inference on $\boldsymbol{\theta}^{*}$. Without further assumptions, there is only one legitimate estimator of $\operatorname{Avar}(\hat{\boldsymbol{\theta}})$ :\\
$\widehat{\operatorname{Avar}(\hat{\boldsymbol{\theta}})}=\left(\sum_{i=1}^{N} \mathbf{H}_{i}(\hat{\boldsymbol{\theta}})\right)^{-1}\left(\sum_{i=1}^{N} \mathbf{s}_{i}(\hat{\boldsymbol{\theta}}) \mathbf{s}_{i}(\hat{\boldsymbol{\theta}})^{\prime}\right)\left(\sum_{i=1}^{N} \mathbf{H}_{i}(\hat{\boldsymbol{\theta}})\right)^{-1}$,\\
where, as before, $\mathbf{s}_{i}(\boldsymbol{\theta})$ is the $P \times 1$ score vector and $\mathbf{H}_{i}(\boldsymbol{\theta})$ is the $P \times P$ Hessian. (As usual, this estimator is "legitimate" in the sense that, when divided by $N$, the righthand side of (13.71) converges in probability to $\operatorname{Avar}\left[\sqrt{N}\left(\hat{\boldsymbol{\theta}}-\boldsymbol{\theta}^{*}\right)\right]=\mathbf{A}^{*-1} \mathbf{B}^{*} \mathbf{A}^{*-1}$, where $\mathbf{A}^{*} \equiv-\mathrm{E}\left[\mathbf{H}_{i}\left(\boldsymbol{\theta}^{*}\right)\right]$ and $\mathbf{B}^{*} \equiv \mathrm{E}\left[\mathbf{s}_{i}\left(\boldsymbol{\theta}^{*}\right) \mathbf{s}_{i}\left(\boldsymbol{\theta}^{*}\right)^{\prime}\right]$; see Section 12.3.) In some caseswe cover an important one in the next subsection-we can use the expected Hessian\\
conditional on $\mathbf{x}_{i}$, but, in general, $\mathrm{E}\left[\mathbf{H}_{i}\left(\boldsymbol{\theta}^{*}\right) \mid \mathbf{x}_{i}\right]$ cannot be computed if $f(\mathbf{y} \mid \mathbf{x} ; \boldsymbol{\theta})$ is misspecified.

The estimate in (13.71) requires computing both first and second derivatives, but these calculations can be done numerically, if necessary. Once we have $\operatorname{Avar}(\hat{\boldsymbol{\theta}})$, forming Wald tests for restrictions on $\boldsymbol{\theta}^{*}$ is straightforward. In particular, individual asymptotic $t$ statistics are easily obtained because the standard errors of the $\hat{\theta}_{j}$ are the square roots of the diagonal elements of (13.71).

Score testing must be carried out using the fully robust framework in Section 12.6.2 because the information matrix equality cannot be assumed to hold. Plus, the conditional expected value of the Hessian cannot be computed in general, so the score test that uses the Hessian (evaluated at the restricted estimates) is the only statistic that can be relied on. Fortunately, the statistic in equation (12.68) is always nonnegative even though $\tilde{\mathbf{A}}$ might not be positive definite (because $\tilde{\mathbf{B}}$ is always at least p.s.d.). Inference based on the LR statistic is very difficult and not advised: the LR statistic no longer has a limiting chi-square distribution (and its limiting distribution depends on unknown parameters).

As an example, consider the probit model we carried along in earlier sections, but where $\mathrm{P}\left(y_{i}=1 \mid \mathbf{x}_{i}\right) \neq \Phi\left(\mathbf{x}_{i} \boldsymbol{\theta}\right)$ for all $K \times 1$ vectors $\boldsymbol{\theta}$. In other words, the probit model is misspecified. If we obtain $\hat{\boldsymbol{\theta}}$ by maximizing the probit log likelihood, under very weak conditions $\hat{\boldsymbol{\theta}}$ converges in probability to some $\boldsymbol{\theta}^{*} \in \mathbb{R}^{K}$, and $\Phi\left(\mathbf{x} \boldsymbol{\theta}^{*}\right)$ provides the "best" approximation to $\mathrm{P}\left(y_{i}=1 \mid \mathbf{x}_{i}=\mathbf{x}\right)$ in the sense of minimizing the Kullback-Leibler distance. Based on the probit model, for a continuous conditioning variable, say $x_{j}$, we would estimate the partial effect of $x_{j}$ on $\mathrm{P}\left(y_{i}=1 \mid \mathbf{x}_{i}=\mathbf{x}\right)$ as the partial derivative $\theta_{j}^{*} \phi\left(\mathbf{x} \boldsymbol{\theta}^{*}\right)$. Therefore, having an estimate of the asymptotic variance of $\hat{\boldsymbol{\theta}}$ is critical for obtaining, say, an asymptotic confidence interval for this approximate partial effect. Some econometrics packages include simple options for computing the sandwich estimator in (13.71), in which case inference under misspecification is straightforward.

We can also allow for complete density misspecification in the context of partial (pooled) MLE. We must allow for a general estimate of the Hessian: for each $i, \mathbf{H}_{i}(\hat{\boldsymbol{\theta}})=\sum_{t=1}^{T} \mathbf{H}_{i t}(\hat{\boldsymbol{\theta}})$. Further, without assuming that $f_{t}\left(\mathbf{y}_{t} \mid \mathbf{x}_{t} ; \boldsymbol{\theta}\right)$ is correctly specified for each $t$, we can no longer conclude that $\mathrm{D}\left(\mathbf{y}_{i t} \mid \mathbf{x}_{i t}\right)= \mathrm{D}\left(\mathbf{y}_{i t} \mid \mathbf{x}_{i t}, \mathbf{y}_{i, t-1}, \mathbf{x}_{i, t-1}, \ldots, \mathbf{y}_{i 1}, \mathbf{x}_{i 1}\right)$ is sufficient for the scores of the log likelihood to be serially uncorrelated (when evaluated now at $\boldsymbol{\theta}^{*}$ ). Therefore, without further analysis, one must use (13.71) as the estimated asymptotic variance, where $\mathbf{s}_{i}(\hat{\boldsymbol{\theta}})= \sum_{t=1}^{T} \mathbf{s}_{i t}(\hat{\boldsymbol{\theta}})$, as in Section 13.8.2. (For commonly used nonlinear models, such as probit models, some econometrics software packages include options to compute the fully robust matrix (13.70) for panel data applications.)

\subsection*{13.11.2 Model Selection Tests}
The properties of maximum likelihood under general misspecification can be used to derive a model selection test due to Vuong (1989). As the name implies, the test is intended to allow one to choose between competing models. Here we treat the case where the two models are, in a sense to be made precise, nonnested. (When one model is a special case of the other, the score approach provides a much simpler way to test an attractive null model against a more general alternative.)

If we are content with just choosing the model with the "best fit" given the data at hand, then it is legitimate to choose the model with the largest value of the log likelihood. After all, whether or not the model is correctly specified, the average log likelihood consistently estimates the negative of the Kullback-Leibler information criterion (KLIC) for that model. And, we know from Section 13.3, the true density maximizes the KLIC. Therefore, a density model cannot be correctly specified if it delivers (asymptotically) a lower average log likelihood than another model. Comparing log-likelihood values is analogous to comparing $R$-squareds in a regression context.

As suggested by Vuong (1989), it is useful to attach statistical significance between the difference in log likelihoods. For nonnested models, this turns out to be remarkably easy. Let $f_{1}\left(\mathbf{y} \mid \mathbf{x} ; \boldsymbol{\theta}_{1}\right)$ and $f_{2}\left(\mathbf{y} \mid \mathbf{x} ; \boldsymbol{\theta}_{2}\right)$ be competing models for the density of $\mathrm{D}\left(\mathbf{y}_{i} \mid \mathbf{x}_{i}\right)$, where both may be misspecified. Let $\hat{\boldsymbol{\theta}}_{1}$ and $\hat{\boldsymbol{\theta}}_{2}$ be the QMLEs converging to $\boldsymbol{\theta}_{1}^{*}$ and $\boldsymbol{\theta}_{2}^{*}$, respectively. Let $\mathscr{L}_{m}=\sum_{i=1}^{N} \ell_{i m}\left(\hat{\boldsymbol{\theta}}_{m}\right)$ be the quasi-log likelihood evaluated at the relevant estimate for $m=1,2$. Then

$$
\left(\mathscr{L}_{1}-\mathscr{L}_{2}\right) / N \xrightarrow{p} \mathrm{E}\left[\log f_{1}\left(\mathbf{y}_{i} \mid \mathbf{x}_{i} ; \boldsymbol{\theta}_{1}^{*}\right)\right]-\mathrm{E}\left[\log f_{2}\left(\mathbf{y}_{i} \mid \mathbf{x}_{i} ; \boldsymbol{\theta}_{2}^{*}\right)\right]
$$

where the expected values are over the joint distribution $\left(\mathbf{x}_{i}, \mathbf{y}_{i}\right)$. We can actually say more. Using a mean value expansion and the $\sqrt{N}$-consistency of $\hat{\boldsymbol{\theta}}_{m}^{*}$ for $\boldsymbol{\theta}_{m}^{*}$, it can be shown that


\begin{align*}
N^{-1 / 2}\left(\mathscr{L}_{1}-\mathscr{L}_{2}\right) & =N^{-1 / 2} \sum_{i=1}^{N}\left[\ell_{i 1}\left(\hat{\boldsymbol{\theta}}_{1}\right)-\ell_{i 2}\left(\hat{\boldsymbol{\theta}}_{2}\right)\right] \\
& =N^{-1 / 2} \sum_{i=1}^{N}\left[\ell_{i 1}\left(\boldsymbol{\theta}_{1}^{*}\right)-\ell_{i 2}\left(\boldsymbol{\theta}_{2}^{*}\right)\right]+o_{p}(1) \tag{13.72}
\end{align*}


Equation (13.72) is the key to obtaining a model specification test because it shows that the estimators $\hat{\boldsymbol{\theta}}_{1}$ and $\hat{\boldsymbol{\theta}}_{2}$ do not affect that asymptotic distribution of $N^{-1 / 2}\left(\mathscr{L}_{1}-\mathscr{L}_{2}\right)$. Therefore, we can obtain an asymptotic normal distribution for\\
$N^{-1 / 2}\left(\mathscr{L}_{1}-\mathscr{L}_{2}\right)$ under the null hypothesis\\
$\mathrm{H}_{0}: \mathrm{E}\left[\ell_{i 1}\left(\boldsymbol{\theta}_{1}^{*}\right)\right]=\mathrm{E}\left[\ell_{i 2}\left(\boldsymbol{\theta}_{2}^{*}\right)\right]$.

In particular, under (13.73),\\
$N^{-1 / 2} \sum_{i=1}^{N}\left[\ell_{i 1}\left(\boldsymbol{\theta}_{1}^{*}\right)-\ell_{i 2}\left(\boldsymbol{\theta}_{2}^{*}\right)\right] \xrightarrow{d} \operatorname{Normal}\left(0, \eta^{2}\right)$,\\
where $\eta^{2}=\operatorname{Var}\left[\ell_{i 1}\left(\boldsymbol{\theta}_{1}^{*}\right)-\ell_{i 2}\left(\boldsymbol{\theta}_{2}^{*}\right)\right]$. A consistent estimator of $\eta^{2}$ is\\
$\hat{\eta}^{2} \equiv N^{-1} \sum_{i=1}^{N}\left[\ell_{i 1}\left(\hat{\boldsymbol{\theta}}_{1}\right)-\ell_{i 2}\left(\hat{\boldsymbol{\theta}}_{2}\right)\right]^{2}$.

Voung's model selection statistic is


\begin{align*}
V M S & =N^{-1 / 2}\left(\mathscr{L}_{1}-\mathscr{L}_{2}\right) / \hat{\eta} \\
& =\frac{N^{-1} \sum_{i=1}^{N}\left[\ell_{i 1}\left(\hat{\boldsymbol{\theta}}_{1}\right)-\ell_{i 2}\left(\hat{\boldsymbol{\theta}}_{2}\right)\right]}{\left\{N^{-1} \sum_{i=1}^{N}\left[\ell_{i 1}\left(\hat{\boldsymbol{\theta}}_{1}\right)-\ell_{i 2}\left(\hat{\boldsymbol{\theta}}_{2}\right)\right]^{2}\right\}^{1 / 2} / \sqrt{N}} \stackrel{d}{\rightarrow} \operatorname{Normal}(0,1), \tag{13.76}
\end{align*}


where the standard normal distribution holds under $\mathrm{H}_{0}$. Note that, while the numerator is simply the difference in log likelihoods from the two estimated models, the denominator requires computing the squared difference in the log likelihoods for each $i$. A simple way to obtain a valid test is to define $\hat{d}_{i}=\ell_{i 1}\left(\hat{\boldsymbol{\theta}}_{1}\right)-\ell_{i 2}\left(\hat{\boldsymbol{\theta}}_{2}\right)$ for each $i$ and then simply to regress $\hat{d}_{i}$ on unity to test that its mean is different from zero. (This version of the statistic subtracts off the sample average of $\left\{\hat{d}_{i}: i=1,2, \ldots N\right\}$ in forming the variance estimate in the denominator.)

In applying Vuong's approach, it is important to understand the scope of its application, including the underlying null hypothesis. We should not use the VMS statistic and its limiting standard normal distribution for testing nested models under correct specification. Recall that the LR statistic is simply $L R=2\left(\mathscr{L}_{u r}-\mathscr{L}_{r}\right)$, and, under the null, $L R$ has a limiting $\chi_{Q}^{2}$ distribution, where $Q$ is the number of restrictions. The important point is that, if the models are nested and correctly specified, then $\ell_{i 1}\left(\boldsymbol{\theta}_{1}^{*}\right)-\ell_{i 2}\left(\boldsymbol{\theta}_{2}^{*}\right)=\ell_{i}\left(\boldsymbol{\theta}_{\mathrm{o}}\right)-\ell_{i}\left(\boldsymbol{\theta}_{\mathrm{o}}\right)=0$. That is, the difference in log likelihoods evaluated at the plims of the estimators is identically zero under the null. This degeneracy makes the asymptotic equivalence in equation (13.72) useless for deriving a test statistic because the variance $\eta^{2}$ would be identically zero. For nested models, we do not divide the difference in $\log$ likelihoods by $\sqrt{N}$ because the result would be a statistic that converges in probability to zero.

The sense in which the models must be nonnested to apply Vuong's approach is that\\
$\mathbf{P}\left[\ell_{i 1}\left(\boldsymbol{\theta}_{1}^{*}\right) \neq \ell_{i 2}\left(\boldsymbol{\theta}_{2}^{*}\right)\right]>0$.

In other words, the $\log$ likelihoods evaluated at the psuedo-true values $\boldsymbol{\theta}_{1}^{*}$ and $\boldsymbol{\theta}_{2}^{*}$ must differ for a nontrivial set of outcomes on $\left(\mathbf{x}_{i}, \mathbf{y}_{i}\right)$. This not only rules out models that are obviously nested, but it rules out other degeneracies, too. For example, if $y_{i}$ is a count variable, $\mathbf{x}_{i}=\left(1, x_{i 2}, \ldots, x_{i K}\right)$, and we specify different Poisson distributions-the first with mean function $\exp \left(\mathbf{x}_{i} \boldsymbol{\theta}\right)$ and the second with mean function $\left(\mathbf{x}_{i} \boldsymbol{\theta}\right)^{2}$-these models are nonnested provided that the mean of $y_{i}$ given $\mathbf{x}_{i}$ actually depends on the nonconstant elements in $\mathbf{x}_{i}$. But if $\mathrm{E}\left(y_{i} \mid \mathbf{x}_{i}\right)=\mathrm{E}\left(y_{i}\right)$, then $f_{1}\left(y \mid \mathbf{x} ; \boldsymbol{\theta}_{1}^{*}\right)$ and $f_{2}\left(y \mid \mathbf{x} ; \boldsymbol{\theta}_{2}^{*}\right)$ are Poisson distributions with the same (constant) means, and the limiting standard normal distribution for Vuong's statistic fails. On the other hand, if the competing models are Poisson and geometric, even with the same mean function, say $\exp \left(\mathbf{x}_{i} \boldsymbol{\theta}\right)$, the models are nonnested no matter what, because the Poisson and geometric distributions differ even if they both have constant means.

Because the models must be nonnested, the null hypothesis in equation (13.73) can only hold if both models are misspecified. If one model were correctly specified, yet the densities differed, then we would have a strict inequality in (13.73) in favor of the correctly specified model. We can summarize when it is appropriate to apply Vuong's test: it applies to nonnested models where the null hypothesis is that both models are misspecified yet fit equally well (in the sense that they have the same expected $\log$ likelihoods).

If we reject model 2 in favor of model 1 because VMS is statistically greater than zero, then we can only conclude that model 1 fits better in the sense that $\mathrm{E}\left[\ell_{i 1}\left(\boldsymbol{\theta}_{1}^{*}\right)\right]>\mathrm{E}\left[\ell_{i 2}\left(\boldsymbol{\theta}_{2}^{*}\right)\right]$. It does not mean that model 1 is correctly specified (although it could be). There are many models that can fit better than a given model, and clearly not all can be correct.

Naturally, Vuong's approach applies directly to panel data methods when two complete densities have been specified for $\mathrm{D}\left(\mathbf{y}_{i 1}, \ldots, \mathbf{y}_{i T} \mid \mathbf{x}_{i 1}, \ldots, \mathbf{x}_{i T}\right)$. But it can also be extended to partial (pooled) MLEs, provided we properly account for the time series dependence. For each $t$, let $f_{t 1}\left(\mathbf{y}_{t} \mid \mathbf{x}_{t} ; \boldsymbol{\theta}_{1}\right)$ and $f_{t 2}\left(\mathbf{y}_{t} \mid \mathbf{x}_{t} ; \boldsymbol{\theta}_{2}\right)$ be competing models of the conditional density in each time period. As in Sections 13.8 and 13.9, the log likelihoods are


\begin{equation*}
\ell_{i m}\left(\boldsymbol{\theta}_{m}\right)=\sum_{t=1}^{T} \log f_{t m}\left(\mathbf{y}_{i t} \mid \mathbf{x}_{i t} ; \boldsymbol{\theta}_{m}\right)=\sum_{t=1}^{T} \ell_{i t m}\left(\boldsymbol{\theta}_{m}\right), \quad m=1,2 . \tag{13.78}
\end{equation*}


The same null hypothesis, (13.73), makes sense in the PMLE setting (and is the weakest sense in which the models fit equally well). Moreover, the convergence result in equation (13.74) still holds under the null. Assuming the nonnested condition (13.77) is satisfied, the variance $\eta^{2}$ is positive. However, in estimating $\eta^{2}$, we must account for the serial dependence in $\left\{\ell_{i t 1}\left(\boldsymbol{\theta}_{1}^{*}\right)-\ell_{i t 2}\left(\boldsymbol{\theta}_{2}^{*}\right): t=1, \ldots, T\right\}$. Let $\hat{d}_{i t}= \ell_{i t 1}\left(\hat{\boldsymbol{\theta}}_{1}\right)-\ell_{i t 2}\left(\hat{\boldsymbol{\theta}}_{2}\right)$ denote the difference in estimated log likelihoods for each $t$, and let $\hat{\lambda}_{t}=N^{-1} \sum_{i=1}^{N} \hat{d}_{i t}$. Then $\hat{\eta}^{2}$ is easily obtained as


\begin{equation*}
\hat{\eta}^{2}=N^{-1} \sum_{i=1}^{N}\left\{\sum_{t=1}^{T}\left(\hat{d}_{i t}-\hat{\lambda}_{t}\right)^{2}+\sum_{t=1}^{T} \sum_{r \neq t}^{T}\left(\hat{d}_{i t}-\hat{\lambda}_{t}\right)\left(\hat{d}_{i r}-\hat{\lambda}_{t}\right)\right\} . \tag{13.79}
\end{equation*}


This variance estimator allows for the possibility that the mean difference in log likelihoods varies across $t$ under the null, but that the averages across $t$ are the same. If the null hypothesis is the stronger version, $\mathrm{E}\left[\ell_{i t 1}\left(\boldsymbol{\theta}_{1}^{*}\right)\right]=\mathrm{E}\left[\ell_{i t 2}\left(\boldsymbol{\theta}_{2}^{*}\right)\right]$ for $t=1, \ldots, T$, then $\hat{\lambda}_{t}$ can be replaced with the average of $\hat{d}_{i t}$ across $i$ and $t$, say $\hat{\lambda}$. In this case, the test statistic is simply the $t$ statistic $\hat{\lambda} / \operatorname{se}(\hat{\lambda})$, where $\operatorname{se}(\hat{\lambda})$ is the heteroskedasticity and serial correlation robust standard error from the pooled regression $\hat{d}_{i t}$ on 1 , $t=1, \ldots, T ; i=1, \ldots, N$.

Vuong's model selection test should not be confused with specification tests in the context of nonnested models. For example, the Cox $(1961,1962)$ approach tests a specified model against a nonnested alternative, and a key component of the test is the average difference in log likelihoods, $\left(\mathscr{L}_{1}-\mathscr{L}_{2}\right) / N$. But with Cox's approach, one model is taken to be the correct model under the null hypothesis. (And, in practice, each of two models is taken to be the null model and then tested against the other.) If one model is correct and the models are truly nonnested, then the expected values of the log likelihoods, evaluated at the plims of the MLE (for the null model) and quasi-MLE (for the alternative model), must differ; computation of the Cox statistic requires estimating the mean of the difference in log likelihoods, conditional on $\mathbf{x}_{i}$, under the null hypothesis. In some cases, including when both distributions are normal but with different means or variances (or both), and binary response models, the Cox statistic is easy to compute. But generally finding the conditional mean difference in the log likelihoods in closed form is intractable. (The mean can be simulated, but how much trouble does one want to go through for a specification test?)

The Cox test can be cast as a conditional moment test when we extend the framework in Section 13.7 to allow for the moment conditions to depend on parameters in addition to $\boldsymbol{\theta}$, as done in White (1994, Chap. 9). For further discussion and several\\
other approaches to testing a model against a nonnested alternative, see Gourieroux and Monfort (1994) and White (1994, Chap. 10).

\subsection*{13.11.3 Quasi-Maximum Likelihood Estimation in the Linear Exponential Family}
Section 13.11.1, and the model selection test in the previous subsection, allowed that nothing about the model $f(\mathbf{y} \mid \mathbf{x} ; \boldsymbol{\theta})$, or $f_{t}\left(\mathbf{y}_{t} \mid \mathbf{x}_{t} ; \boldsymbol{\theta}\right), t=1, \ldots, T$, is correctly specified. Sometimes, a useful middle ground between correct specification and complete misspecification is to allow just one or two features of a conditional density to be correctly specified. In this section, we study a class of QMLEs that consistently estimate the parameters in a correctly specified conditional mean.

As one learns in introductory econometrics-although it may not be stated in quite this way-the ordinary least squares (OLS) estimator is a QMLE: OLS can be obtained by maximimizing the Gaussian (normal) log-likelihood function with a conditional mean linear in parameters; in fact, we can arbitrarily set the variance equal to any fixed value without affecting the estimates of the mean parameters. Not surprisingly, the nonlinear least squares estimator we covered in Chapter 12 has the same property: minimizing the sum of squared residuals is the same as maximizing the Gaussian quasi-log likelihood. Therefore, NLS is a QMLE based on the normal density function.

The normal $\log$ likelihood is not the only log likelihood that identifies the parameters of a conditional mean despite arbitrary misspecification of the remaining features of the distribution. The Bernoulli log likelihood, the Poisson log likelihood, the exponential log likelihood, and others share these feature. These log likelihoods are all members of the linear exponential family (LEF), and it is useful to provide a somewhat general treatment of this class of QMLEs. For simplicity, we consider the case of a scalar response, even though there are results for multiple responses, too. We draw on Gourieroux, Monfort, and Trognon (1984a), or GMT.

A log likelihood in the LEF can be written as a function of the mean as


\begin{equation*}
\log f(y \mid \mu)=a(\mu)+b(y)+y c(\mu), \tag{13.80}
\end{equation*}


for functions $a(\cdot), b(\cdot)$, and $c(\cdot)$, where $\mu$ is a candidate value of the mean of $y_{i}$ and $\mu_{\mathrm{o}}=\mathrm{E}\left(y_{i}\right)$ is the true population mean. Let $M$ denote the set of possible values of the mean. GMT (1984a) show that $\mu_{\mathrm{o}}$ solves


\begin{equation*}
\max _{\mu \in M}\left[a(\mu)+\mathrm{E}\left(y_{i}\right) c(\mu)\right]=\max _{\mu \in M}\left[a(\mu)+\mu_{\mathrm{o}} c(\mu)\right] . \tag{13.81}
\end{equation*}


The functions in (13.80) are easily obtained for the normal, Bernoulli, Poisson, exponential, and other cases; GMT (1984a) contains a summary table. For the Bernoulli distribution,

$$
\begin{aligned}
\log f(y \mid \mu) & =(1-y) \cdot \log (1-\mu)+y \cdot \log (\mu) \\
& =\log (1-\mu)+y \cdot \log [\mu /(1-\mu)], \quad 0<\mu<1
\end{aligned}
$$

Therefore, $a(\mu)=\log (1-\mu), b(y)=0$, and $c(\mu)=\log [\mu /(1-\mu)]$. For some distributions to fit into the LEF, notably the gamma and negative binomial, a nuisance parameter must be fixed at a specific value; see GMT (1984a) for details.

In the most popular examples, it is easy to directly verify that $\mu_{\mathrm{o}}$ solves (13.81). We will consider special cases in Part IV. For now, it is important to understand the meaning of the result. Consider the Bernoulli case but, rather than assuming $y_{i}$ is a zero-one variable, let $y_{i}$ be any random variable with support in the unit interval, $[0,1] . y_{i}$ can be discrete, continuous, or have both features. For example, we could have $\mathrm{P}\left(y_{i}=0\right)>0$ but $\mathrm{P}\left(y_{i}=y\right)=0$ for $y \in(0,1]$, or $y_{i}$ might take on values in $\left\{0,1 / m_{i}, 2 / m_{i}, \ldots, 1\right\}$ for some positive integer $m_{i}$. Regardless of the nature of $y_{i}$, provided its mean $\mu_{\mathrm{o}}$ is in $(0,1), \mu_{\mathrm{o}}$ maximizes the expected value of the Bernoulli log likelihood.

In practice, we are interested in conditional rather than unconditional means, which we parameterize as $m(\mathbf{x}, \boldsymbol{\theta})$. Then the conditional quasi-log likelihood function becomes


\begin{equation*}
\log f(y \mid m(\mathbf{x}, \boldsymbol{\theta}))=a(m(\mathbf{x}, \boldsymbol{\theta}))+b(y)+y c(m(\mathbf{x}, \boldsymbol{\theta})) . \tag{13.82}
\end{equation*}


Because the mean is now assumed to be correctly specified, we assume there is $\boldsymbol{\theta}_{\mathrm{o}} \in \boldsymbol{\Theta}$ such that $\mathrm{E}\left(y_{i} \mid \mathbf{x}_{i}\right)=m\left(\mathbf{x}_{i}, \boldsymbol{\theta}_{\mathrm{o}}\right)$. A simple iterated expectations argument shows that $\boldsymbol{\theta}_{\mathrm{o}}$ solves


\begin{equation*}
\max _{\boldsymbol{\theta} \in \boldsymbol{\Theta}} \mathrm{E}\left[a\left(m\left(\mathbf{x}_{i}, \boldsymbol{\theta}\right)\right)+y_{i} c\left(m\left(\mathbf{x}_{i}, \boldsymbol{\theta}\right)\right)\right], \tag{13.83}
\end{equation*}


regardless of the actual distribution $\mathrm{D}\left(y_{i} \mid \mathbf{x}_{i}\right)$. Again, it is important to understand the meaning of this result. The nature of $y_{i}$ need not even correspond to the chosen density. For example, $y_{i}$ could be a nonnegative, continuous variable, and we use the Poisson quasi-log likelihood. The Poisson QMLE is consistent for the conditional mean parameters provided the mean-with the leading case being an exponential function-is correctly specified. The only restriction is that candidates for $\mathrm{E}\left(y_{i} \mid \mathbf{x}_{i}=\mathbf{x}\right)$ should have the same range as allowed in the chosen LEF density.

One of the most useful characterizations of QMLE in the LEF is based on the score. It can be shown that the score has the form


\begin{equation*}
\mathbf{s}_{i}(\boldsymbol{\theta})=\nabla_{\boldsymbol{\theta}} m\left(\mathbf{x}_{i}, \boldsymbol{\theta}\right)^{\prime}\left[y_{i}-m\left(\mathbf{x}_{i}, \boldsymbol{\theta}\right)\right] / v\left(m\left(\mathbf{x}_{i}, \boldsymbol{\theta}\right)\right) \tag{13.84}
\end{equation*}


where $\nabla_{\boldsymbol{\theta}} m\left(\mathbf{x}_{i}, \boldsymbol{\theta}\right)$ is the $1 \times P$ gradient of the mean function and, of importance, $v(\mu)$ is the variance function associated with the chosen LEF density. For the standard normal, $v(\mu)=1$, for the Bernoulli, $v(\mu)=\mu(1-\mu)$, for the Poisson, $v(\mu)=\mu$, and for the exponential, $v(\mu)=\mu^{2}$. The structure in equation (13.84) shows immediately that the QMLE is Fisher consistent: if $\mathrm{E}\left(y_{i} \mid \mathbf{x}_{i}\right)=m\left(\mathbf{x}_{i}, \boldsymbol{\theta}_{\mathrm{o}}\right)$, then $\mathrm{E}\left[\mathbf{s}_{i}\left(\boldsymbol{\theta}_{\mathrm{o}}\right) \mid \mathbf{x}_{i}\right]=\mathbf{0}$, which in turn implies that the unconditional mean of the score is zero. We derived (13.84) explicitly in the probit example (Example 13.1), and it is also easily seen to be true for the Poisson regression (Example 13.2). But here we emphasize that $\mathrm{E}\left[\mathbf{s}_{i}\left(\boldsymbol{\theta}_{\mathrm{o}}\right) \mid \mathbf{x}_{i}\right]=\mathbf{0}$ holds for arbitrary misspecification of these densities provided the mean is correctly specified.

We can also use the score to compute the expected Hessian conditional on $\mathbf{x}_{i}$ :


\begin{equation*}
\mathbf{A}\left(\mathbf{x}_{i}, \boldsymbol{\theta}_{\mathrm{o}}\right)=-\mathrm{E}\left[\mathbf{H}_{i}\left(\boldsymbol{\theta}_{\mathrm{o}}\right) \mid \mathbf{x}_{i}\right]=\nabla_{\boldsymbol{\theta}} m\left(\mathbf{x}_{i}, \boldsymbol{\theta}_{\mathrm{o}}\right)^{\prime} \nabla_{\boldsymbol{\theta}} m\left(\mathbf{x}_{i}, \boldsymbol{\theta}_{\mathrm{o}}\right) / v\left(m\left(\mathbf{x}_{i}, \boldsymbol{\theta}_{\mathrm{o}}\right)\right) . \tag{13.85}
\end{equation*}


Further,


\begin{equation*}
\mathrm{E}\left[\mathbf{s}_{i}\left(\boldsymbol{\theta}_{\mathrm{o}}\right) \mathbf{s}_{i}\left(\boldsymbol{\theta}_{\mathrm{o}}\right) \mid \mathbf{x}_{i}\right]=\mathrm{E}\left(u_{i}^{2} \mid \mathbf{x}_{i}\right) \nabla_{\boldsymbol{\theta}} m\left(\mathbf{x}_{i}, \boldsymbol{\theta}_{\mathrm{o}}\right)^{\prime} \nabla_{\boldsymbol{\theta}} m\left(\mathbf{x}_{i}, \boldsymbol{\theta}_{\mathrm{o}}\right) /\left[v\left(m\left(\mathbf{x}_{i}, \boldsymbol{\theta}_{\mathrm{o}}\right)\right)\right]^{2}, \tag{13.86}
\end{equation*}


where $u_{i} \equiv y_{i}-m\left(\mathbf{x}_{i}, \boldsymbol{\theta}_{\mathrm{o}}\right)$. It follows immediately that the CIME holds if $\mathrm{E}\left(u_{i}^{2} \mid \mathbf{x}_{i}\right)= v\left(m\left(\mathbf{x}_{i}, \boldsymbol{\theta}_{\mathrm{o}}\right)\right)$, that is,


\begin{equation*}
\operatorname{Var}\left(y_{i} \mid \mathbf{x}_{i}\right)=v\left(m\left(\mathbf{x}_{i}, \boldsymbol{\theta}_{\mathrm{o}}\right)\right) . \tag{13.87}
\end{equation*}


(In the Bernoulli case, $v(m)=m(1-m)$, and in the Poisson case, $v(m)=m$.) In other words, if the chosen LEF density has a conditional variance equal to the actual $\operatorname{Var}\left(y_{i} \mid \mathbf{x}_{i}\right)$, then we can use the usual MLE standard errors and inference (even if features of the distribution other than the first two conditional moments are misspecified). So, for example, in a Poisson regression analysis, if $\operatorname{Var}\left(y_{i} \mid \mathbf{x}_{i}\right)=\mathrm{E}\left(y_{i} \mid \mathbf{x}_{i}\right)$ and the mean function is correctly specified, we can act as if we are using MLE rather than QMLE, even if higher order conditional moments of $y_{i}$ do not match up with the Poisson distribution.

If $\operatorname{Var}\left(y_{i} \mid \mathbf{x}_{i}\right)$ is unrestricted, the IME will not hold, and then the robust sandwich estimator


\begin{equation*}
\widehat{\operatorname{Avar}(\hat{\boldsymbol{\theta}})}=\left(\sum_{i=1}^{N} \mathbf{A}\left(\mathbf{x}_{i}, \hat{\boldsymbol{\theta}}\right)\right)^{-1}\left(\sum_{i=1}^{N} \mathbf{s}_{i}(\hat{\boldsymbol{\theta}}) \mathbf{s}_{i}(\hat{\boldsymbol{\theta}})^{\prime}\right)\left(\sum_{i=1}^{N} \mathbf{A}\left(\mathbf{x}_{i}, \hat{\boldsymbol{\theta}}\right)\right)^{-1}, \tag{13.88}
\end{equation*}


should be used, where $\mathbf{A}\left(\mathbf{x}_{i}, \hat{\boldsymbol{\theta}}\right)$ is given in (13.85) with $\hat{\boldsymbol{\theta}}$ replacing $\boldsymbol{\theta}_{\mathrm{o}}$. Because of the particular structure of the log likelihood and the assumption that the conditional mean is correctly specified, we can find $\mathrm{E}\left[\mathbf{H}_{i}\left(\boldsymbol{\theta}_{\mathrm{o}}\right) \mid \mathbf{x}_{i}\right]$. If we did not want to assume\\
correct specification of the conditional mean, we would be in the setup of Section 13.11.1.

QMLE in the LEF is closely related to the so-called generalized linear model (GLM) literature in statistics. The terminology and some particulars differ, and the early GLM literature did not recognize the robustness of the approach for estimating conditional mean parameters, but in modern applications the key feature is that they both use QMLE to estimate parameters of a conditional mean. The GLM approach is more restrictive in that the conditional mean is assumed to have an index structure. In particular, the mean is assumed to have the form $m(\mathbf{x}, \boldsymbol{\theta})=r(\mathbf{x} \boldsymbol{\theta})$, where the "index" $\mathbf{x} \boldsymbol{\theta}$ is linear in parameters and $r(\cdot)$ is a function of the index. In addition, an important component of the GLM apparatus is the link function, which implicitly defines the mean function. If we let $\eta$ denote the index $\mathbf{x} \boldsymbol{\theta}$, then the link function $g(\cdot)$ is such that $\eta=g(\mu)$. The link function is strictly monotonic and therefore has an inverse, and so $\mu=g^{-1}(\eta)$ or, in the notation of conditional mean functions, $m(\mathbf{x}, \boldsymbol{\theta})=g^{-1}(\mathbf{x} \boldsymbol{\theta})$.

The term "generalized linear model" comes from the underlying linearity of the index function, and then the link function introduces nonlinearity. In most applications, it is more natural to specify the conditional mean function because we want the mean function to be consistent with the nature of $y_{i}$, and $y_{i}$ is the outcome we hope to explain. (So, for example, if $y_{i}$ is nonnegative, we want $m(\mathbf{x}, \boldsymbol{\theta})$ to be positive for all $\mathbf{x}$ and $\boldsymbol{\theta}$; if $0 \leq y_{i} \leq 1$, we want $m(\mathbf{x}, \boldsymbol{\theta})$ to be in the unit interval.) Once the mean function is specified, we use a suitable LEF density; this is the approach taken by GMT (1984a). Directly specifying $m(\mathbf{x}, \boldsymbol{\theta})$ does not wed one to the index structure, although in most applications, $m(\mathbf{x}, \boldsymbol{\theta})$ has an index form. If, say, $m(\mathbf{x}, \boldsymbol{\theta})=\exp (\mathbf{x} \boldsymbol{\theta})$ then the link function is $g(\mu)=\log (\mu)$ for $\mu>0$. If $m(\mathbf{x}, \boldsymbol{\theta})=\exp (\mathbf{x} \boldsymbol{\theta}) /[1+\exp (\mathbf{x} \boldsymbol{\theta})]$, then $g(\mu)=\log [\mu /(1-\mu)]$ for $0<\mu<1$. McCullagh and Nelder (1989) is a good reference for GLM.

The GLM literature recognized early on that assumption (13.87) was too restrictive for many applications. As we discussed, no variance assumption is needed for consistent estimation of $\boldsymbol{\theta}_{\mathrm{o}}$. An assumption that has been used in the GLM literature allows $\operatorname{Var}\left(y_{i} \mid \mathbf{x}_{i}\right)$ to differ from that implied by the LEF distribution by a constant:


\begin{equation*}
\operatorname{Var}\left(y_{i} \mid \mathbf{x}_{i}\right)=\sigma_{\mathrm{o}}^{2} v\left(m\left(\mathbf{x}_{i}, \boldsymbol{\theta}_{\mathrm{o}}\right)\right) \tag{13.89}
\end{equation*}


for some $\sigma_{\mathrm{o}}^{2}>0$, which is often called the dispersion parameter. Because of its historical role in GLM analysis, we refer to (13.89) as the GLM variance assumption. When $\sigma_{\mathrm{o}}^{2}>1$, then we say there is overdispersion (relative to the chosen density); underdispersion is when $\sigma_{\mathrm{o}}^{2}<1$, and both cases arise in practice.

Under (13.89), it is straightforward to estimate $\sigma_{\mathrm{o}}^{2}$. Let $u_{i}=y_{i}-m\left(\mathbf{x}_{i}, \boldsymbol{\theta}_{\mathrm{o}}\right)$ be the additive "errors." Then, because $\mathrm{E}\left(u_{i}^{2} \mid \mathbf{x}_{i}\right)=\sigma_{\mathrm{o}}^{2} v\left(m\left(\mathbf{x}_{i}, \boldsymbol{\theta}_{\mathrm{o}}\right)\right) \equiv \sigma_{\mathrm{o}}^{2} v_{i}$, it follows by iterated expectations that


\begin{equation*}
\mathrm{E}\left(u_{i}^{2} / v_{i}\right)=\mathrm{E}\left[\mathrm{E}\left(u_{i}^{2} / v_{i} \mid \mathbf{x}_{i}\right)\right]=\mathrm{E}\left[\mathrm{E}\left(u_{i}^{2} \mid \mathbf{x}_{i}\right) / v_{i}\right]=\mathrm{E}\left(\sigma_{\mathrm{o}}^{2} v_{i} / v_{i}\right)=\sigma_{\mathrm{o}}^{2} . \tag{13.90}
\end{equation*}


Therefore, by the usual analogy principle argument,


\begin{equation*}
\hat{\sigma}^{2}=(N-P)^{-1} \sum_{i=1}^{N} \hat{u}_{i}^{2} / \hat{v}_{i} \tag{13.91}
\end{equation*}


is consistent for $\sigma_{\mathrm{o}}^{2}$, where $\hat{u}_{i} \equiv y_{i}-m\left(\mathbf{x}_{i}, \hat{\boldsymbol{\theta}}\right)$ are the residuals, $\hat{v}_{i} \equiv v\left(m\left(\mathbf{x}_{i}, \hat{\boldsymbol{\theta}}\right)\right)$ are the estimated conditional variances from the LEF density, and the degrees-offreedom adjustment is common (but, of course, does not affect consistency). In the GLM literature, the standardized residuals $\hat{u}_{i} / \sqrt{\hat{v}_{i}}$ are called the Pearson residuals and the estimate in equation (13.91) is the Pearson dispersion estimator.

Under the GLM variance assumption, it is easily seen that the generalized IME, given in equation (12.53), is satisfied (once we account for the minor difference between a minimization and maximization problem). In fact, a conditional version holds, which we can call the generalized conditional information matrix equality (GCIME):


\begin{equation*}
\mathrm{E}\left[\mathbf{s}_{i}\left(\boldsymbol{\theta}_{\mathrm{o}}\right) \mathbf{s}_{i}\left(\boldsymbol{\theta}_{\mathrm{o}}\right)^{\prime} \mid \mathbf{x}_{i}\right]=-\sigma_{\mathrm{o}}^{2} \mathrm{E}\left[\mathbf{H}_{i}\left(\boldsymbol{\theta}_{\mathrm{o}}\right) \mid \mathbf{x}_{i}\right]=\sigma_{\mathrm{o}}^{2} \mathbf{A}\left(\mathbf{x}_{i}, \boldsymbol{\theta}_{\mathrm{o}}\right) . \tag{13.92}
\end{equation*}


From Chapter 12, we can take


\begin{equation*}
\widehat{\operatorname{Avar}(\hat{\boldsymbol{\theta}})}=\hat{\sigma}^{2}\left(\sum_{i=1}^{N} \mathbf{A}\left(\mathbf{x}_{i}, \hat{\boldsymbol{\theta}}\right)\right)^{-1}=\hat{\sigma}^{2}\left(\sum_{i=1}^{N} \nabla_{\boldsymbol{\theta}} m_{i}(\hat{\boldsymbol{\theta}})^{\prime} \nabla_{\boldsymbol{\theta}} m_{i}(\hat{\boldsymbol{\theta}}) / v_{i}(\hat{\boldsymbol{\theta}})\right)^{-1}, \tag{13.93}
\end{equation*}


where the notation should be clear. Most software packages that have GLM commands-typically requiring one to specify the LEF density and the link function-allow computation of the nonrobust variance matrix estimator that assumes (13.87) (which is the same variance estimate under full MLE), the estimator in (13.93), or the fully robust form in (13.88). Therefore, the GLM framework can be used to obtain QMLEs for the LEF for a certain class of mean functions.

Notice how similar the structure of equation (13.93) is to the asymptotic variance of the weighted nonlinear least squares estimator from Chapter 12; see equation (12.59). In fact, the only difference between (13.93) and (12.59) is that with WNLS we allow the parameters in the variance function to be different from $\boldsymbol{\theta}_{\mathrm{O}}$, and, therefore, to be estimated in a separate stage (usually after an initial NLS estimation). The\\
similarity between (13.93) and (12.59) hints at a close link between QMLEs in the LEF and WNLS. In fact, for every QMLE in the LEF, there is an asymptotically equivalent WNLS estimator when the conditional mean is correctly specified: for WNLS, take $h(\mathbf{x}, \boldsymbol{\gamma})=v(m(\mathbf{x}, \boldsymbol{\theta}))$ and $\hat{\gamma}$ an initial $\sqrt{N}$-consistent estimator of $\boldsymbol{\theta}_{\mathrm{o}}$ (such as NLS). The QMLE and WNLS estimators will be $\sqrt{N}$-equivalent whether or not (13.89) holds.

If we nominally wish to assume (13.89), QMLE is computationally convenient because it is a one-step estimation procedure, and the objective function is usually well behaved; WNLS requires two steps. Therefore, when the variance and mean are thought to be related in a way suggested by an LEF density-more precisely, its GLM extension in (13.89)-QMLE is almost always used. And, of course, we can make inference fully robust to any variance-mean relationship by using the sandwich estimator (13.88).

\subsection*{13.11.4 Generalized Estimating Equations for Panel Data}
Naturally, we can extend QMLE in the LEF to panel data. The simplest approach is to specify conditional mean functions $m_{t}\left(\mathbf{x}_{t}, \boldsymbol{\theta}\right), t=1, \ldots, T$, along with an LEF density, and then to proceed with estimation by ignoring any time dependence. The pooled quasi-likelihood has the same form as equation (13.45), where $f_{t}\left(y_{t} \mid \mathbf{x}_{t} ; \boldsymbol{\theta}\right)$ is in the LEF. (Allowing the mean and density functions to depend on $t$ is probably not necessary, but it reminds us that we often want to allow some parameters to change over $t$-for example, $m_{t}\left(\mathbf{x}_{t}, \boldsymbol{\theta}\right)=\exp \left(\alpha_{t}+\mathbf{x}_{t} \boldsymbol{\beta}\right)$. In practice, we would accomplish different "intercepts" within the exponential function by including a full set of time period dummies among the regressors.)

Correct specification of the mean for each $t$ means that, for some $\boldsymbol{\theta}_{\mathrm{O}}$,


\begin{equation*}
\mathrm{E}\left(y_{i t} \mid \mathbf{x}_{i t}\right)=m_{t}\left(\mathbf{x}_{i t}, \boldsymbol{\theta}_{\mathrm{o}}\right), \quad t=1, \ldots, T \tag{13.94}
\end{equation*}


Notice that (13.94) does not assume strict exogeneity of $\left\{\mathbf{x}_{i t}: t=1, \ldots, T\right\}$. The score for each $t$ has the same form as equation (13.84):


\begin{equation*}
\mathbf{s}_{i t}(\boldsymbol{\theta})=\nabla_{\boldsymbol{\theta}} m_{t}\left(\mathbf{x}_{i t}, \boldsymbol{\theta}\right)^{\prime}\left[y_{i t}-m_{t}\left(\mathbf{x}_{i t}, \boldsymbol{\theta}\right)\right] / v\left(m_{t}\left(\mathbf{x}_{i t}, \boldsymbol{\theta}\right)\right), \tag{13.95}
\end{equation*}


and the partial QMLE (or pooled QMLE) is generally found by solving $\sum_{i=1}^{N} \sum_{t=1}^{T} \mathbf{s}_{i t}(\hat{\boldsymbol{\theta}})=\mathbf{0}$.

As we discussed in Section 13.8 for partial MLE, the scores $\left\{\mathbf{s}_{i t}\left(\boldsymbol{\theta}_{\mathrm{o}}\right): t=1, \ldots, T\right\}$ are generally serially correlated. However, in Section 13.8.3 we saw an important case in which the scores are not serially correlated: the distribution $\mathrm{D}\left(\mathbf{y}_{i t} \mid \mathbf{x}_{i t}\right)$ is dynamically complete in the sense that it also equals $\mathrm{D}\left(\mathbf{y}_{i t} \mid \mathbf{x}_{i t}, \mathbf{y}_{i, t-1}, \mathbf{x}_{i, t-1}, \ldots, \mathbf{y}_{i 1}, \mathbf{x}_{i 1}\right)$.

Here, we are not assuming that we have a full density $f_{t}\left(\mathbf{y}_{t} \mid \mathbf{x}_{t} ; \boldsymbol{\theta}\right)$ correctly specified. Nevertheless, using equation (13.95), we can easily see that if the conditional mean is dynamically complete in the sense that\\
$\mathrm{E}\left(y_{i t} \mid \mathbf{x}_{i t}\right)=\mathrm{E}\left(y_{i t} \mid \mathbf{x}_{i t}, y_{i, t-1}, \mathbf{x}_{i, t-1}, \ldots, y_{i 1}, \mathbf{x}_{i 1}\right)$,\\
then\\
$\left.\mathrm{E}\left[\mathbf{s}_{i t}\left(\boldsymbol{\theta}_{\mathrm{o}}\right) \mid \mathbf{x}_{i t}, y_{i, t-1}, \mathbf{x}_{i, t-1}, \ldots, y_{i 1}, \mathbf{x}_{i 1}\right)\right]=\mathbf{0}$,\\
and so the scores evaluated at $\boldsymbol{\theta}_{\mathrm{o}}$ are necessarily serially uncorrelated. Importantly, this finding has nothing to do with whether other features of the LEF density, such as the conditional variance, are correctly specified. Without any assumptions on $\operatorname{Var}\left(y_{i t} \mid \mathbf{x}_{i t}\right)$, the appropriate asymptotic variance estimator of the pooled QMLE is\\
$\left(\sum_{i=1}^{N} \sum_{i=1}^{N} \nabla_{\boldsymbol{\theta}} \hat{m}_{i t}^{\prime} \nabla_{\boldsymbol{\theta}} \hat{m}_{i t} / \hat{v}_{i t}\right)^{-1}\left(\sum_{i=1}^{N} \sum_{i=1}^{N} \hat{u}_{i t}^{2} \nabla_{\boldsymbol{\theta}} \hat{m}_{i t}^{\prime} \nabla_{\boldsymbol{\theta}} \hat{m}_{i t} / \hat{v}_{i t}^{2}\right)^{-1}\left(\sum_{i=1}^{N} \sum_{i=1}^{N} \nabla_{\boldsymbol{\theta}} \hat{m}_{i t}^{\prime} \nabla_{\boldsymbol{\theta}} \hat{m}_{i t} / \hat{v}_{i t}\right)^{-1}$,\\
where the notation should be clear. This estimator is precisely what is reported from a pooled QMLE analysis-in most cases using a GLM routine-that ignores the time dimension but allows the variance to be misspecified.

If we further add the GLM variance assumption $\operatorname{Var}\left(y_{i t} \mid \mathbf{x}_{i t}\right)=\sigma_{\mathrm{o}}^{2} v\left(m\left(\mathbf{x}_{i t}, \boldsymbol{\theta}_{\mathrm{o}}\right)\right)$ for all $t$-including that the scale factor $\sigma_{\mathrm{o}}^{2}$ is constant across $t$-then the pooled analogue of equation (13.93) is valid (where $\hat{\sigma}^{2}$ is obtained from the sum of squared standardized residuals across $i$ and $t$ ). Some econometrics packages allow computation of these variances as a routine matter, and they also compute the variance matrix that does not assume dynamic completeness in the conditional mean.

Naturally, if the mean is not dynamically complete, then, if we make the stronger assumption of strict exogeneity of the regressors, it is possible to obtain a more efficient estimator. In Section 12.9, we studied the multivariate WNLS estimator, and noted that a common way to choose a nonconstant "working" variance matrix was to specify models for the conditional variances-which may be misspecified-along with a constant "working" correlation matrix. This is common in the generalized estimating equations (GEE) approach to panel data models. For practical purposes, GEE is a special case of WMNLS, where the variance functions are chosen from the LEF of distributions (and the mean functions are commonly chosen from linear, exponential, logistic, and probit, as in standard GLM analysis).

For the $T \times 1$ vector $\mathbf{y}_{i}$, we assume


\begin{equation*}
\mathrm{E}\left(\mathbf{y}_{i} \mid \mathbf{x}_{i}\right)=\mathbf{m}\left(\mathbf{x}_{i}, \boldsymbol{\theta}_{\mathrm{o}}\right), \quad \text { some } \boldsymbol{\theta}_{\mathrm{o}} \in \boldsymbol{\Theta}, \tag{13.98}
\end{equation*}


where $\mathbf{x}_{i}$ is the collection of all regressors across all time periods and $\mathbf{m}(\mathbf{x}, \boldsymbol{\theta})$ is $T \times 1$. Notice that the vector $\mathrm{E}\left(\mathbf{y}_{i} \mid \mathbf{x}_{i}\right)$ has elements $\mathrm{E}\left(y_{i t} \mid \mathbf{x}_{i}\right)$, which means that the regressors are strictly exogenous. When the mean function at time $t$ does not depend on all of $\mathbf{x}_{i}$, (13.98) implies that the regressors excluded at time $t$ can have no partial effect on $\mathrm{E}\left(y_{i t} \mid \mathbf{x}_{i}\right)$. An important case where all elements of $\mathbf{x}_{i}$ appear for each $t$ is in a correlated random coefficient setup, but then restrictions are obtained using a strict exogeneity assumption conditional on the heterogeneity. We return to this situation below. Generally, one should not apply GEE methods with a nondiagonal working variance matrix (see Section 12.9) unless the regressors satisfy a strict exogeneity assumption.

In the most common applications of GEE, we specify, for each $t$, a quasi-log likelihood from the LEF. As in the cross section case, this choice is motivated by the nature of $y_{i t}$. Assuming we have a suitable parametric model $m\left(\mathbf{x}_{i t}, \boldsymbol{\theta}\right)$ for the mean, and have chosen an LEF, GEE analysis is straightforward once we choose the working correlation matrix (see Section 12.9). Typically, this correlation matrix is constant, depending on parameters $\boldsymbol{\rho}$, so we can write $\mathbf{R}(\boldsymbol{\rho})$ as the $T \times T$ matrix of proposed correlations. Because we do not allow this matrix to depend on $\mathbf{x}_{i}$, there can be no presumption that $\operatorname{Corr}\left(y_{i t}, y_{i s} \mid \mathbf{x}_{i}\right)=r_{t s}\left(\boldsymbol{\rho}_{\mathrm{o}}\right)$. A key feature of GEE is that we explicitly recognize that the conditional correlations might not be constant-in fact, often we know almost certainly they are not-but we apply weighted multivariate nonlinear least squares anyway to hopefully improve efficiency over just a pooled QMLE analysis.

As mentioned in Section 12.9, the two most common working correlation matrices for panel data are the unstructured and the exchangeable. The latter conserves on parameters in the working variance matrix, but with large $N$ (and not so large $T$ ) parsimony might not be necessary. In any case, we need to combine a working correlation matrix with a nominal assumption about the variances. The variances naturally come from the chosen LEF density. So, if we are using the Bernoulli quasi-log likelihood with mean function $m_{t}\left(\mathbf{x}_{i t}, \boldsymbol{\theta}\right)$, then we use the nominal variance $m_{t}\left(\mathbf{x}_{i t}, \boldsymbol{\theta}\right)\left[1-m_{t}\left(\mathbf{x}_{i t}, \boldsymbol{\theta}\right)\right]$. For the Poisson QLL, we use $m_{t}\left(\mathbf{x}_{i t}, \boldsymbol{\theta}\right)$. The mean parameters are estimated by pooled QMLE in a first stage, say $\check{\boldsymbol{\theta}}$. Then the working correlation matrix can be estimated using standardized residuals $\check{u}_{i t} / \sqrt{\check{v}_{i t}}$ (see Section 12.9). Then we can form the working variance matrix, $\hat{\mathbf{W}}_{i}=\hat{\mathbf{V}}_{i}^{1 / 2} \hat{\mathbf{R}} \hat{\mathbf{V}}_{i}^{1 / 2}$ where $\hat{\mathbf{V}}_{i}= \operatorname{diag}\left(v_{i 1}(\check{\boldsymbol{\theta}}), \ldots, v_{i T}(\check{\boldsymbol{\theta}})\right)$ depends on the first-step estimates. We then apply WMNLS, or work off the first-order conditions. It is important to use the asymptotic variance matrix in equation (12.98) that allows the working variance matrix to be misspecified, whether it is due to the variances being misspecified or the working correlation ma-\\
trix being misspecified. (As we discussed in Section 12.9, equation (12.98) is not valid if the conditional mean is misspecified, and the GEE literature sometimes refers to this estimate as a "semirobust" estimate.) If one believes the variance matrix is correct up to a dispersion factor, then $\hat{\sigma}^{2}$ is obtained along with other parameter estimates.

Pooled QMLE and GEE methods are attractive for conditional mean models with unobserved heterogeneity, say $c_{i}$, and strictly exogenous regressors conditional on $c_{i}: \mathrm{E}\left(y_{i t} \mid \mathbf{x}_{i 1}, \ldots, \mathbf{x}_{i T}, c_{i}\right)=\mathrm{E}\left(y_{i t} \mid \mathbf{x}_{i t}, c_{i}\right)$. When we combine this assumption with a specific correlated RE assumption, such as $c_{i}=\alpha+\overline{\mathbf{x}}_{i} \boldsymbol{\xi}+a_{i}$, where $\overline{\mathbf{x}}_{i}$ is the vector of time averages and $a_{i}$ is independent of $\left(\mathbf{x}_{i 1}, \ldots, \mathbf{x}_{i T}\right)$, then we can often find $\mathrm{E}\left(y_{i t} \mid \mathbf{x}_{i 1}, \ldots, \mathbf{x}_{i T}\right)$ as a function of $\left(\mathbf{x}_{i t}, \overline{\mathbf{x}}_{i}\right)$ (or simply assert that it takes on a convenient form). GEE (including pooled QMLE) becomes an indispensable tool to ensure our inference is robust to misspecification of the serial correlation structure we adopt. We will not provide a general treatment now, but we draw heavily on this idea in several chapters in Part IV.

\section*{Problems}
13.1. If $f(\mathbf{y} \mid \mathbf{x} ; \boldsymbol{\theta})$ is a correctly specified model for the density of $\mathbf{y}_{i}$ given $\mathbf{x}_{i}$, does $\boldsymbol{\theta}_{\mathrm{o}}$ solve $\max _{\boldsymbol{\theta} \in \boldsymbol{\Theta}} \mathrm{E}\left[f\left(\mathbf{y}_{i} \mid \mathbf{x}_{i} ; \boldsymbol{\theta}\right)\right]$ ?\\
13.2. Suppose that for a random sample, $y_{i} \mid \mathbf{x}_{i} \sim \operatorname{Normal}\left[m\left(\mathbf{x}_{i}, \boldsymbol{\beta}_{\mathrm{o}}\right), \sigma_{\mathrm{o}}^{2}\right]$, where $m(\mathbf{x}, \boldsymbol{\beta})$ is a function of the $K$-vector of explanatory variables $\mathbf{x}$ and the $P \times 1$ parameter vector $\boldsymbol{\beta}$. Recall that $\mathrm{E}\left(y_{i} \mid \mathbf{x}_{i}\right)=m\left(\mathbf{x}_{i}, \boldsymbol{\beta}_{\mathrm{o}}\right)$ and $\operatorname{Var}\left(y_{i} \mid \mathbf{x}_{i}\right)=\sigma_{\mathrm{o}}^{2}$.\\
a. Write down the conditional $\log$-likelihood function for observation $i$. Show that the CMLE of $\boldsymbol{\beta}_{\mathrm{o}}, \hat{\boldsymbol{\beta}}$, solves the problem $\min _{\boldsymbol{\beta}} \sum_{i=1}^{N}\left[y_{i}-m\left(\mathbf{x}_{i}, \boldsymbol{\beta}\right)\right]^{2}$. In other words, the CMLE for $\boldsymbol{\beta}_{\mathrm{o}}$ is the nonlinear least squares estimator.\\
b. Let $\boldsymbol{\theta} \equiv\left(\boldsymbol{\beta}^{\prime}, \sigma^{2}\right)^{\prime}$ denote the $(P+1) \times 1$ vector of parameters. Find the score of the $\log$ likelihood for a generic $i$. Show directly that $\mathrm{E}\left[\mathbf{s}_{i}\left(\boldsymbol{\theta}_{\mathrm{o}}\right) \mid \mathbf{x}_{i}\right]=\mathbf{0}$. What features of the normal distribution do you need to be correctly specified in order to show that the conditional expectation of the score is zero?\\
c. Use the first-order condition to find $\hat{\sigma}^{2}$ in terms of $\hat{\boldsymbol{\beta}}$.\\
d. Find the Hessian of the $\log$-likelihood function with respect to $\boldsymbol{\theta}$.\\
e. Show directly that $-\mathrm{E}\left[\mathbf{H}_{i}\left(\boldsymbol{\theta}_{\mathrm{o}}\right) \mid \mathbf{x}_{i}\right]=\mathrm{E}\left[\mathbf{s}_{i}\left(\boldsymbol{\theta}_{\mathrm{o}}\right) \mathbf{s}_{i}\left(\boldsymbol{\theta}_{\mathrm{o}}\right)^{\prime} \mid \mathbf{x}_{i}\right]$.\\
f. Propose an estimated asymptotic variance of $\hat{\boldsymbol{\beta}}$, and explain how to obtain the asymptotic standard errors.\\
13.3. Consider a general binary response model $\mathrm{P}\left(y_{i}=1 \mid \mathbf{x}_{i}\right)=G\left(\mathbf{x}_{i}, \boldsymbol{\theta}_{\mathrm{o}}\right)$, where $G(\mathbf{x}, \boldsymbol{\theta})$ is strictly between zero and one for all $\mathbf{x}$ and $\boldsymbol{\theta}$. Here, $\mathbf{x}$ and $\boldsymbol{\theta}$ need not have the same dimension; let $\mathbf{x}$ be a $K$-vector and $\boldsymbol{\theta}$ a $P$-vector.\\
a. Write down the log likelihood for observation $i$.\\
b. Find the score for each $i$. Show directly that $\mathrm{E}\left[\mathbf{s}_{i}\left(\boldsymbol{\theta}_{\mathrm{o}}\right) \mid \mathbf{x}_{i}\right]=\mathbf{0}$.\\
c. When $G(\mathbf{x}, \boldsymbol{\theta})=\Phi\left[\mathbf{x} \boldsymbol{\beta}+\delta_{1}(\mathbf{x} \boldsymbol{\beta})^{2}+\delta_{2}(\mathbf{x} \boldsymbol{\beta})^{3}\right]$, find the LM statistic for testing $H_{0}$ : $\delta_{\mathrm{o} 1}=0, \delta_{\mathrm{o} 2}=0$.\\
d. How would you compute a variable addition version of the test in part c?\\
13.4. In the Newey-Tauchen-White specification-testing context, explain why we can take $\mathbf{g}(\mathbf{w}, \boldsymbol{\theta})=a(\mathbf{x}, \boldsymbol{\theta}) \mathbf{s}(\mathbf{w}, \boldsymbol{\theta})$, where $a(\mathbf{x}, \boldsymbol{\theta})$ is essentially any scalar function of $\mathbf{x}$ and $\boldsymbol{\theta}$.\\
13.5. In the context of CMLE, consider a reparameterization of the kind in Section 12.6.2: $\boldsymbol{\phi}=\mathbf{g}(\boldsymbol{\theta})$, where the Jacobian of $\mathbf{g}, \mathbf{G}(\boldsymbol{\theta})$, is continuous and nonsingular for all $\boldsymbol{\theta} \in \boldsymbol{\Theta}$. Let $\mathbf{s}_{i}^{g}(\boldsymbol{\phi})=\mathbf{s}_{i}^{g}[\mathbf{g}(\boldsymbol{\theta})]$ denote the score of the log likelihood in the reparameterized model; thus, from Section 12.6.2, $\mathbf{s}_{i}^{g}(\boldsymbol{\phi})=\left[\mathbf{G}(\boldsymbol{\theta})^{\prime}\right]^{-1} \mathbf{s}_{i}(\boldsymbol{\theta})$.\\
a. Using the conditional information matrix equality, find $\mathbf{A}_{i}^{g}\left(\boldsymbol{\phi}_{\mathrm{o}}\right) \equiv \mathrm{E}\left[\mathbf{s}_{i}^{g}\left(\boldsymbol{\phi}_{\mathrm{o}}\right) \mathbf{s}_{i}^{g}\left(\boldsymbol{\phi}_{\mathrm{o}}\right)^{\prime} \mid \mathbf{x}_{i}\right]$ in terms of $\mathbf{G}\left(\boldsymbol{\theta}_{\mathrm{o}}\right)$ and $\mathbf{A}_{i}\left(\boldsymbol{\theta}_{\mathrm{o}}\right) \equiv \mathrm{E}\left[\mathbf{s}_{i}\left(\boldsymbol{\theta}_{\mathrm{o}}\right) \mathbf{s}_{i}\left(\boldsymbol{\theta}_{\mathrm{o}}\right)^{\prime} \mid \mathbf{x}_{i}\right]$.\\
b. Show that $\tilde{\mathbf{A}}_{i}^{g}=\tilde{\mathbf{G}}^{\prime-1} \tilde{\mathbf{A}}_{i} \tilde{\mathbf{G}}^{-1}$, where these are all evaluated at the restricted estimate, $\tilde{\boldsymbol{\theta}}$.\\
c. Use part $b$ to show that the expected Hessian form of the LM statistic is invariant to reparameterization.\\
13.6. Suppose that for a panel data set with $T$ time periods, $y_{i t}$ given $\mathbf{x}_{i t}$ has a Poisson distribution with mean $\exp \left(\mathbf{x}_{i t} \boldsymbol{\theta}_{\mathrm{o}}\right), t=1, \ldots, T$.\\
a. Do you have enough information to construct the joint distribution of $\mathbf{y}_{i}$ given $\mathbf{x}_{i}$ ? Explain.\\
b. Write down the partial log likelihood for each $i$ and find the score, $\mathbf{s}_{i}(\boldsymbol{\theta})$.\\
c. Show how to estimate $\operatorname{Avar}(\hat{\boldsymbol{\theta}})$; it should be of the form (13.53).\\
d. How does the estimator of $\operatorname{Avar}(\hat{\boldsymbol{\theta}})$ simplify if the conditional mean is dynamically complete?\\
13.7. Suppose that you have two parametric models for conditional densities: $g\left(y_{1} \mid y_{2}, \mathbf{x} ; \boldsymbol{\theta}\right)$ and $h\left(y_{2} \mid \mathbf{x} ; \boldsymbol{\theta}\right)$; not all elements of $\boldsymbol{\theta}$ need to appear in both densities. Denote the true value of $\boldsymbol{\theta}$ by $\boldsymbol{\theta}_{\mathrm{o}}$.\\
a. What is the joint density of $\left(y_{1}, y_{2}\right)$ given $\mathbf{x}$ ? How would you estimate $\boldsymbol{\theta}_{\mathrm{o}}$ given a random sample on $\left(\mathbf{x}, y_{1}, y_{2}\right)$ ?\\
b. Suppose now that a random sample is not available on all variables. In particular, $y_{1}$ is observed only when $\left(\mathbf{x}, y_{2}\right)$ satisfies a known rule. For example, when $y_{2}$ is binary, $y_{1}$ is observed only when $y_{2}=1$. We assume $\left(\mathbf{x}, y_{2}\right)$ is always observed. Let $r_{2}$ be a binary variable equal to one if $y_{1}$ is observed and zero otherwise. A partial MLE is obtained by defining\\
$\ell_{i}(\boldsymbol{\theta})=r_{i 2} \log g\left(y_{i 1} \mid y_{i 2}, \mathbf{x}_{i} ; \boldsymbol{\theta}\right)+\log h\left(y_{i 2} \mid \mathbf{x}_{i} ; \boldsymbol{\theta}\right) \equiv r_{i 2} \ell_{i 1}(\boldsymbol{\theta})+\ell_{i 2}(\boldsymbol{\theta})$\\
for each $i$. This formulation ensures that first part of $\ell_{i}$ only enters the estimation when $y_{i 1}$ is observed. Verify that $\boldsymbol{\theta}_{\mathrm{o}}$ maximizes $\mathrm{E}\left[\ell_{i}(\boldsymbol{\theta})\right]$ over $\boldsymbol{\Theta}$.\\
c. Show that $-\mathrm{E}\left[\mathbf{H}_{i}\left(\boldsymbol{\theta}_{\mathrm{o}}\right)\right]=\mathrm{E}\left[\mathbf{s}_{i}\left(\boldsymbol{\theta}_{\mathrm{o}}\right) \mathbf{s}_{i}\left(\boldsymbol{\theta}_{\mathrm{o}}\right)^{\prime}\right]$, even though the problem is not a true conditional MLE problem (and therefore a conditional information matrix equality does not hold).\\
d. Argue that a consistent estimator of Avar $\sqrt{N}\left(\hat{\boldsymbol{\theta}}-\boldsymbol{\theta}_{\mathrm{o}}\right)$ is

$$
\left[N^{-1} \sum_{i=1}^{N}\left(r_{i 2} \hat{\mathbf{A}}_{i 1}+\hat{\mathbf{A}}_{i 2}\right)\right]^{-1}
$$

where $\mathbf{A}_{i 1}\left(\boldsymbol{\theta}_{\mathrm{o}}\right)=-\mathrm{E}\left[\nabla_{\theta}^{2} \ell_{i 1}\left(\boldsymbol{\theta}_{\mathrm{o}}\right) \mid y_{i 2}, \mathbf{x}_{i}\right], \mathbf{A}_{i 2}\left(\boldsymbol{\theta}_{\mathrm{o}}\right)=-\mathrm{E}\left[\nabla_{\theta}^{2} \ell_{i 2}\left(\boldsymbol{\theta}_{\mathrm{o}}\right) \mid \mathbf{x}_{i}\right]$, and $\hat{\boldsymbol{\theta}}$ replaces $\boldsymbol{\theta}_{\mathrm{o}}$ in obtaining the estimates.\\
13.8. Suppose that for random vectors $\mathbf{y}_{i}, \mathbf{x}_{i}$, and $\mathbf{w}_{i}$,\\
$\mathrm{D}\left(\mathbf{y}_{i} \mid \mathbf{x}_{i}, \mathbf{w}_{i}\right)=\mathrm{D}\left(\mathbf{y}_{i} \mid \mathbf{x}_{i}, \mathbf{g}\left(\mathbf{w}_{i}, \gamma_{\mathrm{o}}\right)\right) \equiv \mathrm{D}\left(\mathbf{y}_{i} \mid \mathbf{x}_{i}, \mathbf{g}_{i}\right)$,\\
where $\mathbf{g}(\mathbf{w}, \boldsymbol{\gamma})$ is a known function of $\mathbf{w}$ and the $J$ parameters $\boldsymbol{\gamma}$ and $\mathbf{g}_{i} \equiv \mathbf{g}\left(\mathbf{w}_{i}, \gamma_{\mathrm{o}}\right)$ (which is unobserved because we generally do not know $\gamma_{\mathrm{o}}$ ). Assume that we have a correctly specified density for $\mathrm{D}\left(\mathbf{y}_{i} \mid \mathbf{x}_{i}, \mathbf{g}_{i}\right), f(\mathbf{y} \mid \mathbf{x}, \mathbf{g} ; \boldsymbol{\theta})$, and let the $P \times 1$ vector $\boldsymbol{\theta}_{\mathrm{o}}$ denote the true value. Let $\hat{\gamma}$ be an $\sqrt{N}$-asymptotically normal estimator of $\gamma_{\mathrm{o}}$ that satisfies\\
$\sqrt{N}\left(\hat{\gamma}-\gamma_{\mathrm{o}}\right)=N^{-1 / 2} \sum_{i=1}^{N} \mathbf{r}\left(\mathbf{w}_{i}, \gamma_{\mathrm{o}}\right)+o_{p}(1)$,\\
and let $\hat{\boldsymbol{\theta}}$ be the two-step nonlinear least squares estimator that solves\\
$\max _{\theta \in \Theta} \sum_{i=1}^{N} \log f\left(\mathbf{y}_{i} \mid \mathbf{x}_{i}, \mathbf{g}\left(\mathbf{w}_{i}, \hat{\gamma}\right) ; \boldsymbol{\theta}\right)$.\\
a. Show that, under standard regularity conditions, the asymptotic variance of $\sqrt{N}\left(\hat{\boldsymbol{\theta}}-\boldsymbol{\theta}_{\mathrm{o}}\right)$ is at least as large as the asymptotic variance if we knew $\boldsymbol{\gamma}_{\mathrm{o}}$.\\
b. Propose a consistent estimator of $\mathrm{Avar} \sqrt{N}\left(\hat{\boldsymbol{\theta}}-\boldsymbol{\theta}_{\mathrm{o}}\right)$.\\
c. Suppose that $y_{i}$ is a scalar and $\mathrm{P}\left(y_{i}=1 \mid \mathbf{x}_{i}, \mathbf{w}_{i}\right)=\Phi\left(\mathbf{x}_{i} \boldsymbol{\beta}_{\mathrm{o}}+\rho_{\mathrm{o}} g_{i}\right)$, where $\mathbf{w}_{i}= \left(\mathbf{z}_{i}, h_{i}\right), g_{i}=h_{i}-\mathbf{z}_{i} \gamma_{\mathrm{o}}$, and $\mathrm{E}\left(\mathbf{z}_{i}^{\prime} g_{i}\right)=\mathbf{0}$. Let $\hat{\boldsymbol{\theta}}=\left(\hat{\boldsymbol{\beta}}^{\prime}, \hat{\rho}\right)^{\prime}$ be the two-step probit estimators from probit of $y_{i}$ on $\mathbf{x}_{i}$, and $\hat{g}_{i}$, where $\hat{g}_{i}=h_{i}-\mathbf{z}_{i} \hat{\gamma}$ are OLS residuals from $h_{i}$ on $\mathbf{z}_{i}, i=1, \ldots, N$. How does part a apply in this case?\\
d. Show that, when $\rho_{\mathrm{o}}=0$, the usual probit variance matrix estimator of $\hat{\boldsymbol{\theta}}$ is asymptotically valid. That is, valid inference is obtained by ignoring the first-stage estimation of $\gamma_{\mathrm{o}}$.\\
e. How would you test $\mathrm{H}_{\mathrm{o}}: \rho_{\mathrm{o}}=0$ ?\\
13.9. Let $\left\{y_{t}: t=0,1, \ldots, T\right\}$ be an observable time series representing a population, where we use the convention that $t=0$ is the first time period for which $y$ is observed. Assume that the sequence follows a Markov process: $\mathrm{D}\left(y_{t} \mid y_{t-1}, y_{t-2}, \ldots y_{0}\right) =\mathrm{D}\left(y_{t} \mid y_{t-1}\right)$ for all $t \geq 1$. Let $f_{t}\left(y_{t} \mid y_{t-1} ; \boldsymbol{\theta}\right)$ denote a correctly specified model for the density of $y_{t}$ given $y_{t-1}, t \geq 1$, where $\boldsymbol{\theta}_{\mathrm{o}}$ is the true value of $\boldsymbol{\theta}$.\\
a. Show that, to obtain the joint distribution of $\left(y_{0}, y_{2}, \ldots, y_{T}\right)$, you need to correctly model the density of $y_{0}$.\\
b. Given a random sample of size $N$ from the population, that is, $\left(y_{i 0}, y_{i 1}, \ldots, y_{i T}\right)$ for each $i$, explain how to consistently etimate $\boldsymbol{\theta}_{\mathrm{o}}$ without modeling $\mathrm{D}\left(y_{0}\right)$.\\
c. How would you estimate the asymptotic variance of the estimator from part $b$ ? Be specific.\\
13.10. Let $\mathbf{y}$ be a $G \times 1$ random vector with elements $y_{g}, g=1,2, \ldots, G$. These could be different response variables for the same cross section unit or responses at different points in time. Let $\mathbf{x}$ be a $K$-vector of observed conditioning variables, and let $c$ be an unobserved conditioning variable. Let $f_{g}(\cdot \mid \mathbf{x}, c)$ denote the density of $y_{g}$ given $(\mathbf{x}, c)$. Further, assume that the $y_{1}, y_{2}, \ldots, y_{G}$ are independent conditional on $(\mathbf{x}, c)$.\\
a. Write down the joint density of $\mathbf{y}$ given $(\mathbf{x}, c)$.\\
b. Let $h(\cdot \mid \mathbf{x})$ be the density of $c$ given $\mathbf{x}$. Find the joint density of $\mathbf{y}$ given $\mathbf{x}$.\\
c. If each $f_{g}(\cdot \mid \mathbf{x}, c)$ is known up to a $P_{g}$-vector of parameters $\gamma_{\mathrm{o}}^{g}$ and $h(\cdot \mid \mathbf{x})$ is known up to an $M$-vector $\boldsymbol{\delta}_{\mathrm{o}}$, find the log likelihood for any random draw $\left(\mathbf{x}_{i}, y_{i}\right)$ from the population.\\
d. Is there a relationship between this setup and a linear SUR model?\\
13.11. Consider the dynamic, linear unobserved effects model\\
$y_{i t}=\rho y_{i, t-1}+c_{i}+e_{i t}, \quad t=1,2, \ldots, T$\\
$\mathrm{E}\left(e_{i t} \mid y_{i, t-1}, y_{i, t-2}, \ldots, y_{i 0}, c_{i}\right)=0$\\
In Section 11.6.2 we discussed estimation of $\rho$ by instrumental variables methods after differencing. The deficiencies of the IV approach for large $\rho$ may be overcome by applying the conditional MLE methods in Section 13.9.2.\\
a. Make the stronger assumption that $y_{i t} \mid\left(y_{i, t-1}, y_{i, t-2}, \ldots, y_{i 0}, c_{i}\right)$ is normally distributed with mean $\rho y_{i, t-1}+c_{i}$ and variance $\sigma_{e}^{2}$. Find the density of $\left(y_{i 1}, \ldots, y_{i T}\right)$ given $\left(y_{i 0}, c_{i}\right)$. Is it a good idea to use the log of this density, summed across $i$, to estimate $\rho$ and $\sigma_{e}^{2}$ along with the "fixed effects" $c_{i}$ ?\\
b. If $c_{i} \mid y_{i 0} \sim \operatorname{Normal}\left(\alpha_{0}+\alpha_{1} y_{i 0}, \sigma_{a}^{2}\right)$, where $\sigma_{a}^{2} \equiv \operatorname{Var}\left(a_{i}\right)$ and $a_{i} \equiv c_{i}-\alpha_{0}-\alpha_{1} y_{i 0}$, write down the density of $\left(y_{i 1}, \ldots, y_{i T}\right)$ given $y_{i 0}$. How would you estimate $\rho, \alpha_{0}, \alpha_{1}$, $\sigma_{e}^{2}$, and $\sigma_{a}^{2}$ ?\\
c. Under the same assumptions in parts a and b , extend the model to $y_{i t}=\rho y_{i, t-1}+ c_{i}+\delta c_{i} y_{i, t-1}+e_{i t}$. Explain how to estimate the parameters of this model, and propose a consistent estimator of the average partial effect of the lag, $\rho+\delta \mathrm{E}\left(c_{i}\right)$.\\
d. Now extend part b to the case where $\mathbf{z}_{i t} \boldsymbol{\beta}$ is added to the conditional mean function, where the $\mathbf{z}_{i t}$ are strictly exogenous conditional on $c_{i}$. Assume that $c_{i} \mid y_{i 0}, \mathbf{z}_{i} \sim \operatorname{Normal}\left(\alpha_{0}+\alpha_{1} y_{i 0}+\overline{\mathbf{z}}_{i} \boldsymbol{\delta}, \sigma_{a}^{2}\right)$, where $\overline{\mathbf{z}}_{i}$ is the vector of time averages.\\
13.12. In the context of GLM, for a given LEF density, there are many ways to characterize what is known as the canonical link function. One is that $g(\cdot)$ is the canonical link if the first-order conditions for the QMLE can be expressed as\\
$\sum_{i=1}^{n} \mathbf{x}_{i}^{\prime}\left[y_{i}-g^{-1}\left(\mathbf{x}_{i} \hat{\boldsymbol{\theta}}\right)\right] \equiv \sum_{i=1}^{n} \mathbf{x}_{i}^{\prime}\left[y_{i}-m\left(\mathbf{x}_{i}, \hat{\boldsymbol{\theta}}\right)\right]=\mathbf{0}$,\\
where $m(\mathbf{x}, \boldsymbol{\theta})$ is the corresponding mean function and $\hat{\boldsymbol{\theta}}$ is the QMLE. Alternatively, using the canonical link implies that the Hessian of the quasi-log likelihood (QLL) does not depend on $y_{i}$; consequently, the Hessian and the expected value of the Hessian given $\mathbf{x}_{i}$ are identical.\\
a. Define residuals as usual, $\hat{u}_{i}=y_{i}-m\left(\mathbf{x}_{i}, \hat{\boldsymbol{\theta}}\right)$. Argue that if $\mathbf{x}_{i}$ contains a constant-as it almost always does in GLM analysis-then the residuals obtained by using the canonical link always average to zero. (This extends the well-known result for linear regression.)\\
b. Show that for the Bernoulli QLL, the canonical link function is $g(\mu)= \log [\mu /(1-\mu)]$, so that the mean function is the logistic function, $\exp (\mathbf{x} \boldsymbol{\theta}) / [1+\exp (\mathbf{x} \boldsymbol{\theta})]$.\\
c. Show that for the Poisson QLL, the canonical link is $g(\mu)=\log (\mu)$. What is the mean function?\\
d. Evaluate the following statement: "In a GLM analysis using the canonical link, our estimate of the asymptotic variance of $\hat{\boldsymbol{\theta}}$ is the same whether or not we assume the mean function is correctly specified."\\
13.13. Let $\boldsymbol{\theta}^{*}$ be the solution to the maximization problem

$$
\max _{\boldsymbol{\theta} \in \boldsymbol{\Theta}} \mathrm{E}\left[\ell\left(\mathbf{w}_{i}, \boldsymbol{\theta}\right)\right],
$$

where $\ell(\mathbf{w}, \boldsymbol{\theta})$ is a quasi-log likelihood satisfying the conditions of Theorem 12.3 (but with maximization rather than minimization). Assume that $\boldsymbol{\theta}^{*}$ is in the interior of $\boldsymbol{\Theta}$ and that the gradients and expectation can be interchanged. Let $\hat{\boldsymbol{\theta}}$ be the quasi-MLE. Show that

$$
N^{-1 / 2} \sum_{i=1}^{N} \ell\left(\mathbf{w}_{i}, \hat{\boldsymbol{\theta}}\right)=N^{-1 / 2} \sum_{i=1}^{N} \ell\left(\mathbf{w}_{i}, \boldsymbol{\theta}^{*}\right)+o_{p}(1) .
$$

\section*{Appendix 13A}
In this appendix we cover some important properties of conditional distributions and conditional densities. Billingsley (1979) is a good reference for this material. For random vectors $\mathbf{y} \in \mathscr{Y} \subset \mathbb{R}^{G}$ and $\mathbf{x} \in \mathscr{X} \subset \mathbb{R}^{K}$, the conditional distribution of $\mathbf{y}$ given $\mathbf{x}$ always exists and is denoted $\mathrm{D}(\mathbf{y} \mid \mathbf{x})$. For each $\mathbf{x}$ this distribution is a probability measure and completely describes the behavior of the random vector $\mathbf{y}$ once $\mathbf{x}$ takes on a particular value. In econometrics, we almost always assume that this distribution is described by a conditional density, which we denote by $p(\cdot \mid \mathbf{x})$. The density is with respect to a measure defined on the support $\mathscr{Y}$ of $\mathbf{y}$. A conditional density makes sense only when this measure does not change with the value of $\mathbf{x}$. In practice, this assumption is not very restrictive, as it means that the nature of $\mathbf{y}$ is not dramatically different for different values of $\mathbf{x}$. Let $v$ be this measure on $\mathbb{R}^{J}$. If $\mathrm{D}(\mathbf{y} \mid \mathbf{x})$ is discrete, $v$ can be the counting measure and all integrals are sums. If $\mathrm{D}(\mathbf{y} \mid \mathbf{x})$ is absolutely continuous, then $v$ is the familiar Lebesgue measure appearing in elementary integration theory. In some cases, $\mathrm{D}(\mathbf{y} \mid \mathbf{x})$ has both discrete and continuous characteristics.

The important point is that all conditional probabilities can be obtained by integration:\\
$\mathrm{P}(\mathbf{y} \in \mathscr{A} \mid \mathbf{x}=x)=\int_{A} p(y \mid x) v(\mathrm{~d} y)$,\\
where $y$ is the dummy argument of integration. When $\mathbf{y}$ is discrete, taking on the values $y_{1}, y_{2}, \ldots$, then $p(\cdot \mid x)$ is a probability mass function and $\mathbf{P}\left(\mathbf{y}=y_{j} \mid \mathbf{x}=x\right)= p\left(y_{j} \mid x\right), j=1,2, \ldots$.

Suppose that $f$ and $g$ are nonnegative functions on $\mathbb{R}^{M}$, and define $\mathscr{S}_{f} \equiv \left\{z \in \mathbb{R}^{M}: f(z)>0\right\}$. Assume that\\
$1=\int_{\mathscr{S}_{f}} f(z) v(d z) \geq \int_{\mathscr{S}_{f}} g(z) v(\mathrm{~d} z)$,\\
where $v$ is a measure on $\mathbb{R}^{M}$. The equality in expression (13.99) implies that $f$ is a density on $\mathbb{R}^{M}$, while the inequality holds if $g$ is also a density on $\mathbb{R}^{M}$. An important result is that\\
$\mathscr{I}(f ; g) \equiv \int_{\mathscr{S}_{f}} \log [f(z) / g(z)] f(z) v(\mathrm{~d} z) \geq 0$.\\
(Note that $\mathscr{I}(f ; g)=\infty$ is allowed; one case where this result can occur is $f(z)>0$ but $g(z)=0$ for some $z$. Also, the integrand is not defined when $f(z)=g(z)=0$, but such values of $z$ have no effect because the integrand receives zero weight in the integration.) The quantity $\mathscr{I}(f ; g)$ is called the Kullback-Leibler information criterion (KLIC). Another way to state expression (13.100) is\\
$\mathrm{E}\{\log [f(\mathbf{z})]\} \geq \mathrm{E}\{\log [g(\mathbf{z})]\}$,\\
where $\mathbf{z} \in \mathscr{Z} \subset \mathbb{R}^{M}$ is a random vector with density $f$.\\
Conditional MLE relies on a conditional version of inequality (13.99):\\
PROPERTY CD.1: Let $\mathbf{y} \in \mathscr{Y} \subset \mathbb{R}^{G}$ and $\mathbf{x} \in \mathscr{X} \subset \mathbb{R}^{K}$ be random vectors. Let $p(\cdot \mid \cdot)$ denote the conditional density of $\mathbf{y}$ given $\mathbf{x}$. For each $\mathbf{x}$, let $\mathscr{Y}(\mathbf{x}) \equiv\{y: p(y \mid \mathbf{x})>0\}$ be the conditional support of $\mathbf{y}$, and let $v$ be a measure that does not depend on $\mathbf{x}$. Then, for any other function $g(\cdot \mid \mathbf{x}) \geq 0$ such that

$$
1=\int_{\mathscr{Y}(\mathbf{x})} p(y \mid \mathbf{x}) v\left(\mathrm{~d}_{y}\right) \geq \int_{\mathscr{Y}(\mathbf{x})} g(y \mid \mathbf{x}) v\left(\mathrm{~d}_{y}\right),
$$

the conditional KLIC is nonnegative.\\
$\mathscr{I}_{\mathbf{x}}(p ; g) \equiv \int_{\mathscr{Y}(\mathbf{x})} \log [p(y \mid \mathbf{x}) / g(y \mid \mathbf{x})] p(y \mid \mathbf{x}) v\left(\mathrm{~d}_{y}\right) \geq 0$.\\
That is,\\
$\mathrm{E}\{\log [p(\mathbf{y} \mid \mathbf{x})] \mid \mathbf{x}\} \geq \mathrm{E}\{\log [g(\mathbf{y} \mid \mathbf{x})] \mid \mathbf{x}\}$\\
for any $\mathbf{x} \in \mathscr{X}$. The proof uses the conditional Jensen's inequality (Property CE. 7 in Chapter 2). See Manski (1988, Sect. 5.1).

PROPERTY CD.2: For random vectors $\mathbf{y}, \mathbf{x}$, and $\mathbf{z}$, let $p(y \mid \mathbf{x}, \mathbf{z})$ be the conditional density of $\mathbf{y}$ given ( $\mathbf{x}, \mathbf{z}$ ) and let $p(x \mid \mathbf{z})$ denote the conditional density of $\mathbf{x}$ given $\mathbf{z}$. Then the density of $(\mathbf{y}, \mathbf{x})$ given $\mathbf{z}$ is\\
$p(y, x \mid \mathbf{z})=p(y \mid x, \mathbf{z}) p(x \mid \mathbf{z})$,\\
where the script variables are placeholders.\\
property CD.3: For random vectors $\mathbf{y}, \mathbf{x}$, and $\mathbf{z}$, let $p(y \mid \mathbf{x}, \mathbf{z})$ be the conditional density of $\mathbf{y}$ given $(\mathbf{x}, \mathbf{z})$, let $p(y \mid \mathbf{x})$ be the conditional density of $\mathbf{y}$ given $\mathbf{x}$, and let $p(z \mid \mathbf{x})$ denote the conditional density of $\mathbf{z}$ given $\mathbf{x}$ with respect to the measure $v(\mathrm{~d} z)$.\\
Then\\
$p(y \mid \mathbf{x})=\int_{\mathscr{Z}} p(y \mid \mathbf{x}, z) p(z \mid \mathbf{x}) v(\mathrm{~d} z)$.\\
In other words, we can obtain the density of $\mathbf{y}$ given $\mathbf{x}$ by integrating the density of $\mathbf{y}$ given the larger conditioning set, $(\mathbf{x}, \mathbf{z})$, against the density of $\mathbf{z}$ given $\mathbf{x}$.\\
property CD.4: Suppose that the random variable, $u$, with $\mathrm{cdf}, F$, is independent of the random vector $\mathbf{x}$. Then, for any function $a(\mathbf{x})$ of $\mathbf{x}, \mathrm{P}[u \leq a(\mathbf{x}) \mid \mathbf{x}]=F[a(\mathbf{x})]$.

\section*{14 Generalized Method of Moments and Minimum Distance Estimation}

\end{document}
